% Options for packages loaded elsewhere
\PassOptionsToPackage{unicode}{hyperref}
\PassOptionsToPackage{hyphens}{url}
\documentclass[
  11pt,
]{article}
\usepackage{xcolor}
\usepackage[margin=1in]{geometry}
\usepackage{amsmath,amssymb}
\setcounter{secnumdepth}{-\maxdimen} % remove section numbering
\usepackage{iftex}
\ifPDFTeX
  \usepackage[T1]{fontenc}
  \usepackage[utf8]{inputenc}
  \usepackage{textcomp} % provide euro and other symbols
\else % if luatex or xetex
  \usepackage{unicode-math} % this also loads fontspec
  \defaultfontfeatures{Scale=MatchLowercase}
  \defaultfontfeatures[\rmfamily]{Ligatures=TeX,Scale=1}
\fi
\usepackage{lmodern}
\ifPDFTeX\else
  % xetex/luatex font selection
\fi
% Use upquote if available, for straight quotes in verbatim environments
\IfFileExists{upquote.sty}{\usepackage{upquote}}{}
\IfFileExists{microtype.sty}{% use microtype if available
  \usepackage[]{microtype}
  \UseMicrotypeSet[protrusion]{basicmath} % disable protrusion for tt fonts
}{}
\makeatletter
\@ifundefined{KOMAClassName}{% if non-KOMA class
  \IfFileExists{parskip.sty}{%
    \usepackage{parskip}
  }{% else
    \setlength{\parindent}{0pt}
    \setlength{\parskip}{6pt plus 2pt minus 1pt}}
}{% if KOMA class
  \KOMAoptions{parskip=half}}
\makeatother
\setlength{\emergencystretch}{3em} % prevent overfull lines
\providecommand{\tightlist}{%
  \setlength{\itemsep}{0pt}\setlength{\parskip}{0pt}}
\usepackage{mathtools,mathrsfs}
\emergencystretch=2em
\AtBeginDocument{\hypersetup{pdftitle={Magnetic continuation and the interacting Wilson vacuum},pdfauthor={}}}
\usepackage{bookmark}
\IfFileExists{xurl.sty}{\usepackage{xurl}}{} % add URL line breaks if available
\urlstyle{same}
\hypersetup{
  hidelinks,
  pdfcreator={LaTeX via pandoc}}

\author{}
\date{}

\begin{document}

\section{Magnetic continuation, Gauss projection, and the interacting
vacuum}\label{magnetic-part-1-magnetic-continuation-gauss-projection-and-the-interacting-vacuum}

8 September 2026. Complete finite-regulator calculation. The connection
to the original cusp is a map of its magnetic background into the full
Wilson Hilbert space. Every spatial plaquette and the original electric
coefficient are retained. No continuum vacuum or mass-gap conclusion is
assumed. Only the explicit regular-fibre period and magnetic-bundle data
are used; the claimed identification of a completed complex threefold
with S6 is not an input to this calculation.

\subsection{1. Original objects and the reference-frame
map}\label{magnetic-part-1-original-objects-and-the-reference-frame-map}

Let \(L\ge2\) be an integer. Vertices are \(n\in\{-L,\ldots,L\}^3\); a
physical edge \(e=(n,i)\) has positive direction i and exists whenever n
and \(n+e_i\) lie in the box. Its independent variable is
\(U_i(n)\in SU(2)\). Negative traversal uses its inverse. Every
elementary face \(p=(n;i,j)\), \(i<j\), is present, with \[
 W_p(U)=\operatorname{tr}\bigl(U_i(n)U_j(n+e_i)
                 U_i(n+e_j)^{-1}U_j(n)^{-1}\bigr).
\] There are \(N=3(2L)(2L+1)^2\) links and \(M=3(2L)^2(2L+1)\) faces.
The configuration manifold is \(Q=SU(2)^N\) and dU is product Haar
probability measure. The Hilbert inner product is conjugate-linear in
its first argument. Put \(T_a=-i\sigma_a/2\) and define \[
 X_{e,a}f(U)=\left.\frac{d}{dt}\right|_{t=0}
 f(\ldots,e^{tT_a}U_e,\ldots),\qquad
 \Delta_e^T=\sum_{a=1}^3X_{e,a}^2.
\] Haar invariance proves X is skew-adjoint. The original Lie metric is
\(c(A,B)=-\operatorname{tr}(AB)/2\): its orthonormal frame is 2T\_a, so
\(\Delta_e^c=4\Delta_e^T\). The Hamiltonian is \[
 H=-\kappa\sum_e\Delta_e^T+V,\qquad
 V=b\sum_p(2-W_p),\qquad
 \kappa=\frac{2g_{\rm YM}^2}{a},\quad
 b=\frac{1}{2g_{\rm YM}^2a},\quad \xi=\frac b\kappa.
\tag{1.1}
\] Equivalently its kinetic coefficient in the original metric is
\(-g_{\rm YM}^2/(2a)\) times \(\sum_e\Delta_e^c\). The full constant 2bM
is part of H.

Gauge transformations act as \((g\cdot U)_e=g_{s(e)}U_eg_{t(e)}^{-1}\).
Let \(R_gf(U)=f(g^{-1}\cdot U)\), and
\(\Pi=\int_{SU(2)^{\mathsf V}}R_g\,dg\). Haar invariance gives
\(R_g^*=R_{g^{-1}}\); averaging proves \(\Pi^*=\Pi^2=\Pi\). Its range is
the gauge-invariant Hilbert space. The kinetic energy is bi-invariant on
each link and the face traces are invariant, so H and its closed
quadratic form commute with Pi.

It is important to specify what is translated. Let M\_e and R\_e be,
respectively, the original magnetic and a fixed reference parallel
transport, in the same endpoint frames. Define \[
 h_e=M_eR_e^{-1},\qquad (T_hf)(U)=f((h_e^{-1}U_e)_e).
\tag{1.2}
\] Here the reference symbol R\_e is not the gauge-representation
operator R\_g. Under a frame change both transports change by endpoint
conjugation; therefore \[
 M'_e=g_sM_eg_t^{-1},\quad R'_e=g_sR_eg_t^{-1},\quad
 h'_e=g_sh_eg_s^{-1},\quad R_gT_hR_g^{-1}=T_{h'}.
\tag{1.3}
\] This proves the exact background-to-state covariance. Translating by
the magnetic transport while changing its frames but not those of the
reference would not be the same construction. In the original fixed
frames the flat reference is R\_e=I. All statements below retain that
reference; no assertion that left multiplication is addition of
arbitrary non-Abelian connections is needed.

The original nonwrapping magnetic links, in the physical plane
identified explicitly below, are \[
 h_1(n)=\exp(-\theta n_2 H_c),\quad h_2(n)=h_3(n)=I,
 \qquad H_c=i\sigma_3=-2T_3,\qquad
 \theta=\frac{2\pi a^2}{D},\quad D=L_Tq_T+6m_T^2.
\tag{1.4}
\] The period parameters are \(L_T=\operatorname{Im}\tau_T\),
\(q_T=-\operatorname{Im}\beta_T\), \(m_T=\operatorname{Im}\mu_T\), and
D\textgreater0. The box is a nonwrapping physical patch; it is not the
entire periodic fibre with transition edges removed. In the original
periodic coordinates u, the underlying connection and transition are \[
 d+2\pi u_1H_c\,du_2,\qquad
 G_k(u)=\exp(-2\pi k_1u_2H_c).
\] Its exact finite-vertex frame change is
\(h(n)=\exp(-2\pi H_cn_1n_2/N_0^2)\), where N\_0 is the number of
subdivisions in each original compact direction, not the number N of
links. It gives U\_1=exp(-2pi H\_c n\_2/N\_0\^{}2) and U\_2=I on
nonwrapping edges; on n\_2=N\_0-1 the second-direction link is exp(2pi
H\_c n\_1/N\_0). These wrap factors and the full period-matrix map are
retained in the companion cusp calculation. The physical-coordinate
dictionary for the patch used here can also be written exactly. In
original real coordinates
\((x_1,x_2,x_3,x_4)=(\operatorname{Re}z_1,\operatorname{Im}z_1,
\operatorname{Re}z_2,\operatorname{Im}z_2)\), the imaginary period block
is \[
 B=\begin{pmatrix}6m_T&L_T\\-q_T&m_T\end{pmatrix},\qquad
 \begin{pmatrix}u_1\\u_2\end{pmatrix}
 =\frac1D\begin{pmatrix}m_T&-L_T\\q_T&6m_T\end{pmatrix}
       \begin{pmatrix}x_2\\x_4\end{pmatrix}.
\] The real period block and the covered real periods are not set to
zero: they determine the other coordinates in the full period map; they
do not enter this displayed exact inverse for u\_1,u\_2. Choose
Euclidean temporal coordinate y\_0=x\_3 and spatial coordinates
\((y_1,y_2,y_3)=(x_2,x_4,x_1)\). The full permutation
\((y_0,y_1,y_2,y_3)=(x_3,x_2,x_4,x_1)\) has determinant +1 and preserves
the Euclidean metric; thus no scale or orientation is suppressed. Direct
exterior multiplication gives \(du_1\wedge du_2=D^{-1}dy_1\wedge dy_2\).

The original connection on this patch is exactly \[
 A_{\rm orig}=\frac{2\pi H_c}{D^2}
 (m_Ty_1-L_Ty_2)(q_Tdy_1+6m_Tdy_2).
\] Set \(\mathfrak b=2\pi/D\) and \[
 f(y)=\frac{2\pi}{D^2}
 \left(\frac{m_Tq_T}{2}y_1^2-L_Tq_Ty_1y_2-3m_TL_Ty_2^2\right).
\] Expanding both components proves
\(A_{\rm orig}=\mathfrak bH_c y_1dy_2+H_cdf\). The frame change
exp(-H\_c f), under \(A\mapsto g^{-1}Ag+g^{-1}dg\), therefore gives
\(A_0=\mathfrak bH_c y_1dy_2\). The further frame change
\(\exp(-\mathfrak bH_c y_1y_2)\) gives \(A_1=-\mathfrak bH_c y_2dy_1\).
These are exact frame maps of the same magnetic connection, with all
m\_T,L\_T,q\_T terms shown. Using the stated inverse-parallel-transport
convention, integrating A\_1 over the edge from y=a n to y+a e\_i gives
precisely (1.4). The reference R\_e=I in (1.2) is chosen in this final
displayed frame; subsequent frame changes act on it as well as on M\_e
by (1.3). The 12-face holonomy is exp(theta H\_c); the 13 and 23
holonomies are I.

\subsection{2. Actual vacuum and exact finite-translation
identities}\label{magnetic-part-1-actual-vacuum-and-exact-finite-translation-identities}

The product Laplacian on this compact connected manifold is elliptic,
with domain \(H^2\) and closed form domain \(H^1\). V is smooth and
bounded, so H has the same domains and compact resolvent. Its bottom
eigenvector can be chosen nonnegative: taking the absolute value does
not increase the gradient form, as follows by differentiating
\(\sqrt{|u|^2+\varepsilon^2}\) and passing to the limit. Elliptic
regularity and the strong maximum principle give a smooth strictly
positive eigenvector psi. Choose \(\int\psi^2dU=1\), let E be its
eigenvalue, and put \(\mathcal A=H-E\).

For a smooth function F, integration by parts and H psi=E psi give \[
 \mathfrak q_{\mathcal A}[\psi F]
 =\kappa\sum_{e,a}\int |X_{e,a}F|^2\psi^2\,dU.
\tag{2.1}
\] The same formula after polarization holds on H1 by density. If
another ground eigenfunction existed, its quotient by the positive psi
would make the right hand side zero, hence be constant on Q. This proves
simplicity. Gauge transformations preserve the positive normalized
ground state, so \(\Pi\psi=\psi\). Compactness and simplicity give a
positive finite-regulator first excited energy, but no uniform bound as
L or a varies is asserted here.

Translation T\_h is unitary by Haar invariance, preserves smooth
functions and all Sobolev spaces, and commutes with the full electric
Laplacian. Define \[
 c_h=\langle\psi,T_h\psi\rangle,
 \qquad \chi_h=\Pi T_h\psi-c_h\psi.
\tag{2.2}
\] The scalar c\_h is real and positive, since both pointwise functions
in the integral are positive. The state is gauge invariant, smooth, and
vacuum orthogonal. Its exact norm and excitation form energy are \[
 d_h=\|\chi_h\|^2
 =\langle T_h\psi,\Pi T_h\psi\rangle-c_h^2,
\qquad
 n_h=\mathfrak q_{\mathcal A}[\Pi T_h\psi].
\tag{2.3}
\] When d\_h\textgreater0 its Rayleigh quotient is n\_h/d\_h. No
division by a presumed nonzero norm is made.

There is also an exact comparison with the unprojected state: \[
 0\le n_h\le\mathfrak q_{\mathcal A}[T_h\psi]
 =\int\psi(U)^2\,[V((h_eU_e)_e)-V(U)]\,dU.
\tag{2.4}
\] Indeed Pi commutes with the nonnegative operator A, so its retained
and orthogonal components have additive nonnegative energies. The last
identity follows from \(T_h^*HT_h=H+V(hU)-V(U)\) and H psi=E psi. Thus
even the unprojected magnetic energy is an average in the actual vacuum,
not the action of the magnetic configuration alone. Formula (2.4)
supplies an energy numerator, not a lower bound on d\_h.

\subsection{3. General infinitesimal projection, with the zero branch
resolved}\label{magnetic-part-1-general-infinitesimal-projection-with-the-zero-branch-resolved}

For real coefficients alpha\_\{e,a\}, set \(A_e=\sum_a\alpha_{e,a}T_a\),
\(Y=\sum_{e,a}\alpha_{e,a}X_{e,a}\), and \(h_e(t)=\exp(tA_e)\). Then
\(T_t=\exp(-tY)\) exactly. All expansions in this section are applied to
smooth psi in any fixed Sobolev norm. Taylor's formula for the smooth
translation flow, with its bounded integral remainder on the compact
group, justifies every term and its form-domain use.

The conjugation rule obtained by differentiating (1.3) is \[
 R_gX_e(A)R_g^{-1}=X_e(\operatorname{Ad}_{g_{s(e)}}A).
\tag{3.1}
\] The adjoint SU(2) representation acts by all rotations on the three
real coefficients in the chosen orthonormal Lie basis. Haar averaging of
a vector is zero. Averaging a pair gives
\(\int(\operatorname{Ad}_g u)_a(\operatorname{Ad}_g v)_b,dg
=(u\cdot v)\delta_{ab}/3\): invariance makes the matrix scalar, and its
trace is u dot v. Different vertices have independent Haar variables.
Applying these identities to invariant psi proves \[
 \Pi Y\psi=0,\qquad
 K_2\psi:=\Pi Y^2\psi
 =\frac13\sum_v\sum_{s(e)=s(f)=v}
 (\alpha_e\cdot\alpha_f)\sum_a X_{e,a}X_{f,a}\psi.
\tag{3.2}
\] The second sum is ordered: e=f terms and both orders of unequal edges
are retained. The grouping is by the sources of the chosen positive
edges, not by an unrecorded change to incoming-edge derivatives.

Since Y is skew-adjoint and psi is real, its overlap has
\(c'(0)=\langle\psi,-Y\psi\rangle=0\) and
\(c''(0)=\langle\psi,Y^2\psi\rangle=-\|Y\psi\|^2\). Define \[
 S=\|Y\psi\|^2,\qquad \eta=K_2\psi+S\psi.
\] Then the exact centered physical state obeys \[
 \chi_t=\frac{t^2}{2}\eta+O_{H^1}(t^3),\qquad
 \|\chi_t\|^2=\frac{t^4}{4}\|\eta\|^2+O(t^5),
\quad
 \mathfrak q_{\mathcal A}[\chi_t]
 =\frac{t^4}{4}\mathfrak q_{\mathcal A}[\eta]+O(t^5).
\tag{3.3}
\] The energy remainder follows from continuity of the form on H1; it is
not obtained by treating an unbounded Hamiltonian as bounded.

There is no unexamined branch where eta vanishes but a later derivative
is the first excitation. The Haar integral of every X derivative is
zero, so \(\int K_2\psi\,dU=0\). If eta=0, integrating its defining
identity gives \(S\int\psi\,dU=0\); positivity of psi gives S=0. Hence Y
psi=0, the translation flow fixes psi, and chi\_t=0 for every t.
Conversely that identity implies eta=0. Therefore a translation either
fixes the vacuum exactly or has strictly positive quartic leading norm.
In the nontrivial case the finite-regulator limiting excitation quotient
is \[
 \lim_{t\to0}\frac{\mathfrak q_{\mathcal A}[\chi_t]}{\|\chi_t\|^2}
 =\frac{\mathfrak q_{\mathcal A}[\eta]}{\|\eta\|^2}>0.
\tag{3.4}
\] Strict positivity follows because eta is nonzero, gauge invariant and
orthogonal to psi, whose eigenvalue is simple. This statement holds at
fixed regulator and arbitrary positive original coupling. It does not
take L to infinity or a to zero.

\subsection{4. Exact class-convolution map for the native magnetic
translation}\label{magnetic-part-1-exact-class-convolution-map-for-the-native-magnetic-translation}

For (1.4), there is at most one nonidentity translated edge with any
given source vertex. Conjugating T\_h by R\_g and averaging therefore
factorizes into independent averages for these edges. On invariant
inputs, \[
 \Pi T_h\Pi=K_\theta\Pi,\qquad
 K_\theta=\prod_{e=(n,1)}C_{e,\theta n_2},
\tag{4.1}
\] where the full single-edge operator is \[
 C_{e,s}f(U)=\int_{SU(2)}
 f(\ldots,(g e^{-sH_c}g^{-1})^{-1}U_e,\ldots)\,dg.
\tag{4.2}
\] Different link factors commute. Every C is a contraction and
preserves positive functions, since it averages unitary translations
with a probability measure. Its conjugacy-class measure is central and
invariant under inversion: in SU(2), e\^{}\{sH\_c\} and e\^{}\{-sH\_c\}
are conjugate, for example by i sigma\_1. Consequently C is self-adjoint
and even in s. These facts do not claim it has nonnegative spectral
eigenvalues.

Here is its exact Peter--Weyl action, including all representations. On
a spin-j matrix coefficient on edge e, j=0,1/2,1,\ldots, averaging the
conjugated representation matrix commutes with the irreducible SU(2)
representation. It is scalar by Schur's lemma, and its trace fixes that
scalar to the character divided by the dimension. The weights
m=-j,-j+1,\ldots,j give \[
 c_j(s)=\frac{1}{2j+1}\sum_{m=-j}^j e^{2ims}
 =\frac{\sin((2j+1)s)}{(2j+1)\sin s}.
\tag{4.3}
\] The sum is the definition at sin s=0 and gives its continuous value;
no singular denominator is used there. Since the product coefficients
are a complete orthogonal basis, (4.3) defines K\_theta on the entire
Hilbert space. It commutes with each left and right link translation,
with Pi, and with the full electric energy.

Thus (2.2)--(2.3) become the exact full-vacuum formulas \[
 c_\theta=\langle\psi,K_\theta\psi\rangle,\quad
 d_\theta=\langle\psi,K_\theta^2\psi\rangle-c_\theta^2,
\quad n_\theta=\mathfrak q_{\mathcal A}[K_\theta\psi].
\tag{4.4}
\] On smooth psi the numerator also has the exact double-commutator form
\[
 n_\theta=\frac12\langle\psi,
 [K_\theta,[H,K_\theta]]\psi\rangle
 =\frac12\langle\psi,[K_\theta,[V,K_\theta]]\psi\rangle.
\tag{4.5}
\] Expanding the first double commutator gives
2KHK-K\textsuperscript{2H-HK}2, and applying H psi=E psi proves the
first equality. The second uses the proved commutation with the electric
Laplacian. No magnetic term disappears: V is still the sum over all M
faces.

\subsection{5. Leading physical state for the magnetic angle, at every
b\textgreater0}\label{magnetic-part-1-leading-physical-state-for-the-magnetic-angle-at-every-b0}

In (3.2) the native generator has \(\alpha_{(n,1),3}=2n_2\) and all
other components zero. Indeed \(-n_2H_c=2n_2T_3\) in h\_1=exp(-theta
n\_2 H\_c), and T\_theta uses the inverse translation exp(-theta Y).
Since a source has only one such edge, (3.2) yields \[
 K_2=\frac43\sum_{e=(n,1)}n_2^2\Delta_e^T
 =-\frac43\Gamma,
\qquad \Gamma=\sum_{e=(n,1)}n_2^2(-\Delta_e^T).
\tag{5.1}
\] The same coefficient follows from the retained character sum:
symmetry gives sum m=0 and \(\sum_{m=-j}^jm^2=(2j+1)j(j+1)/3\), so
\(c_j(s)=1-2j(j+1)s^2/3+O_j(s^4)\). For the full operator the correct
expansion is on smooth vectors, or as a map H\^{}\{k+4\} to H\^{}k: \[
 K_\theta\psi=\psi-\frac23\theta^2\Gamma\psi
                  +O_{H^k}(\theta^4).
\tag{5.2}
\] One need not sum unbounded representation-wise errors to justify it.
Apply Taylor's integral formula to the averaged link-translation flow:
four derivatives are bounded H\^{}\{k+4\} to H\^{}k, uniformly over the
compact averaging group and small theta. Evenness of each C, hence of
their product, removes the odd terms. There is no same-Sobolev-space
bounded operator expansion with the unbounded Gamma silently treated as
bounded.

Put \[
 \gamma=\langle\psi,\Gamma\psi\rangle,\qquad
 v=(\Gamma-\gamma)\psi.
\] Equations (4.4) and (5.2), with the actual vacuum scalar retained,
give \[
 \chi_\theta=-\frac23\theta^2v+O_{H^1}(\theta^4),
\quad d_\theta=\frac49\theta^4\|v\|^2+O(\theta^6),
\quad n_\theta=\frac49\theta^4\mathfrak q_{\mathcal A}[v]
                         +O(\theta^6).
\tag{5.3}
\] In particular \(\|v\|^2=\langle\psi,\Gamma^2\psi\rangle-\gamma^2\) is
the actual quantum variance of this specified weighted electric
operator, not a variance in an independently postulated link measure.

For this native translation v is nonzero for every b\textgreater0. To
prove it, suppose Gamma psi=gamma psi. Integrate against Haar, not psi:
every Laplacian has zero Haar integral, and \(\int\psi>0\), so gamma=0.
Taking its inner product with psi yields \[
 0=\langle\psi,\Gamma\psi\rangle
   =\sum_{e=(n,1),a}n_2^2\|X_{e,a}\psi\|^2.
\] Thus psi is independent of each direction-1 link with n\_2 nonzero;
the three derivative fields span that connected SU(2) factor. Select
e=((0,1,0),1). It exists in every box L\textgreater=2 and has exactly
four incident faces. In the eigen-equation, after this independence is
proved, all kinetic terms and the factor psi are independent of U\_e.
Fix all other link variables to I; smoothness permits this pointwise
evaluation and strict positivity makes the fixed value of psi nonzero.
The potential as a function of U\_e is precisely \[
 V(U_e,I_{\ne e})=4b(2-\operatorname{tr}U_e).
\] Nonincident faces have trace 2; every incident face has trace tr U\_e
or tr U\_e\^{}\{-1\}, which are equal in SU(2). This nonconstant
potential contradicts the eigen-equation for b\textgreater0. We have
proved v!=0, without a small-b assumption and without calculating the
full vacuum by Haar substitution.

Consequently for every fixed box, a, and g\_YM\textgreater0, \[
 \lim_{\theta\to0}\frac{n_\theta}{d_\theta}
 =\frac{\mathfrak q_{\mathcal A}[v]}{\|v\|^2}>0.
\tag{5.4}
\] The original magnetic angle and scale are still theta=2pi a\^{}2/D.
Both numerator and norm tend to zero at the same fourth order, so a
small classical magnetic action does not by itself determine this
physical Rayleigh quotient. Formula (5.4) does not bound that quotient
uniformly when L, a, g\_YM and D vary together. Those joint dependences
remain to be calculated rather than inferred from the fixed-box limit.

\subsection{6. Exact derivatives retaining every magnetic
plaquette}\label{magnetic-part-1-exact-derivatives-retaining-every-magnetic-plaquette}

For completeness the numerator in (5.4) has no hidden electric-operator
commutator. Gamma commutes with the full electric part of H, and hence
\[
 (H-E)\Gamma\psi=[V,\Gamma]\psi,
\qquad
 [V,\Gamma]F
 =\sum_{e=(n,1),a}n_2^2
       \bigl[(X_{e,a}^2V)F+2(X_{e,a}V)X_{e,a}F\bigr].
\tag{6.1}
\] The formula is the product rule for each original derivative. Thus \[
 \mathfrak q_{\mathcal A}[v]
 =\langle\Gamma\psi,[V,\Gamma]\psi\rangle
 =\frac12\langle\psi,[\Gamma,[V,\Gamma]]\psi\rangle.
\tag{6.2}
\] All functions are smooth, so these unbounded-operator products are
well-defined on the vector being used. Defining
\(q_p=(3/4)\sum_{e\in\partial p,\ e=(n,1)}n_2^2\), the four-edge
fundamental Casimir gives \[
 \Gamma W_p=q_pW_p,\quad
 \Gamma V=-b\sum_pq_pW_p,\quad
 X_{e,a}V=-b\sum_{p\ni e}X_{e,a}W_p.
\tag{6.3}
\] These are exact expressions over the same full face set, including
the 13-plane faces for which the chosen background has identity
curvature. They provide explicit coefficients for calculating the vacuum
moments; they do not assign those moments their free-Haar values.

\subsection{7. Exact map to the discrete electric representation
variables}\label{magnetic-part-1-exact-map-to-the-discrete-electric-representation-variables}

Choose an orthonormal product Peter--Weyl basis with a spin j\_e on each
edge. The electric eigenvalue is \(\kappa\sum_ej_e(j_e+1)\). Each joint
edge-spin subspace is finite dimensional and invariant under vertex
gauge transformations, since left and right multiplication do not change
the edge representation. Applying Pi in that subspace and choosing an
orthonormal basis of its range gives vectors \(\Phi_s\), where the index
s records all edge spins and the invariant multiplicity label. Their
union is a complete orthonormal basis of the physical Hilbert space: the
full Peter--Weyl expansion is complete, Pi is bounded, and its
finite-spin images are dense in its range.

Expand the actual vacuum, without changing it, as \[
 \psi=\sum_s a_s\Phi_s,\quad \sum_s|a_s|^2=1,\qquad
 k_s(\theta)=\prod_{e=(n,1)}
 \frac{\sin((2j_e+1)\theta n_2)}{(2j_e+1)\sin(\theta n_2)}.
\tag{7.1}
\] Each factor has the finite-sum interpretation of (4.3) at a zero
denominator. K acts by k\_s on the entire invariant multiplicity space,
because it already acts by that scalar on the surrounding full joint
edge-spin subspace. Hence the exact vacuum overlap and physical state
norm become \[
 c_\theta=\sum_s|a_s|^2k_s(\theta),\quad
 \chi_\theta=\sum_s a_s(k_s(\theta)-c_\theta)\Phi_s,
\quad
 d_\theta=\sum_s|a_s|^2(k_s(\theta)-c_\theta)^2.
\tag{7.2}
\] The scalar sums converge absolutely since
\textbar k\_s\textbar\textless=1; the vector sum converges in Hilbert
norm. It also converges in the form domain because psi is smooth and K
commutes with the electric Sobolev operator. Thus (7.2) is a unitary
coordinate map, not an identification of the unknown coefficients a\_s
with independent free-field weights.

Retain the entire magnetic matrix \[
 V_{st}=\langle\Phi_s,V\Phi_t\rangle
 =2bM\delta_{st}-b\sum_p\langle\Phi_s,W_p\Phi_t\rangle.
\] For a growing sequence of finite joint-spin cutoffs F\_R tending to
the identity, (4.5) gives the exact convergent expression \[
 n_\theta=-\frac12\lim_{R\to\infty}
 \sum_{s,t\in F_R}\overline{a_s}a_t\,
 (k_s(\theta)-k_t(\theta))^2 V_{st}.
\tag{7.3}
\] Here the notation s in F\_R means basis indices in that cutoff. To
justify the limit, expand the finite sum as the expectation on F\_R psi
of the bounded operator \([K,[V,K]]/2\). Explicitly, the matrix entry of
\([K,[V,K]]\) is \(-(k_s-k_t)^2V_{st}\). Thus the finite sum is the
expectation of \([K,[V,K]]/2\) on F\_R psi. The operator is bounded
because K and V are bounded, and F\_R psi tends to psi. No unproved
absolute convergence of an infinite double matrix sum is required. The
full scalar 2bM has zero contribution in (7.3), since its matrix is
diagonal and the displayed difference then vanishes; it remains in H and
in E.

Equations (7.1)--(7.3) retain the exact relation between the continuous
magnetic angle and all discrete electric channels. The non-Abelian
plaquettes couple those channels through the specified V\_st. Their
individual summands are not assumed nonnegative and are not dropped. A
conclusion about arbitrarily low physical energies requires the limiting
ratio of these two actual-vacuum expressions, not only the smallness of
each change in k\_s or of the classical field strength.

\subsection{8. A finite-range energy kernel and the full covariance
denominator}\label{magnetic-part-1-a-finite-range-energy-kernel-and-the-full-covariance-denominator}

The same state can be expressed directly through local electric
operators, without knowing their vacuum expectations in advance. For
each physical edge let \(E_e=-\Delta_e^T\),
\(\gamma_e=\langle\psi,E_e\psi\rangle\), and \(v_e=(E_e-\gamma_e)\psi\).
These are smooth gauge-invariant vectors orthogonal to psi. Define the
two exact finite matrices \[
 C_{ef}=\langle v_e,v_f\rangle,
 \qquad N_{ef}=\mathfrak q_{\mathcal A}(v_e,v_f).
\tag{8.1}
\] They are positive semidefinite by their Gram definitions, with no
assumption about signs of individual entries. Because psi, E\_e and H
preserve real functions, both matrices are real and symmetric. Different
link Casimirs commute, and an operator commutes with itself, so
\([E_e,E_f]=0\) for every pair. Expanding the double commutator on the
smooth vacuum gives \[
 \boxed{N_{ef}=\frac12\langle\psi,[E_e,[H,E_f]]\psi\rangle
              =\frac12\langle\psi,[E_e,[V,E_f]]\psi\rangle.}
\tag{8.2}
\] Indeed its four terms have expectation
\(\langle E_e\psi,HE_f\psi\rangle+
\langle E_f\psi,HE_e\psi\rangle-
2E\langle E_e\psi,E_f\psi\rangle\). Reality and symmetry make this
exactly twice the centered form in (8.1). The second equality retains
the original kinetic coefficient: every E\_e commutes with its whole
electric sum, so that commutator is identically zero, not approximated
by zero.

This gives an exact finite-range support statement for the numerator: \[
 N_{ef}=\frac b2\sum_{p:\ e,f\in\partial p}
 \langle\psi,[E_e,[2-W_p,E_f]]\psi\rangle.
\tag{8.3}
\] To prove it, if f does not belong to p then E\_f differentiates a
different group variable and commutes with multiplication by W\_p. If f
belongs but e does not, the first commutator differentiates only f and
has coefficients depending only on p's variables; it therefore commutes
with E\_e. These statements follow term by term from the product rule,
including all first-derivative terms. Thus a term can remain only when
both edges belong to the same face, including the diagonal case e=f.~No
vacuum-factorization or correlation-decay assumption is involved.

For the native state set w\_e=n\_2\^{}2 on e=(n,1), and w\_e=0
otherwise. Equations (5.1) and (8.1) imply exactly \[
 v=\sum_e w_ev_e,\qquad
 \|v\|^2=\sum_{e,f}w_ew_fC_{ef},\qquad
 \mathfrak q_{\mathcal A}[v]=\sum_{e,f}w_ew_fN_{ef}.
\tag{8.4}
\] The denominator is positive by Section 5. The small-angle physical
excitation quotient (5.4) is therefore the fully specified ratio \[
 \boxed{\frac{\sum_{e,f}w_ew_fN_{ef}}
              {\sum_{e,f}w_ew_fC_{ef}}.}
\tag{8.5}
\] Its numerator couples only edges sharing an elementary plaquette; its
denominator retains all electric correlations of the actual interacting
vacuum, including separated edges. This is an exact
local-energy/global-covariance relation. It neither asserts that the
denominator grows faster nor rules that out. Estimating these two sums
along the original geometric and coupling sequence is the precise
remaining spectral calculation for this family of physical states.

\subsection{9. Context, sources, and
scope}\label{magnetic-part-1-context-sources-and-scope}

This uses the Hamiltonian formulation of lattice gauge theory of J.
Kogut and L. Susskind, \emph{Physical Review D} 11 (1975), 395--408,
\url{https://doi.org/10.1103/PhysRevD.11.395}. The coefficient
dictionary from the original anisotropic Wilson transfer and the full
compact-group ground-state form are proved in the companion
temporal-transfer and interacting-Wilson sources. Haar projection and
Schur averaging are established compact-group tools; no novelty of those
general tools is claimed here. The particular original background, its
reference-frame map, all orientations, the physical state, and the
nonzero all-b magnetic variance are explicitly calculated above.

The four-dimensional target includes construction of the interacting
quantum theory and a statement about its physical spectrum; see A. Jaffe
and E. Witten, \emph{Quantum Yang--Mills Theory},
\url{https://www.claymath.org/wp-content/uploads/2022/06/yangmills.pdf}.
The present theorem tests an explicit background-to-quantum-state map at
a finite spatial regulator. It neither assumes a uniform volume estimate
nor claims a construction or disproof of the continuum theory. The next
calculation evaluates the actual vacuum coefficients of this state with
all box plaquettes present; its full derivation accompanies this note.

\clearpage
% Complete next live proof; its original equation tags are retained.
\section{Exact gauge-projected magnetic translation in a full Wilson
box}\label{magnetic-part-2-exact-gauge-projected-magnetic-translation-in-a-full-wilson-box}

8 September 2026. All links and all elementary spatial plaquettes are
retained. The vacuum below is the actual normalized positive ground
state, not the one-time Wilson density. This proof is a finite-box
calculation; the geometric diagonal in Section 9 concerns its explicitly
computed strong-coupling coefficient, not an interchange of quantum and
volume limits.

\subsection{1. Original objects, constants, and
operators}\label{magnetic-part-2-original-objects-constants-and-operators}

For an integer \(L\ge2\), take the open vertex box
\(\{-L,\ldots,L\}^3\), all positively directed physical links
\(e=(n,i)\) from \(n\) to \(n+\mathbf e_i\) contained in the box, and
all elementary faces \(p=(n;i,j)\), \(i<j\). Write these sets as
\(\mathsf E_L,\mathsf P_L\). Their complete counts are

\[
N=3(2L)(2L+1)^2,\qquad M=3(2L)^2(2L+1).
\tag{1.1}
\]

A negative traversal uses the inverse of the same physical link, not an
independent variable. The face word and its real \(SU(2)\) trace are

\[
U_p=U_i(n)U_j(n+\mathbf e_i)U_i(n+\mathbf e_j)^{-1}U_j(n)^{-1},
\qquad W_p=\operatorname{tr}U_p.
\tag{1.2}
\]

Work in \(\mathcal H=L^2(SU(2)^{\mathsf E_L},dU)\), with product Haar
probability measure and inner product conjugate-linear in its first
entry. Put

\[
T_a=-i\sigma_a/2,\quad -2\operatorname{tr}(T_aT_b)=\delta_{ab},\qquad
X_{e,a}f=\left.\frac{d}{dt}f(\ldots,e^{tT_a}U_e,\ldots)\right|_{t=0},
\quad \Delta_e^T=\sum_aX_{e,a}^2.
\tag{1.3}
\]

For \(\kappa,b>0\), define

\[
H_0=-\sum_{e\in\mathsf E_L}\Delta_e^T,\quad
\mathcal V=\sum_{p\in\mathsf P_L}W_p,\quad \xi=b/\kappa,
\] \[
H_{\kappa,b}=\kappa H_0+b\sum_p(2-W_p)
=\kappa(H_0-\xi\mathcal V)+2bM I.
\tag{1.4}
\]

The exact Wilson constant is \(2bM\). The original metric and physical
parameter dictionary remain

\[
c(X,Y)=-\tfrac12\operatorname{tr}(XY),\quad
\Delta_e^c=4\Delta_e^T,\quad
\kappa=\frac{2g_{\mathrm{YM}}^2}{a},\quad
b=\frac1{2g_{\mathrm{YM}}^2a},\quad
\xi=\frac1{4g_{\mathrm{YM}}^4}.
\tag{1.5}
\]

Thus the kinetic term is also
\(-g_{\mathrm{YM}}^2(2a)^{-1}\sum_e\Delta_e^c\). Nothing is absorbed
into a redefined coupling.

For fixed \(h=(h_e)\in SU(2)^{\mathsf E_L}\), set

\[
(\mathcal T_hf)(U)=f((h_e^{-1}U_e)_e).
\tag{1.6}
\]

This is unitary and preserves Haar integrals. Its adjoint is
\(\mathcal T_{h^{-1}}\). At each vertex use an independent gauge
variable \(g_v\), acting by \[
(\alpha_gU)_e=g_{s(e)}U_eg_{t(e)}^{-1},\qquad
(\Pi f)(U)=\int f(\alpha_gU)\prod_vdg_v.
\tag{1.7}
\]

Haar invariance proves \(\Pi^2=\Pi=\Pi^*\); its range consists precisely
of invariant functions. The original convention
\(U_e\mapsto h_{s(e)}^{-1}U_eh_{t(e)}\) is the same action with vertex
variables \(g_v=h_v^{-1}\); these vertex variables are distinct from the
fixed link assignments in (1.6).

Both \(\mathcal T_h\) and gauge transformations commute with \(H_0\).
For the first assertion, differentiation replaces the basis \(T_a\) by
\(\operatorname{Ad}(h_e^{-1})T_a\), another orthonormal basis, leaving
the sum of squares unchanged. Right multiplication in a gauge
transformation commutes with left derivatives, and its left
multiplication has the same orthonormal-basis argument. Consequently
\(\mathcal T_h\) is an isometry and \(\Pi\) a contraction on every
Sobolev space defined by \(H_0+1\). Since each \(W_p\) is invariant,
\(\Pi\) commutes with \(H_{\kappa,b}\).

\subsection{2. Full vertex integral, with the shared-source factor
retained}\label{magnetic-part-2-full-vertex-integral-with-the-shared-source-factor-retained}

For \(p=(n;i,j)\), abbreviate

\[
\begin{array}{llll}
U=U_i(n),&V=U_j(n+\mathbf e_i),&
Z=U_i(n+\mathbf e_j),&Y=U_j(n),\\
a_0=h_i(n),&d_0=h_j(n+\mathbf e_i),&
c_0=h_i(n+\mathbf e_j),&b_0=h_j(n).
\end{array}
\tag{2.1}
\]

These subscripted matrices are not the scalar couplings or mesh. The
translated word is exactly

\[
\mathcal T_hW_p=
\operatorname{tr}(a_0^{-1}U\,d_0^{-1}V\,Z^{-1}c_0\,Y^{-1}b_0).
\tag{2.2}
\]

In particular, an inverse physical link contributes
\((h_e^{-1}U_e)^{-1}=U_e^{-1}h_e\); its insertion must remain on that
side.

Denote the four vertex variables by \(g_0,g_i,g_{ij},g_j\). Substituting
(1.7) into (2.2), cancelling adjacent factors, and cyclically rotating
the trace gives

\[
\operatorname{tr}\!\left[
(g_0^{-1}b_0a_0^{-1}g_0)\,
U(g_i^{-1}d_0^{-1}g_i)\,VZ^{-1}
(g_j^{-1}c_0g_j)\,Y^{-1}
\right].
\tag{2.3}
\]

The \(g_{ij}\) factors cancel completely. The two links with common
canonical source \(n\) produce the one product \(b_0a_0^{-1}\),
conjugated by the one variable \(g_0\). They cannot be averaged as
independent insertions.

For every complex \(2\times2\) matrix \(C\),

\[
\int g^{-1}Cg\,dg=\tfrac12\operatorname{tr}(C)I.
\tag{2.4}
\]

Indeed the Haar integral commutes with every \(SU(2)\) matrix.
Commutation with all three Pauli one-parameter groups forces it to be
scalar, and taking the trace fixes that scalar. Apply (2.4)
independently to the three remaining variables in (2.3). The exact
answer for every face is

\[
\boxed{\Pi\mathcal T_hW_{n;i,j}=\alpha_{n;i,j}(h)W_{n;i,j},}
\] \[
\boxed{\alpha_{n;i,j}(h)=
\frac{
\operatorname{tr}[h_j(n)h_i(n)^{-1}]
\operatorname{tr}[h_j(n+\mathbf e_i)^{-1}]
\operatorname{tr}[h_i(n+\mathbf e_j)]
}{8}.}
\tag{2.5}
\]

The three traces are real and lie in \([-2,2]\), so
\(\alpha_p\in[-1,1]\). The factor \(2^{-3}\) is one dimension factor for
each nonempty vertex insertion. The fourth vertex contributes 1.

This applies without modification to every coordinate plane and boundary
face. Four vertices and four physical edges are distinct in this open
box. A cyclic change of the starting point leaves the Wilson function
unchanged. Reversing the traversal gives the same function because
\(\operatorname{tr}C^{-1}=\operatorname{tr}C\) on \(SU(2)\). The \(h_e\)
remain attached to canonical physical sources, not to sources invented
by reorienting the traversal.

Two exact checks show what would go wrong with different prescriptions.

\begin{itemize}
\tightlist
\item
  If \(a_0=b_0=k\) and \(c_0=d_0=I\), then \(\alpha_p=1\). Independently
  twirling the two shared-source factors would instead give
  \((\operatorname{tr}k/2)^2\).
\item
  If all four matrices are \(k\), then
  \(\alpha_p=(\operatorname{tr}k/2)^2\), whereas their background face
  word \(kkk^{-1}k^{-1}\) is \(I\). At \(k=i\sigma_3\), the two alleged
  coefficients are 0 and 1.
\end{itemize}

Thus the projected translated face is not generally the original face
times half the trace of a background holonomy.

We shall also need the exact inverse relation

\[
\alpha_p(h^{-1})=\alpha_p(h).
\tag{2.6}
\]

For the first factor, \(\operatorname{tr}(h_j^{-1}h_i)=
\operatorname{tr}[(h_jh_i^{-1})^{-1}]=\operatorname{tr}(h_jh_i^{-1})\),
using cyclicity and the \(SU(2)\) inverse-trace identity. The other two
factors have the same inverse-trace identity separately. This proves
(2.6), including noncommuting assignments.

\subsection{3. The true all-plaquette vacuum and its analytic
expansion}\label{magnetic-part-2-the-true-all-plaquette-vacuum-and-its-analytic-expansion}

The compact product Laplacian has compact resolvent. Adding the smooth
bounded real potential preserves self-adjointness on its operator domain
and compact resolvent. A ground minimizer can be taken nonnegative by
\(|\nabla|f||\le|\nabla f|\), with zeros treated by the approximation
\(\sqrt{|f|^2+\varepsilon^2}\). Elliptic regularity makes the ground
eigenfunction smooth, and the strong maximum principle makes it strictly
positive on the connected configuration manifold. If \(\psi>0\) is a
ground eigenfunction, integration by parts and its eigen-equation give

\[
\langle\psi F,(H-E)\psi F\rangle
=\kappa\sum_{e,a}\int\psi^2|X_{e,a}F|^2\,dU.
\tag{3.1}
\]

The identity extends from smooth \(F\) to the form domain by density;
the positive smooth function \(\psi\) and its reciprocal are bounded on
this finite compact manifold. A second ground eigenfunction divided by
\(\psi\) would make every derivative on the right zero, and thus be
constant. Hence the normalized positive vacuum \(\psi_\xi\) is unique.
Gauge invariance of \(H\) gives \(\Pi\psi_\xi=\psi_\xi\).

Here the needed local expansion follows from the full potential, as can
be proved directly. A fundamental matrix coefficient has free Casimir
\(3/4\) since \(-\sum_aT_a^2=3I/4\). On spin \(j\), the usual matrices
\(J_3v_m=mv_m\), \(J_\pm v_m=\sqrt{(j\mp m)(j\pm m+1)}v_{m\pm1}\) give
Casimir \(j(j+1)\) by multiplication. The Peter--Weyl decomposition and
its tensor products give the full product spectrum: a simple constant
eigenfunction at zero and first nonzero eigenvalue \(3/4\). Each \(W_p\)
uses four distinct fundamental links, so

\[
H_0W_p=3W_p,\qquad
\int W_p\,dU=0,\qquad
\int W_pW_q\,dU=\delta_{pq}.
\tag{3.2}
\]

For the Haar assertions, integration over a singly occurring physical
link kills its fundamental entries: multiplication by the center \(-I\)
changes their signs. Different elementary squares have different
four-edge supports, so one can integrate an edge occurring in one but
not the other. A four-edge support determines the two coordinate
directions and the lower corner of an elementary square, proving the
asserted difference of supports. For equal faces, Haar invariance
reduces the integral to \(\int(\operatorname{tr}U)^2dU=1\): Haar
\(SU(2)\) is the unit sphere in \(\mathbb R^4\), and its scalar
coordinate has second moment \(1/4\).

The interaction norm is exactly \(\|\mathcal V\|_\infty=2M\). The upper
bound follows from \(|W_p|\le2\); equality follows at the all-identity
configuration and then for the essential supremum by continuity and
positive measure of neighborhoods. On \(|z|=3/8\), the free resolvent
norm is at most \(8/3\). Its Neumann series under perturbation by
\(-\xi\mathcal V\) therefore converges for

\[
\boxed{|\xi|<\frac3{16M}.}
\tag{3.3}
\]

The contour projection is analytic and rank one on this disk, since its
rank is locally constant and is one at zero. For real \(\xi\), the
bottom eigenvalue of \(H_0-\xi\mathcal V\) lies in \([-2M|\xi|,0]\),
while the next is at least \(3/4-2M|\xi|>3/8\); the contour thus
contains the actual ground eigenvalue \(e(\xi)\).

A Haar-mean-one analytic eigenvector \(\phi_\xi\) exists locally at
zero. Explicitly, if \(q=16M|\xi|/3\), the contour resolvent difference
bounds the projection difference by \(q/(1-q)\); for \(q<1/2\), the
projection of 1 has nonzero Haar mean and supplies this normalization.
Differentiating the eigen-equation gives

\[
e(0)=e'(0)=0,\qquad \phi'_0=\mathcal V/3,\qquad
\int\phi_\xi^2dU=1+(M/9)\xi^2+O_L(\xi^3).
\tag{3.4}
\]

In the derivative equation the Haar mean forces \(e'(0)=0\), and (3.2)
solves the mean-zero equation \(H_0\phi'_0=\mathcal V\). The squared
normalization follows because all higher coefficients of \(\phi_\xi\)
have zero Haar mean and \(\|\mathcal V\|_2^2=M\).

Analyticity and bounded Taylor remainders hold in every fixed Sobolev
norm. The equation \[
\phi_\xi=(H_0+1)^{-1}[(e(\xi)+1+\xi\mathcal V)\phi_\xi]
\] upgrades analyticity by two derivatives at a time; multiplication by
the fixed smooth \(\mathcal V\) is bounded on each integer Sobolev
space. Sobolev embedding in the fixed dimension \(3N\) also gives
\(C^k\) expansions for every fixed \(k\). Normalize positively for real
\(\xi\). The actual vacuum then has

\[
\psi_\xi=1+\xi A+\xi^2D+O_L(\xi^3),\qquad
A=\frac13\sum_pW_p,\qquad \int D\,dU=-\frac12\|A\|_2^2=-\frac M{18}.
\tag{3.5}
\]

The last identity follows exactly from \(\|\psi_\xi\|_2=1\). The
coefficient \(D\) here is a function, not an assumed truncation of the
interacting vacuum. The physical ground energy is exactly \[
E_\xi=\kappa e(\xi)+2bM.
\tag{3.6}
\]

All constants in these local analytic estimates may depend on the full
box. The disk (3.3) is the guaranteed spectral-projection disk, not a
uniform-volume bound on every Taylor remainder.

\subsection{4. Exact physical state, norm, energy, and vanishing
criterion}\label{magnetic-part-2-exact-physical-state-norm-energy-and-vanishing-criterion}

Define

\[
F_{\xi,h}=\Pi\mathcal T_h\psi_\xi,\quad
c_{\xi,h}=\langle\psi_\xi,\mathcal T_h\psi_\xi\rangle,\quad
\chi_{\xi,h}=F_{\xi,h}-c_{\xi,h}\psi_\xi.
\tag{4.1}
\]

These functions are real; the first two uncentered functions and the
vacuum are strictly positive. Hence \(0<c_{\xi,h}\le1\), the upper bound
following from unitarity. Since \(\Pi\psi_\xi=\psi_\xi\),

\[
\langle\psi_\xi,\chi_{\xi,h}\rangle=0,\qquad
\|\chi_{\xi,h}\|_2^2=\|F_{\xi,h}\|_2^2-c_{\xi,h}^2.
\tag{4.2}
\]

The state is smooth and invariant. Its exact excitation quotient when it
is nonzero is

\[
\boxed{\mathcal Q_{\xi,h}
=\frac{\|(H_{\kappa,b}-E_\xi)^{1/2}\chi_{\xi,h}\|_2^2}
{\|\chi_{\xi,h}\|_2^2}
=\frac{\displaystyle\kappa\sum_{e,a}\int\psi_\xi^2
\left|X_{e,a}\left(\frac{F_{\xi,h}}{\psi_\xi}\right)\right|^2dU}
{\displaystyle\|F_{\xi,h}\|_2^2-c_{\xi,h}^2}.}
\tag{4.3}
\]

Apply (3.1) to \(F_{\xi,h}/\psi_\xi-c_{\xi,h}\) to prove this; the
derivative of the centering constant vanishes. Smoothness and strict
positivity on a compact configuration space justify the form domain and
all displayed integrals. The true vacuum is kept in every factor.

Write the finite real coefficients

\[
\mathcal S_h=\sum_p(1-\alpha_p),\qquad
\mathcal R_h=\sum_p(1-\alpha_p^2),\qquad
\mathcal D_h=\sum_p(1-\alpha_p)^2.
\tag{4.4}
\]

All three are nonnegative and
\(\mathcal D_h=2\mathcal S_h-\mathcal R_h\).

For every \(b>0\), not just perturbatively, the following complete
classification holds:

\[
\boxed{
\chi_{\xi,h}=0
\ \Longleftrightarrow\ \mathcal T_h\psi_\xi=\psi_\xi
\ \Longleftrightarrow\ \mathcal T_h\mathcal V=\mathcal V
\ \Longleftrightarrow\ (\alpha_p=1\text{ for every }p)
\ \Longleftrightarrow\ \mathcal D_h=0.}
\tag{4.5}
\]

If \(\chi=0\), then \(F=c\psi\). Haar integration, preserved by both
\(\Pi\) and \(\mathcal T_h\), yields \(\int\psi=c\int\psi\); strict
positivity forces \(c=1\). Equality
\(\|\Pi\mathcal T_h\psi\|=\|\mathcal T_h\psi\|=1\) in an orthogonal
projection norm implies \((I-\Pi)\mathcal T_h\psi=0\) by Pythagoras.
Thus \(\mathcal T_h\psi=\psi\).

If this last equality holds, apply \(\mathcal T_h\) to the
eigen-equation. It commutes with \(H_0\), and translates multiplication
by \(\mathcal V\) to multiplication by \(\mathcal T_h\mathcal V\).
Subtracting the original equation gives
\(b(\mathcal T_h\mathcal V-\mathcal V)\psi=0\). Since \(b>0\) and
\(\psi>0\), the potentials agree. Project that equality by \(\Pi\), use
(2.5), and pair with \(W_q\) in (3.2); every \(\alpha_q\) is 1.

Conversely, if every \(\alpha_p=1\), then
\(\|\Pi\mathcal T_hW_p\|_2=\|W_p\|_2=\|\mathcal T_hW_p\|_2\). Projection
norm equality gives \(\mathcal T_hW_p=W_p\), in \(L^2\) and hence
pointwise by smoothness. Therefore \(\mathcal T_h\) preserves
\(\mathcal V\), commutes with the full Hamiltonian, and fixes its unique
positive normalized vacuum. This gives \(\chi=0\). A sum of real squares
is zero exactly when each square is zero, completing (4.5). At \(b=0\)
the vacuum is constant and every such centered state is zero; this
explains why \(b>0\) is required in the classification.

There is also a compact exact unprojected energy identity, expressed in
the true moments \(m_{p,\xi}=\int W_p\psi_\xi^2dU\):

\[
\boxed{\langle\mathcal T_h\psi_\xi,
(H_{\kappa,b}-E_\xi)\mathcal T_h\psi_\xi\rangle
=b\sum_p(1-\alpha_p)m_{p,\xi}.}
\tag{4.6}
\]

The kinetic expectation is unchanged by \(\mathcal T_h\), so subtracting
the vacuum energy leaves \(b\sum_p(m_{p,\xi}-\langle\psi_\xi,
(\mathcal T_h^{-1}W_p)\psi_\xi\rangle)\). The density \(\psi_\xi^2\) is
invariant, so the last multiplication function may be averaged by
\(\Pi\). Equations (2.5)--(2.6) replace it exactly by \(\alpha_pW_p\),
proving (4.6). No sign of any individual moment is assumed. The
aggregate is nonnegative because \(H-E_\xi\) is nonnegative.

\subsection{5. First nonzero norm and excitation quotient, with exact
coupling
parity}\label{magnetic-part-2-first-nonzero-norm-and-excitation-quotient-with-exact-coupling-parity}

Equations (2.5) and (3.5) give the actual first-order vector

\[
\chi_{\xi,h}=\xi V_h+O_L(\xi^2),\qquad
V_h=\frac13\sum_p(\alpha_p-1)W_p.
\tag{5.1}
\]

The error is valid in \(H^2\) and every fixed higher Sobolev space. At
this stage its constant may be chosen independently of \(h\) on the
fixed box, because translations are Sobolev isometries and \(\Pi\) a
Sobolev contraction. Equation (3.2) gives

\[
\|V_h\|_2^2=\mathcal D_h/9,\qquad H_0V_h=3V_h.
\tag{5.2}
\]

The second vacuum coefficient is not needed as a function to compute the
scalar centering, but its exact normalization mean is needed. From
(3.5), Haar preservation, and (3.2),

\[
\begin{split}
c_{\xi,h}
&=1+\xi^2\big(\langle A,\mathcal T_hA\rangle-\|A\|_2^2\big)+O_L(\xi^3)
=1-\frac{\xi^2}{9}\mathcal S_h+O_L(\xi^3),\\
\|F_{\xi,h}\|_2^2
&=1+\xi^2\big(\|\Pi\mathcal T_hA\|_2^2-\|A\|_2^2\big)+O_L(\xi^3)
=1-\frac{\xi^2}{9}\mathcal R_h+O_L(\xi^3).
\end{split}
\tag{5.3}
\]

For example, invariance of \(A\) gives
\(\langle A,\mathcal T_hA\rangle=\langle A,\Pi\mathcal T_hA\rangle
=\sum_p\alpha_p/9\). The loss of norm under Gauss projection, measured
here by \(\mathcal R_h\), is not the same as the overlap loss
\(\mathcal S_h\); both enter (4.2).

An exact central link transformation eliminates all odd terms of these
scalar expansions. Define

\[
z_i(n)=(-1)^{\sum_{k<i}n_k}I,\qquad \mathcal Z=\mathcal T_z.
\tag{5.4}
\]

For a face \(i<j\), shifting \(n\) by \(\mathbf e_i\) changes \(z_j\) by
a sign, whereas shifting by \(\mathbf e_j\) does not change \(z_i\).
Multiplying the four link signs around the face therefore gives \(-1\).
Each inverse sign equals the sign itself. Thus \textbf{every} face in
\textbf{every} plane satisfies

\[
\mathcal ZW_p=-W_p,\qquad
\mathcal Z(H_0-\xi\mathcal V)\mathcal Z^{-1}=H_0+\xi\mathcal V.
\tag{5.5}
\]

The central translation \(\mathcal Z\) is a real unitary involution.
Centrality implies it commutes with every gauge transformation and every
\(\mathcal T_h\), hence with \(\Pi\). For the real analytic auxiliary
family near zero, uniqueness of its positive normalized vacuum gives

\[
e(-\xi)=e(\xi),\quad
\psi_{-\xi}=\mathcal Z\psi_\xi,\quad
c_{-\xi,h}=c_{\xi,h},\quad
F_{-\xi,h}=\mathcal ZF_{\xi,h},\quad
\chi_{-\xi,h}=\mathcal Z\chi_{\xi,h}.
\tag{5.6}
\]

Negative \(\xi\) is only an analytic device here; the physical
statements have \(b>0\). The exact excitation operator is
\(\kappa(H_0-\xi\mathcal V-e(\xi))\), since the scalar in (3.6) cancels
exactly. It is conjugated to its negative-\(\xi\) version by
\(\mathcal Z\). Therefore the squared norm of \(\chi\), its excitation
quadratic form, the centering scalar, and the squared norm of the
uncentered projection are even real-analytic scalar functions of
\(\xi\). In particular their cubic coefficients vanish. Equations (4.2),
(5.2), and (5.3) prove

\[
\boxed{\begin{split}
c_{\xi,h}&=1-\frac{\xi^2}{9}\mathcal S_h+O_L(\xi^4),\\
\|F_{\xi,h}\|_2^2&=1-\frac{\xi^2}{9}\mathcal R_h+O_L(\xi^4),\\
\|\chi_{\xi,h}\|_2^2&=\frac{\xi^2}{9}\mathcal D_h+O_L(\xi^4).
\end{split}}
\tag{5.7}
\]

To compute the energy, use the \(H^2\) expansion (5.1) in the exact
excitation operator and \(e(\xi)=O_L(\xi^2)\). The first term is
\(\kappa\xi^2\langle V_h,H_0V_h\rangle\), and evenness removes the cubic
term. Hence

\[
\boxed{\|(H_{\kappa,b}-E_\xi)^{1/2}\chi_{\xi,h}\|_2^2
=\frac{\kappa\xi^2}{3}\mathcal D_h+O_L(\kappa\xi^4).}
\tag{5.8}
\]

For each fixed \(h\) with \(\mathcal D_h>0\), the vector is nonzero for
every \(b>0\) by (4.5). Its first nonzero norm and leading exact
excitation quotient as \(\xi\downarrow0\) are

\[
\boxed{
\|\chi_{\xi,h}\|_2=\frac{\xi}{3}\sqrt{\mathcal D_h}+O_{L,h}(\xi^3),
\qquad \mathcal Q_{\xi,h}=3\kappa+O_{L,h}(\kappa\xi^2).}
\tag{5.9}
\]

The dependence after division must not be discarded near zero
\(\mathcal D_h\). More explicitly, take a fixed-box constant \(C_L\)
bounding both remainders in the last line of (5.7) and in (5.8), after
extraction of \(\xi^4\) and \(\kappa\xi^4\), on a sufficiently small
real interval. Such a constant can be chosen uniformly in \(h\): the
previous Sobolev and analytic bounds hold uniformly for the compact
family of translations, and the spectral expansion is the same vacuum.
If \(C_L\xi^2\le\mathcal D_h/18\), then

\[
|\mathcal Q_{\xi,h}-3\kappa|
\le\frac{72C_L}{\mathcal D_h}\,\kappa\xi^2.
\tag{5.10}
\]

Indeed write the quotient as
\(\kappa(\mathcal D_h/3+\xi^2r_1)/(\mathcal D_h/9+\xi^2r_2)\) with
\(|r_i|\le C_L\). Subtracting \(3\kappa\) gives a numerator of magnitude
at most \(4C_L\kappa\xi^2\); the denominator is at least
\(\mathcal D_h/18\).

In the original physical constants, the leading squared norm is
\(\mathcal D_h/(144g_{\mathrm{YM}}^8)\), the leading unnormalized
excitation form is \(\mathcal D_h/(24g_{\mathrm{YM}}^6a)\), and the
leading quotient is \(6g_{\mathrm{YM}}^2/a\). The exact remainder scale
of the quotient before its dimensionless constant is
\(\kappa\xi^2=1/(8g_{\mathrm{YM}}^6a)\). These are substitutions into
the proved finite-box expressions, not claims about weak coupling or a
continuum limit.

\subsection{6. Native nonwrapping flux: exact coefficients for all three
planes}\label{magnetic-part-2-native-nonwrapping-flux-exact-coefficients-for-all-three-planes}

Now take precisely

\[
h_1(n)=\exp(-\theta n_2H_c),\qquad h_2(n)=h_3(n)=I,\qquad
H_c=i\sigma_3=-2T_3,\qquad \theta\in\mathbb R.
\tag{6.1}
\]

Here \(c(H_c,H_c)=1\) and its squared norm in (1.3) is 4. Its
eigenvalues are \(i,-i\), giving
\(\operatorname{tr}\exp(-tH_c)/2=\cos t\). No integer coordinate has
been wrapped or rescaled, and no factor of 2 has been absorbed into
\(\theta\).

Substitution into the exact full-vertex formula (2.5), with \(m=n_2\),
yields

\[
\boxed{
\begin{array}{ll}
\alpha_{n;1,2}=\cos(m\theta)\cos((m+1)\theta),&-L\le m\le L-1,\\
\alpha_{n;1,3}=\cos^2(m\theta),&-L\le m\le L,\\
\alpha_{n;2,3}=1.&
\end{array}}
\tag{6.2}
\]

The \(13\)-coefficient is not always one, despite the trivial background
holonomy in that plane. The \(12\)-coefficient is not generally the
background normalized trace \(\cos\theta\). Both facts follow from the
actual independent vertex twirls, rather than from a curvature
substitution.

For each permitted \(m\), the number of \(12\)-faces is \(2L(2L+1)\):
\(n_1\) has \(2L\) choices and \(n_3\) has \(2L+1\). The number of
\(13\)-faces for each \(m\) is \(4L^2\). Thus the exact coefficient
entering (5.7)--(5.10) is

\[
\boxed{\begin{split}
\mathcal D_L(\theta)
={}&2L(2L+1)\sum_{m=-L}^{L-1}
[1-\cos(m\theta)\cos((m+1)\theta)]^2\\
&+4L^2\sum_{m=-L}^{L}\sin^4(m\theta).
\end{split}}
\tag{6.3}
\]

The separately relevant overlap coefficient is

\[
\begin{split}
\mathcal S_L(\theta)
={}&2L(2L+1)\sum_{m=-L}^{L-1}
[1-\cos(m\theta)\cos((m+1)\theta)]\\
&+4L^2\sum_{m=-L}^{L}\sin^2(m\theta).
\end{split}
\tag{6.4}
\]

The projection loss is exactly
\(\mathcal R_L=2\mathcal S_L-\mathcal D_L\). The first-order vector
itself is

\[
\begin{split}
V_\theta=\frac13\bigg[
&\sum_{p=(n;1,2)}(\cos(n_2\theta)\cos((n_2+1)\theta)-1)W_p\\
&-\sum_{p=(n;1,3)}\sin^2(n_2\theta)W_p
\bigg],\qquad \chi_{\xi,\theta}=\xi V_\theta+O_L(\xi^2).
\end{split}
\tag{6.5}
\]

Every \(23\)-plaquette is still present in the full potential and actual
vacuum; only its displayed first-order coefficient is zero.

\subsubsection{6.1 Finite Fourier form with no undefined
resonances}\label{magnetic-part-2-finite-fourier-form-with-no-undefined-resonances}

Define the finite trigonometric polynomials

\[
\mathcal O_L(t)=\sum_{m=-L}^{L-1}\cos((2m+1)t),\qquad
\mathcal K_L(t)=\sum_{m=-L}^{L}\cos(mt).
\tag{6.6}
\]

Finite geometric summation gives, whenever their denominators are
nonzero,

\[
\mathcal O_L(t)=\frac{\sin(2Lt)}{\sin t},\qquad
\mathcal K_L(t)=\frac{\sin((L+\tfrac12)t)}{\sin(t/2)}.
\tag{6.7}
\]

At all other arguments, their values are the finite sums (6.6),
equivalently the continuous extensions. In particular
\(\mathcal O_L(k\pi)=2L(-1)^k\) and \(\mathcal K_L(2k\pi)=2L+1\).

With \(a_\theta=1-\tfrac12\cos\theta\), product-to-sum and the identity
\(\sin^4x=(3-4\cos2x+\cos4x)/8\) give exactly

\[
\begin{split}
\sum_{m=-L}^{L-1}[1-\cos(m\theta)\cos((m+1)\theta)]^2
&=2L(a_\theta^2+\tfrac18)-a_\theta\mathcal O_L(\theta)
+\tfrac18\mathcal O_L(2\theta),\\
\sum_{m=-L}^{L}\sin^4(m\theta)
&=\tfrac38(2L+1)-\tfrac12\mathcal K_L(2\theta)
+\tfrac18\mathcal K_L(4\theta).
\end{split}
\tag{6.8}
\]

Multiplying by the respective face counts in (6.3) gives the exact
closed Fourier form of \(\mathcal D_L\). The two unsquared sums in (6.4)
are respectively

\[
2La_\theta-\tfrac12\mathcal O_L(\theta),\qquad
\tfrac12(2L+1)-\tfrac12\mathcal K_L(2\theta).
\tag{6.9}
\]

These expressions are finite polynomial identities and remain valid at
every resonance with (6.6); they introduce no singular angles.

\subsubsection{6.2 Exact zero angles, special values, and elementary
bounds}\label{magnetic-part-2-exact-zero-angles-special-values-and-elementary-bounds}

The coefficient \(\mathcal D_L\) is even and \(2\pi\)-periodic. Its
\(m=0,-1\) terms in the first sum give

\[
4L(2L+1)(1-\cos\theta)^2\le\mathcal D_L(\theta).
\tag{6.10}
\]

Each term in that sum is at most 4 and each sine fourth power at most 1,
giving the nonoptimal but explicit full-box upper bound

\[
\mathcal D_L(\theta)\le5(2L)^2(2L+1)=5M/3.
\tag{6.11}
\]

At \(\theta\in2\pi\mathbb Z\), every \(h_e=I\). Together with (6.10),
this proves

\[
\mathcal D_L(\theta)=0\ \Longleftrightarrow\ \theta\in2\pi\mathbb Z.
\tag{6.12}
\]

The exact all-coupling classification (4.5) therefore gives

\[
\boxed{\chi_{\xi,\theta}=0\ \Longleftrightarrow\
\theta\in2\pi\mathbb Z,\qquad b>0.}
\tag{6.13}
\]

At those angles the quotient of this zero vector is undefined, not zero.
At every other angle it is the quotient of a nonzero, invariant,
vacuum-orthogonal physical vector.

Two explicit nonzero values are

\[
\mathcal D_L(\pi)=16L^2(2L+1),\qquad
\mathcal D_L(\pi/2)=4L^2(3L+1+\varepsilon_L),\quad
\varepsilon_L=L\bmod2.
\tag{6.14}
\]

At \(\pi\), every \(12\)-coefficient is \(-1\) and the \(13,23\)
coefficients are 1, proving the first value by counting faces. At
\(\pi/2\), consecutive cosine factors have zero product, giving
contribution 1 from every \(12\)-face. Among \(-L,\ldots,L\) there are
\(L+\varepsilon_L\) odd integers, which are exactly the nonzero
fourth-power sine terms. This proves the second value with its parity
dependence.

\subsection{7. Exact small-angle polynomial and a signed explicit
remainder}\label{magnetic-part-2-exact-small-angle-polynomial-and-a-signed-explicit-remainder}

For each integer \(m\), put \(r_m=m^2+m+\tfrac12\). Elementary Taylor
expansion gives

\[
1-\cos(m\theta)\cos((m+1)\theta)=r_m\theta^2+O_L(\theta^4),
\qquad \sin^2(m\theta)=m^2\theta^2+O_L(\theta^4)
\tag{7.1}
\]

uniformly over the finite index ranges of the fixed box. Thus

\[
\mathcal D_L(\theta)=\mathcal A_L\theta^4+O_L(\theta^6),
\] \[
\begin{split}
\mathcal A_L
&=2L(2L+1)\sum_{m=-L}^{L-1}r_m^2+4L^2\sum_{m=-L}^{L}m^4\\
&=\boxed{\frac{L^2(2L+1)}{15}
(24L^4+24L^3+8L^2-4L+3)}\\
&=\frac{16}{5}L^7+\frac{24}{5}L^6+\frac83L^5+\frac2{15}L^3+\frac15L^2.
\end{split}
\tag{7.2}
\]

Here are the exact finite sums needed to verify every coefficient:

\[
\sum_{m=-L}^{L-1}r_m^2=\frac{L(4L^4+1)}{10},\qquad
\sum_{m=-L}^{L}m^4
=\frac{L(L+1)(2L+1)(3L^2+3L-1)}{15}.
\tag{7.3}
\]

For the first identity, expand \(r_m^2=m^4+2m^3+2m^2+m+\tfrac14\). Each
polynomial identity in (7.3) is proved directly by its value at \(L=1\)
and by subtracting its value at \(L\) from its value at \(L+1\): the
left side adds exactly the two new endpoint terms, and expansion of the
displayed right side gives those same terms. Substituting (7.3) and
multiplying gives (7.2), so its degree-seven dependence has not been
hidden in a remainder.

The unsquared overlap coefficient has the distinct exact leading
polynomial

\[
\mathcal S_L(\theta)=\mathcal B_L\theta^2+O_L(\theta^4),\qquad
\mathcal B_L=\frac{2L^2(2L+1)(4L^2+2L+1)}3.
\tag{7.4}
\]

In fact \(\sum_{m=-L}^{L-1}r_m=L(2L^2+1)/3\) and
\(\sum_{m=-L}^{L}m^2=L(L+1)(2L+1)/3\); the same endpoint-increment proof
establishes both identities. Multiplication by the respective face
counts in (6.4) proves (7.4).

We now replace the unspecified remainder in (7.2) by an explicit bound.
For every real \(t\),

\[
0\le\tfrac12t^2-(1-\cos t)\le t^4/24.
\tag{7.5}
\]

The lower bound is \(1-\cos t=2\sin^2(t/2)\le t^2/2\); the upper bound
is Taylor's theorem with the fourth derivative of cosine bounded by 1.
For real \(x,y\), set \[
u=1-\cos x\cos y,\qquad q=(x^2+y^2)/2.
\] Using \(u=\tfrac12[(1-\cos(x+y))+(1-\cos(x-y))]\) in (7.5) gives

\[
0\le u\le q,\quad
0\le q-u\le\frac{(x+y)^4+(x-y)^4}{48}
=\frac{x^4+6x^2y^2+y^4}{24}\le q^2/3.
\tag{7.6}
\]

The last inequality is equivalent to \((x^2-y^2)^2\ge0\). It follows
that

\[
0\le q^2-u^2=(q-u)(q+u)\le\tfrac23q^3.
\tag{7.7}
\]

Similarly, with \(u=\sin^2x=(1-\cos2x)/2\) and \(q=x^2\), (7.5) gives
\(0\le q-u\le q^2/3\), hence (7.7) again. Substituting
\(x=m\theta,y=(m+1)\theta\) in the first case and \(x=m\theta\) in the
second proves the signed exact bound

\[
\boxed{0\le\mathcal A_L\theta^4-\mathcal D_L(\theta)
\le\mathcal E_L|\theta|^6,}
\] \[
\mathcal E_L=\frac23\left[
2L(2L+1)\sum_{m=-L}^{L-1}r_m^3+
4L^2\sum_{m=-L}^{L}m^6\right].
\tag{7.8}
\]

This proof is valid for every real \(\theta\), in particular for the
requested regime \(|\theta|(L+1)\le1\).

The finite polynomial sum in (7.8) can be written exactly as

\[
\mathcal E_L=
\frac{L^2(2L+1)}{315}
(240L^6+360L^5+156L^4-120L^3-10L^2+40L-1).
\tag{7.9}
\]

For explicit verification, the two sums before multiplication are

\[
\begin{split}
\sum_{m=-L}^{L-1}r_m^3
&=\frac27L^7-\frac15L^5+\frac16L^3-\frac1{420}L,\\
\sum_{m=-L}^{L}m^6
&=\frac{L(L+1)(2L+1)(3L^4+6L^3-3L+1)}{21}.
\end{split}
\tag{7.10}
\]

The first polynomial has value \(1/4\) at \(L=1\), and its increment
from \(L\) to \(L+1\) is \(2(L^2+L+\tfrac12)^3\), exactly the new two
terms. The second has value 2 at \(L=1\) and increment \(2(L+1)^6\).
This proves both formulas by induction, and substitution into (7.8)
gives (7.9).

A simpler upper bound needs no signed polynomial comparison. On
\(-L\le m\le L-1\), \(0<r_m\le L^2\); on \(-L\le m\le L\),
\(m^6\le L^6\). Thus \[
\begin{split}
\mathcal E_L
&\le\frac23\big[2L(2L+1)(2L)L^6+
4L^2(2L+1)L^6\big]\\
&=\frac{16}{3}L^8(2L+1)\le16L^9
\end{split}
\] for \(L\ge1\). We have established the fully explicit form

\[
\boxed{|\mathcal D_L(\theta)-\mathcal A_L\theta^4|
\le16L^9|\theta|^6,\qquad L\ge2,\ \theta\in\mathbb R,}
\tag{7.11}
\] with the stronger sign information in (7.8).

\subsection{8. The order of the finite-box quantum and angular
statements}\label{magnetic-part-2-the-order-of-the-finite-box-quantum-and-angular-statements}

For each fixed box and angle, the proved actual-vacuum expansion gives

\[
\lim_{\xi\downarrow0}\frac{\|\chi_{\xi,\theta}\|_2^2}{\xi^2}
=\frac{\mathcal D_L(\theta)}9.
\tag{8.1}
\]

Combining it with (7.2) gives the iterated, and therefore unambiguous,
small-angle statement

\[
\lim_{\theta\to0}\frac1{\theta^4}
\left(\lim_{\xi\downarrow0}\frac{\|\chi_{\xi,\theta}\|_2^2}{\xi^2}\right)
=\frac{\mathcal A_L}{9}.
\tag{8.2}
\]

For every fixed \(\theta\notin2\pi\mathbb Z\), the quotient tends to
\(3\kappa\) as \(\xi\downarrow0\), by (5.9). Its leading angular
coefficient cancels because the entire first-order vector lies in the
proved free eigenvalue-3 eigenspace (5.2). The coefficient has not
vanished from the norm, the overlap, or the explicit
denominator-sensitive error estimate (5.10).

The smallness of a centered state's norm is not by itself smallness of
the energy quotient after normalization. In particular (8.2) is not a
zero-gap conclusion. No fixed-\(b\) small-angle theorem is inferred by
interchanging the limits in (8.2); the parent operator calculation
treats that separate order directly.

\subsection{9. Original admissible cusp diagonal: exact coefficient
scaling}\label{magnetic-part-2-original-admissible-cusp-diagonal-exact-coefficient-scaling}

Retain the original geometric sequence, with no alteration of its units
or integer coordinates:

\[
L_j=j^2,\qquad a_j=\frac1{100j},\qquad
T_j^{\mathrm{cusp}}=j^2,\qquad M_j^{\mathrm{cover}}=j,
\qquad D_j^{\mathrm{geom}}=j^4(1+O(j^{-2})).
\tag{9.1}
\]

The determinant \(D_j^{\mathrm{geom}}\) and its asymptotic in (9.1) are
the original retained-cusp geometric data from the parent calculation,
not the number \(M\) of lattice faces and not the function
\(\mathcal D_L\). Their map to the native link parameter is

\[
\theta_j=\frac{2\pi a_j^2}{D_j^{\mathrm{geom}}}
=Cj^{-6}(1+O(j^{-2})),\qquad C=\frac{2\pi}{10^4}.
\tag{9.2}
\]

The second equality follows by substituting \(a_j^2=1/(10^4j^2)\) and
inverting \(1+O(j^{-2})\), which is bounded away from zero for all
sufficiently large \(j\). Consequently \(|\theta_j|(L_j+1)=O(j^{-4})\),
so the requested small-angle box regime holds eventually.

The exact polynomial (7.2) gives

\[
\mathcal A_{j^2}=\frac{16}{5}j^{14}(1+O(j^{-2})).
\tag{9.3}
\]

Meanwhile (7.11), with its explicit constant rather than a box-dependent
Taylor symbol, gives

\[
0\le j^{10}\big(\mathcal A_{j^2}\theta_j^4
-\mathcal D_{j^2}(\theta_j)\big)
\le16j^{28}|\theta_j|^6=O(j^{-8}).
\tag{9.4}
\]

The fourth power in (9.2) and (9.3) gives
\(j^{10}\mathcal A_{j^2}\theta_j^4
=(16/5)C^4(1+O(j^{-2}))\). Squeezing with (9.4) proves the exact
geometric coefficient limit

\[
\boxed{\lim_{j\to\infty}j^{10}\mathcal D_{j^2}(\theta_j)
=\frac{16}{5}\left(\frac{2\pi}{10^4}\right)^4.}
\tag{9.5}
\]

In particular the sequential true-vacuum coefficient has \[
\lim_{j\to\infty}j^{10}
\left(\lim_{\xi\downarrow0}
\frac{\|\chi_{\xi,\theta_j}\|_2^2}{\xi^2}\right)
=\frac{16}{45}\left(\frac{2\pi}{10^4}\right)^4.
\tag{9.6}
\] The inner limit in (9.6) is taken separately on each finite box
\(L_j\), exactly as in (8.1). Equation (9.6) is not a simultaneous
quantum limit.

Indeed the number of all faces on this diagonal is \[
M_j^{\mathrm{faces}}=12j^4(2j^2+1),
\] so the explicit spectral-projection disk (3.3) is exactly

\[
|\xi|<\frac1{64j^4(2j^2+1)}
\sim\frac1{128j^6}.
\tag{9.7}
\]

The Sobolev/Taylor constants \(C_{L_j}\) in (5.7)--(5.10) also remain
box-dependent, and (5.10) further divides by
\(\mathcal D_{j^2}(\theta_j)\), which is of order \(j^{-10}\). These
facts are displayed to prevent use of an unproved uniform remainder.

For the original mesh and fixed \(g_{\mathrm{YM}}>0\), \[
\kappa_j=200g_{\mathrm{YM}}^2j,\qquad
b_j=\frac{50j}{g_{\mathrm{YM}}^2},\qquad
\xi_j=\frac1{4g_{\mathrm{YM}}^4},\qquad
3\kappa_j=600g_{\mathrm{YM}}^2j.
\tag{9.8}
\]

The ratio \(\xi_j\) is fixed, not shrinking, and eventually exceeds the
sufficient disk (9.7). Consequently neither the leading quotient
\(3\kappa_j\) nor its fixed-box remainder justifies a
fixed-\(g_{\mathrm{YM}}\), growing-box quantum limit. Equations
(9.5)--(9.6) concern the explicit coefficient and the displayed order of
limits only. They neither prove nor refute a four-dimensional continuum
gap.

\subsection{10. Exact verification and
provenance}\label{magnetic-part-2-exact-verification-and-provenance}

The companion checker, check\_magnetic\_translation.py, uses integer and
rational quaternion arithmetic. The four conjugations by
\(I,-i\sigma_1,-i\sigma_2,-i\sigma_3\) average every \(2\times2\) matrix
exactly to its scalar part: each imaginary coefficient has two positive
and two negative conjugation signs, while the scalar coefficient is
fixed. Thus their discrete average is precisely (2.4), not a numerical
Haar sample. After the cancellations in (2.3), each independent vertex
appears in only one such matrix twirl. The checker's full
\(4^4=256\)-term vertex average is therefore exact for the original
plaquette integrand. It evaluates the original translated,
gauge-transformed word before using the projected formula, including
noncommuting rational assignments, shared-source cancellation, and the
explicit failure of the background-holonomy rule.

For \(L=2,\ldots,7\), the checker tests the whole-box angular sums at
the rational-cosine angles \(0,\pi/2,\pi,\pi/3\), the Fourier
identities, special angles, both leading \(L\)-polynomials, and the sign
of the central cochain on every face in every plane. It also tests the
sixth-order error polynomial and its explicit upper bound. These are
exact algebra/incidence diagnostics, not a vacuum simulation or a
continuum computation.

The full-box vacuum construction and metric dictionary agree with
BULK\_PLAQUETTE\_EXTENSION.md; the ground-state form agrees with
spatial\_vacuum\_elimination.tex. All specific definitions and proofs
used in this calculation are written above. Standard compact-group and
elliptic spectral facts are used with the displayed compact state space,
exact Casimir, domain, and bounded potential; no external numerical
claim or unproved marginal closure enters. No remote action, new agent,
PDF generation, or heavy job was used. The result is an exact
full-vertex projection and a controlled finite interacting quantum
excitation calculation, together with the explicit geometric coefficient
limit (9.5).

\clearpage
% Complete next live proof; its original equation tags are retained.
\section{Local energy currents and the first surviving quantum
state}\label{magnetic-part-3-local-energy-currents-and-the-first-surviving-quantum-state}

8 September 2026. This calculation keeps every link and every elementary
plaquette of the interacting three-dimensional spatial Wilson box. Time
is the Hamiltonian time; no fourth spatial coordinate is added. The
weighted electric operator from the magnetic continuation is extended by
the matching magnetic energy, with the extension written explicitly. The
first plaquette contribution then cancels. The next contribution is
calculated below, rather than inferred to have zero norm or zero energy.

\subsection{1. Full Hamiltonian, state space, and weighted
operator}\label{magnetic-part-3-full-hamiltonian-state-space-and-weighted-operator}

Fix an integer \(L\ge2\). Vertices are \(\{-L,\ldots,L\}^3\). A
positively oriented physical link \(e=(n,i)\) exists when both \(n\) and
\(n+\mathbf e_i\) are vertices. Its variable is \(U_e\in SU(2)\);
inverse traversal uses \(U_e^{-1}\). For every elementary face
\(p=(n;i,j)\), \(i<j\), put \[
 W_p=\operatorname{tr}\{U_i(n)U_j(n+\mathbf e_i)
                  U_i(n+\mathbf e_j)^{-1}U_j(n)^{-1}\}.
\tag{1.1}
\] The full link and face sets are denoted \(\mathsf E_L,\mathsf P_L\).
Their cardinalities are \[
 N=3(2L)(2L+1)^2,\qquad M=3(2L)^2(2L+1).
\] All sums below use these sets, including boundary links and faces in
all three planes. On \(\mathcal H=L^2(SU(2)^N,dU)\), with product Haar
probability measure, use \[
 T_a=-i\sigma_a/2,\quad
 X_{e,a}F=\left.\frac{d}{dt}F(\ldots,e^{tT_a}U_e,\ldots)\right|_{t=0},
 \quad E_e=-\sum_{a=1}^3X_{e,a}^2.
\] These conventions give fundamental Casimir \(3/4\). The original
metric \(c(A,B)=-\operatorname{tr}(AB)/2\) has orthonormal basis
\(2T_a\), so \(-\Delta_e^c=4E_e\). Retain \[
 H=\kappa H_0+V,\quad H_0=\sum_e E_e,\quad
 V=\sum_p V_p,\quad V_p=b(2-W_p),
 \quad \kappa=\frac{2g_{\rm YM}^2}{a},\quad
 b=\frac{1}{2g_{\rm YM}^2a},\quad \xi=\frac b\kappa.
\tag{1.2}
\] In particular \(\xi=1/(4g_{\rm YM}^4)\), and the constant \(2bM\) is
part of H. It is not removed from the energy origin.

Gauge transformations act by \(U_e\mapsto g_{s(e)}U_eg_{t(e)}^{-1}\).
Each \(E_e\) is bi-invariant and each \(W_p\) is gauge invariant. Hence
every operator used below preserves the gauge-invariant subspace.

The compact connected configuration space, elliptic kinetic operator,
and bounded smooth V give a compact-resolvent self-adjoint H with domain
\(H^2\) and form domain \(H^1\). Its bottom eigenvector \(\psi>0\) is
smooth and unique after fixing \(\int\psi^2dU=1\). Here is the argument
needed below. Taking the modulus of a form minimizer does not increase
its gradient energy, by approximation with
\(\sqrt{|u|^2+\varepsilon^2}\). Elliptic regularity and the strong
maximum principle make a nonnegative minimizer smooth and strictly
positive. If its eigenvalue is \(\mathcal E\), integration by parts
gives, for smooth F, \[
 \mathfrak q_{H-\mathcal E}[\psi F]
 =\kappa\sum_{e,a}\int |X_{e,a}F|^2\psi^2dU.
\tag{1.3}
\] A second bottom eigenvector divided by \(\psi\) would have every
derivative zero and hence be constant. This proves uniqueness, and gauge
transformations therefore fix \(\psi\). Write
\(\mathcal A=H-\mathcal E\).

For arbitrary real link weights \(f_e\), including signed weights,
define \[
 f_p=\frac14\sum_{e\in\partial p}f_e,\quad
 D_f=\kappa\sum_e f_eE_e+\sum_p f_pV_p,
 \quad \zeta_f=(D_f-\langle\psi,D_f\psi\rangle)\psi.
\tag{1.4}
\] Thus the coefficient \(1/4\) distributes each \emph{same} face energy
among its four boundary links. The magnetic extension is explicit: if
\(\Gamma=\sum_{e=(n,1)}n_2^2E_e\) is the original magnetic-translation
operator, take \(f_{(n,1)}=n_2^2\), \(f_{(n,2)}=f_{(n,3)}=0\). Then
\(D_f=\kappa\Gamma+\sum_p f_pV_p\); it is a new physical trial operator
in the original Hamiltonian, not an assertion that its magnetic term was
already present in \(\Gamma\).

The signed kinetic sum in (1.4) is self-adjoint on its diagonal
Peter--Weyl spectral domain: a coefficient with spins \(j_e\) has
eigenvalue \(\kappa\sum_e f_ej_e(j_e+1)\). Its domain is the set whose
squared coefficients times this eigenvalue squared are summable. This
domain need not equal \(H^2\) when weights are signed. Adding the
bounded real potential in (1.4) preserves self-adjointness on that
domain. All smooth vectors lie in it, and D\_f maps smooth vectors to
smooth vectors. Thus \(\zeta_f\) is a well-defined smooth,
gauge-invariant, vacuum-orthogonal state. Every product and commutator
below is evaluated on smooth vectors; no bounded-operator claim about
D\_f is needed.

\subsection{2. Exact energy current and spectral moments at arbitrary
positive
coupling}\label{magnetic-part-3-exact-energy-current-and-spectral-moments-at-arbitrary-positive-coupling}

The product rule with the above sign of \(E_e\) gives \[
 [E_e,V_p]F=-\sum_a\{(X_{e,a}^2V_p)F
                       +2(X_{e,a}V_p)X_{e,a}F\}.
\tag{2.1}
\] It vanishes when \(e\notin\partial p\). Since all link Casimirs
commute and all potential multiplication operators commute, \[
 [H,D_f]=\kappa\sum_{p}\sum_{e\in\partial p}
                         (f_p-f_e)[E_e,V_p].
\tag{2.2}
\] This is an exact weighted-difference formula, not an expansion in the
background angle or coupling. All derivatives of every V\_p in (2.1)
remain, including the first-derivative terms acting on F.

Define local energy operators and antisymmetric currents by \[
 h_e=\kappa E_e+\frac14\sum_{p\ni e}V_p,\qquad
 J_{ef}=i[h_e,h_f]
       =\frac{i\kappa}{4}\sum_{p:\ e,f\in\partial p}
                         ([E_e,V_p]-[E_f,V_p]).
\tag{2.3}
\] For \(e=f\) both sides of the last formula are zero. Each face is
counted four times in \(\sum_e h_e\), so \(\sum_e h_e=H\). The exact
Heisenberg derivative and its weighted form are therefore \[
 i[H,h_e]=-\sum_fJ_{ef},\qquad
 i[H,D_f]=\frac12\sum_{e,f}(f_f-f_e)J_{ef}.
\tag{2.4}
\] The second equality follows by exchanging e and f in one of the two
terms, using \(J_{ef}=-J_{fe}\). Expanding (2.3) also gives (2.2).
Currents between links that share no face vanish exactly. A constant
weight gives D\_f equal to that constant times H and \(\zeta_f=0\)
exactly, not merely to a computed order.

Let \(d_f=\|\zeta_f\|^2\), \(n_f=\mathfrak q_{\mathcal A}[\zeta_f]\).
The vacuum eigen-equation and expansion of the double commutator give \[
 \mathcal A\zeta_f=[H,D_f]\psi,\qquad
 n_f=\frac12\langle\psi,[D_f,[H,D_f]]\psi\rangle,
 \qquad
 \|\mathcal A\zeta_f\|^2=\|i[H,D_f]\psi\|^2.
\tag{2.5}
\] For example the double commutator expands to
\(2D_fHD_f-D_f^2H-HD_f^2\); its expectation is twice
\(\langle D_f\psi,(H-\mathcal E)D_f\psi\rangle\). There is no change of
vacuum in these formulas.

To specify their spectral content, choose a complete orthonormal
eigenbasis \(u_0=\psi,u_1,u_2,\ldots\) of the physical Hamiltonian, with
\(\mathcal A u_\ell=\lambda_\ell u_\ell\), \(\lambda_\ell>0\) for
\(\ell\ge1\). Set \[
 d\mu_f(\lambda)=\sum_{\ell\ge1}
 |\langle u_\ell,D_f\psi\rangle|^2\delta_{\lambda_\ell}(d\lambda).
\] Completeness, smoothness, and (2.5) prove \[
 d_f=\int d\mu_f,\quad n_f=\int\lambda\,d\mu_f,\quad
 \|i[H,D_f]\psi\|^2=\int\lambda^2d\mu_f.
\tag{2.6}
\] Thus a small current matrix element is accompanied by a specified
state norm, rather than being identified with an excitation energy on
its own.

\subsection{3. The second coefficient of the actual all-plaquette
vacuum}\label{magnetic-part-3-the-second-coefficient-of-the-actual-all-plaquette-vacuum}

Put \(S=\sum_pW_p\). The exact Hamiltonian is
\(H=\kappa(H_0-\xi S)+2bM\). The first nonzero full-Hilbert-space
eigenvalue of H\_0 is \(3/4\). Also \(\|S\|_\infty=2M\): the upper bound
follows from \(|W_p|\le2\), and the all-identity configuration and its
positive-measure neighborhoods give equality. On the circle \(|z|=3/8\),
the free resolvent has norm at most \(8/3\). Factoring \(z-H_0+\xi S\)
therefore yields a convergent Neumann series when \[
 |\xi|<\frac{3}{16M}.
\tag{3.1}
\] The contour spectral projection is analytic and rank one there, by
continuity from \(\xi=0\). For real \(\xi\) in this interval, min--max
bounds place the lowest eigenvalue in \([-2M|\xi|,0]\) and every other
eigenvalue above \(3/4-2M|\xi|>3/8\), so this is the actual ground
projection. Its positive unit-norm eigenvector is real analytic in a
neighborhood of zero. Elliptic regularity applied successively to the
eigen-equation gives analyticity in each fixed Sobolev space.
Neighborhood and remainder constants here may depend on L; no
volume-uniform disk is asserted.

Integrating a single boundary link proves \[
 \int W_p=0,\qquad \int W_pW_q=\delta_{pq},\qquad H_0W_p=3W_p.
\tag{3.2}
\] For distinct faces there is an edge in only one of them;
multiplication of that link by \(-I\) changes the integrand's sign. For
equal faces, Haar invariance reduces the integral to
\(\int_{SU(2)}(\operatorname{tr}U)^2dU=1\). The latter follows from
\(U=u_0I+i\sum_au_a\sigma_a\), uniform on the unit 3-sphere, where
\(\int u_0^2=1/4\). Four fundamental Casimirs give the last identity in
(3.2).

Write the positive unit-norm vacuum and dimensionless bottom energy as
\[
 \psi_\xi=1+\xi A+\xi^2B+O_{H^k,L}(\xi^3),\quad
 A=\frac13S,\quad e(\xi)=-\frac M3\xi^2+O_L(\xi^3).
\] Coefficient comparison in \((H_0-\xi S)\psi_\xi=e(\xi)\psi_\xi\) and
the retained unit-norm condition give \[
 H_0B=\frac13(S^2-M),\qquad \int B=-\frac M{18}.
\tag{3.3}
\] Let \(R_0\) act as zero on constants and as the inverse of H\_0 on
their orthogonal complement. Then B is the completely specified finite
polynomial \(R_0(S^2-M)/3-M/18\). Its fully resolved terms are as
follows.

For an unordered pair of distinct faces sharing an edge e, define \[
 P_{pq,0}=\int_{SU(2)}W_pW_q\,dU_e,\qquad
 P_{pq,1}=W_pW_q-P_{pq,0}.
\tag{3.4}
\] The first expression is a function of the other links. It is one half
of the fundamental trace around the six-edge boundary, with either
orientation giving the same SU(2) trace. To see the coefficient, reverse
one of the face words if necessary so the shared link is oppositely
traversed, and cyclically write the traces as
\(\operatorname{tr}(U_eA)\) and \(\operatorname{tr}(U_e^{-1}B)\). The
fundamental Haar identity
\(\int U_{ij}\overline{U}_{kl}dU=\delta_{ik}\delta_{jl}/2\) gives
\(P_{pq,0}=\operatorname{tr}(AB)/2\). A,B use six distinct outer links.
Their products are independent Haar variables when those links are
integrated. The shared-link tensor product has spins 0 and 1, while
every outer link has spin \(1/2\). Consequently \[
 H_0P_{pq,0}=\frac92P_{pq,0},\quad
 H_0P_{pq,1}=\frac{13}2P_{pq,1},\quad
 \|P_{pq,0}\|^2=\frac14,\quad \|P_{pq,1}\|^2=\frac34.
\tag{3.5}
\] The total squared norm of \(W_pW_q\) is 1 by integrating A and B
first; the Haar projection is orthogonal, proving the last norm.

A repeated face gives \(W_p^2-1\), the spin-1 character, with H\_0
eigenvalue 8. An edge-disjoint pair gives \(W_pW_q\) with H\_0
eigenvalue 6, even if its faces share a vertex. These statements exhaust
the face pairs. Thus no term is hidden in the inverse: \[
 \boxed{B=-\frac M{18}
 +\sum_p\frac{W_p^2-1}{24}
 +\sum_{\{p,q\}:\partial p\cap\partial q=\varnothing}\frac{W_pW_q}{9}
 +\sum_{\{p,q\}:|\partial p\cap\partial q|=1}
       \left(\frac4{27}P_{pq,0}+\frac4{39}P_{pq,1}\right).}
\tag{3.6}
\] Each sum over distinct pairs is unordered. Every independent link of
the full Hamiltonian is still present; (3.6) is a coefficient of its
vacuum, not a replacement by separate two-face Hamiltonians.

\subsection{4. Cancellation and the first surviving state, with all
constants}\label{magnetic-part-3-cancellation-and-the-first-surviving-state-with-all-constants}

Set \(F_f=\sum_pf_p\), \(S_f=\sum_pf_pW_p\), \(E_f=\sum_ef_eE_e\).
Equation (1.4) is exactly \[
 D_f=\kappa\{E_f-\xi S_f+2\xi F_fI\}.
\tag{4.1}
\] The link Casimir gives \(E_fA=S_f\), since
\(E_fW_p=(3/4)\sum_{e\in\partial p}f_eW_p=3f_pW_p\). Thus the
first-order vector of the centered state cancels for \emph{every} chosen
weight profile. Its scalar expectation, retaining its origin, is \[
 \langle\psi_\xi,D_f\psi_\xi\rangle
 =\kappa\left(2\xi F_f-\frac{\xi^2}3F_f\right)+O_L(\kappa\xi^3),
\] because \(\langle A,E_fA\rangle=F_f/3\) and
\(\langle S_f\rangle_{\psi_\xi^2}=2\xi F_f/3+O_L(\xi^2)\). Subtracting
this scalar times the same vacuum gives \[
 \zeta_f=\kappa\xi^2 Z_f+O_{H^k,L,f}(\kappa\xi^3),\qquad
 Z_f=E_fB-S_fA+\frac{F_f}{3}.
\tag{4.2}
\]

Inserting (3.6) resolves every cancellation. For a repeated face,
\(E_f(W_p^2-1)=8f_p(W_p^2-1)\), so its contribution in (4.2) is
\(f_p(W_p^2-1)/3-f_pW_p^2/3+f_p/3=0\). For an edge-disjoint pair,
\(E_f(W_pW_q)=3(f_p+f_q)W_pW_q\); its B term cancels the two
corresponding ordered terms of S\_fA.

For an adjacent pair with shared edge e, all six outer edges have
Casimir \(3/4\) and the shared edge has Casimir 0 or 2. Therefore \[
 E_fP_{pq,0}=\left(3(f_p+f_q)-\frac32f_e\right)P_{pq,0},\qquad
 E_fP_{pq,1}=\left(3(f_p+f_q)+\frac12f_e\right)P_{pq,1}.
\] Define \(\delta_{pq}(f)=f_p+f_q-2f_e\). Subtracting the two terms
from S\_fA leaves the exact finite sum \[
 \boxed{Z_f=\sum_{\{p,q\}:\partial p\cap\partial q=\{e\}}
 \delta_{pq}(f)\left(\frac{P_{pq,0}}9-\frac{P_{pq,1}}{39}\right).}
\tag{4.3}
\] The denominators 9 and 39 arise from the distinct \(9/2\) and
\(13/2\) electric channels, not from dropping their interaction.

Different unordered adjacent pairs are orthogonal in both channels. Here
is a direct proof covering coplanar and hinged pairs. Both channel
functions are odd under \(U_l\mapsto -U_l\) on each of the six outer
boundary edges and even on the shared edge. Two distinct pairs have
different six-edge boundaries: otherwise the symmetric difference of
their two face sets would be a nonempty mod-2 two-cycle with at most
four faces. Choose a face in that cycle. Each of its four boundary edges
needs another face. Two distinct elementary square faces share at most
one edge, so four distinct other faces would be needed, at least five in
total, a contradiction. Multiplication by \(-I\) on an edge in just one
of the boundaries now annihilates the cross integral. The same holds
with H\_0 applied, because of (3.5).

Put \(\mathcal D_f=\sum_{\{p,q\}\ \text{adjacent}}\delta_{pq}(f)^2\).
Using both norms and both electric eigenvalues in (3.5) gives \[
 \boxed{\|Z_f\|^2=\frac{49}{13689}\mathcal D_f,\qquad
 \langle Z_f,H_0Z_f\rangle=\frac{2}{117}\mathcal D_f.}
\tag{4.4}
\] Explicitly the first coefficient is
\(1/(4\cdot9^2)+3/(4\cdot39^2)=49/13689\); the second is
\((9/2)/(4\cdot9^2)+3(13/2)/(4\cdot39^2)=2/117\).

All odd scalar powers vanish. Indeed the central link transformation
\(z_i(n)=(-1)^{\sum_{j<i}n_j}I\) gives an involutive Haar unitary
\(\mathcal Z\) sending \emph{every} W\_p to \(-W_p\), and commuting with
every E\_e. It sends \(H_0-\xi S\) to \(H_0+\xi S\), so uniqueness gives
\(\psi_{-\xi}=\mathcal Z\psi_\xi\). Equation (4.1) gives
\(D_{f,-\xi}=\mathcal ZD_{f,\xi}\mathcal Z-4\kappa\xi F_fI\). This
scalar is exactly removed by centering, and hence
\(\zeta_{f,-\xi}=\mathcal Z\zeta_{f,\xi}\). The centered H transforms in
the same way. Consequently the squared norm and excitation form are even
real-analytic functions locally at zero. Combining this with
(4.2)--(4.4) proves \[
 d_f=\kappa^2\xi^4\frac{49}{13689}\mathcal D_f
                         +O_{L,f}(\kappa^2\xi^6),\qquad
 n_f=\kappa^3\xi^4\frac{2}{117}\mathcal D_f
                         +O_{L,f}(\kappa^3\xi^6).
\tag{4.5}
\] For each fixed profile with \(\mathcal D_f>0\), this proves a nonzero
state for a genuine sufficiently small positive interval of \(\xi\), and
\[
 \frac{n_f}{d_f}=\frac{234}{49}\kappa+O_{L,f}(\kappa\xi^2).
\tag{4.6}
\] The squared norm, not the unsquared norm, is analytic in the
preceding statement. When \(\mathcal D_f=0\), formula (4.3) says exactly
that this coefficient vanishes; it does not by itself assert the full
state vanishes. Constant f was resolved exactly in Section 2.

\subsection{5. An exact spatial profile and its full boundary
count}\label{magnetic-part-3-an-exact-spatial-profile-and-its-full-boundary-count}

Take a fixed coordinate direction r and the link-midpoint weights \[
 f_{(n,i)}=n_r+\frac12\delta_{ir}.
\tag{5.1}
\] These are the original lattice coordinates, without replacing the
signed profile by its absolute value. Their face average is the face
center coordinate \(f_{(n;i,j)}=n_r+(\delta_{ir}+\delta_{jr})/2\). If
physical midpoint length is wanted, it is exactly a times (5.1); then
D\_f and \(\zeta_f\) are multiplied by a and both (4.5) scalars by
\(a^2\). The Rayleigh quotient is unchanged by this stated map.

For an edge of direction i, the centers of two incident faces are
\(x(e)+s\mathbf e_j/2\) and \(x(e)+t\mathbf e_k/2\), where j,k differ
from i and the signs indicate which available side is used. For coplanar
distinct faces j=k and s=-t, so \(\delta_{pq}=0\). For hinged pairs j
differs from k, \[
 \delta_{pq}=\frac12(s\delta_{jr}+t\delta_{kr}).
\] It is zero when i=r and has squared value \(1/4\) when i differs from
r. At a transverse coordinate \(n_j\), the number of available incident
j-faces is \(a_j(n_j)=1\) at either endpoint \(\pm L\) and 2 otherwise.
Its sum over \(-L,\ldots,L\) is \(4L\). Thus the complete number of
hinged pairs sharing direction-i edges is \((2L)(4L)^2=32L^3\). There
are two directions i differing from r. This proves, including every
boundary case, \[
 \mathcal D_f=\frac14\cdot2\cdot32L^3=16L^3.
\tag{5.2}
\] The explicit leading results are therefore \[
 d_f=\frac{784}{13689}\kappa^2\xi^4L^3
                         +O_L(\kappa^2\xi^6),\qquad
 n_f=\frac{32}{117}\kappa^3\xi^4L^3
                         +O_L(\kappa^3\xi^6).
\tag{5.3}
\] The profile's energy expectation is actually zero for every positive
b. Reflection \(n_r\mapsto-n_r\) induces a Haar-preserving link-variable
map; links whose orientation reverses are inverted. It preserves H,
every link Casimir, and the unique positive vacuum, and changes all the
midpoint and face-center weights in (5.1) to their negatives. Thus its
unitary conjugates D\_f to \(-D_f\). Taking the vacuum expectation
proves \(\langle D_f\rangle=0\) exactly. This does not assert the
nonzero-state conclusion uniformly at every coupling.

There is also an exact all-coupling statement for \emph{all} affine
profiles, not only their computed Taylor coefficient. Write \[
 f_e=c+\sum_{r=1}^3t_r x_r(e),\qquad
 x_r((n,i))=n_r+\frac12\delta_{ir},\qquad t_r,c\in\mathbb R.
\] Let \(\zeta_r=D_{x_r}\psi\); the reflection argument just proved its
expectation is zero. Since the constant-weight operator is cH,
\(\zeta_f=\sum_rt_r\zeta_r\) exactly. For every Borel set
\(B\subset(0,\infty)\), let \(P_B=1_B(\mathcal A)\). Reflection in
coordinate r commutes with \(P_B\), changes \(\zeta_r\) to its negative,
and fixes \(\zeta_s\) when \(s\ne r\). Consequently
\(\langle\zeta_r,P_B\zeta_s\rangle=0\) for \(r\ne s\). Coordinate
permutations also preserve the box and H and exchange the three midpoint
profiles, so the three diagonal expectations agree. Thus the
\emph{entire} spectral measures, not only their first moments, obey \[
 \boxed{\mu_f(B)=(t_1^2+t_2^2+t_3^2)\,\mu_{x_1}(B)}
 \qquad\text{for every }b>0.
\tag{5.4}
\] The group-variable inversion needed on an orientation-reversed link
is Haar preserving and preserves its Casimir; reversed face traces equal
the original traces in SU(2). These facts justify both symmetry actions
on the actual Hilbert space. All original affine coefficients are
retained in (5.4). Its content is that varying a slope's direction, size
or constant term within this specified family changes the spectral
weight by the displayed scalar, not its positive-energy support. It
makes no assertion that a non-affine profile has the same measure. In
particular \(\mathcal D_f=16L^3(t_1^2+t_2^2+t_3^2)\), also obtained by
taking the coefficient of \(\xi^4\) in this exact identity.

The physical coefficient dictionary remains \[
 \kappa^2\xi^4=\frac{1}{64g_{\rm YM}^{12}a^2},\quad
 \kappa^3\xi^4=\frac{1}{32g_{\rm YM}^{10}a^3},\quad
 \frac{234}{49}\kappa=\frac{468g_{\rm YM}^2}{49a}.
\tag{5.5}
\] The common \(L^3\) in (5.3) is a computed Taylor coefficient, not a
statement that the all-coupling remainder is uniform in L. The proved
disk (3.1) itself shrinks with M. Sending L to infinity at a fixed
positive coupling cannot be justified by inserting that limit in (5.3).

\subsection{6. Relation to the spectral target and
literature}\label{magnetic-part-3-relation-to-the-spectral-target-and-literature}

Equations (2.2)--(2.6) give the exact energy-current map and spectral
measure for these physical states at arbitrary positive coupling.
Equations (3.6)--(5.5) calculate the actual first surviving coefficient
after the local electric/magnetic cancellation. In this explicit regime
the cancellation produces the two adjacent-face channels, with quotient
\(234\kappa/49\), rather than an arbitrarily small excitation energy. It
does not decide the joint large-volume, small-spacing quantum limit.
That limit requires control of the actual spectral measure (2.6) with
the original coupling and geometry retained; neither the smoothness of a
classical background nor the exact constant-weight zero state supplies
its nonzero low-energy weight.

The Hamiltonian conventions are those of J. Kogut and L. Susskind,
\emph{Hamiltonian formulation of Wilson's lattice gauge theories},
Physical Review D \textbf{11} (1975), 395--408,
\url{https://doi.org/10.1103/PhysRevD.11.395}, with the explicit metric
and coefficient dictionary in (1.2). The use of an observable applied to
the true vacuum and a ratio of its spectral moments belongs to the
variational tradition exemplified by R. P. Feynman, \emph{Atomic Theory
of the Two-Fluid Model of Liquid Helium}, Physical Review \textbf{94}
(1954), 262--277, \url{https://doi.org/10.1103/PhysRev.94.262}. No
helium structure factor or helium dispersion is imported into this gauge
theory.

For compact-graph gauge projection, a primary comparison is B. Bahr and
T. Thiemann, \emph{Gauge-invariant coherent states for loop quantum
gravity. II. Non-Abelian gauge groups}, Classical and Quantum Gravity
\textbf{26} (2009), 045012,
\url{https://doi.org/10.1088/0264-9381/26/4/045012},
\url{https://arxiv.org/abs/0709.4636}. Its product-Haar vertex
projection acts on the same kind of compact link Hilbert space. Here the
seed is the actual interacting Wilson vacuum, not a heat-kernel coherent
state; no gravitational dynamics is identified with H.

The four-dimensional quantum target is the construction and spectral
statement in A. Jaffe and E. Witten, \emph{Quantum Yang--Mills Theory},
\url{https://www.claymath.org/wp-content/uploads/2022/06/yangmills.pdf}.
No continuum construction or mass-gap disproof is claimed in this note.

\clearpage
% Complete next live proof; its original equation tags are retained.
\section{Actual-vacuum electric covariance, the full unsigned incidence
kernel, and its first surviving
states}\label{magnetic-part-4-actual-vacuum-electric-covariance-the-full-unsigned-incidence-kernel-and-its-first-surviving-states}

8 September 2026. The Hamiltonian contains every link and every
elementary plaquette of the full open three-dimensional \(SU(2)\)
spatial box. This proof computes its true-vacuum covariance and energy
matrices, gives the entire finite unsigned incidence kernel, and
calculates the first nonzero interacting vacuum state for every nonzero
real weight in that leading kernel. The kernel is not identified with an
exact all-coupling zero-state space.

\subsection{1. Complete finite model and
domains}\label{magnetic-part-4-complete-finite-model-and-domains}

Fix an integer \(L\ge2\), put \(m=2L\), and let the vertices be
\(n\in\{-L,\ldots,L\}^3\). The physical links are all \(e=(n,i)\),
directed from \(n\) to \(n+\mathbf e_i\) when both endpoints are in the
box. A reversed traversal uses \(U_e^{-1}\). The elementary faces
\(p=(n;i,j)\), \(i<j\), have Wilson functions

\[
W_p=\operatorname{tr}\big[
U_i(n)U_j(n+\mathbf e_i)U_i(n+\mathbf e_j)^{-1}U_j(n)^{-1}
\big].
\tag{1.1}
\]

The full edge and face sets are \(\mathsf E_L,\mathsf P_L\), with

\[
N=3m(m+1)^2,\qquad M=3m^2(m+1).
\tag{1.2}
\]

These counts follow by choosing the direction of a link, or plane of a
face, then its allowed coordinate ranges. There are no periodic
identifications.

Use product Haar probability measure \(dU\) on \(SU(2)^{\mathsf E_L}\),
and

\[
T_a=-i\sigma_a/2,\quad -2\operatorname{tr}(T_aT_b)=\delta_{ab},\qquad
X_{e,a}f=\left.\frac{d}{dt}f(\ldots,e^{tT_a}U_e,\ldots)\right|_{t=0},
\] \[
E_e=-\Delta_e^T=-\sum_aX_{e,a}^2,\quad
H_0=\sum_eE_e,\quad \mathcal V=\sum_pW_p,\quad \xi=b/\kappa,
\] \[
H_{\kappa,b}=\kappa H_0+b\sum_p(2-W_p)
=\kappa(H_0-\xi\mathcal V)+2bM I,\qquad \kappa,b>0.
\tag{1.3}
\]

All \(E_e\) are nonnegative self-adjoint link Casimirs and commute with
one another on the common smooth core and in their joint product
spectral decomposition. The original metric and physical constants
remain

\[
c(X,Y)=-\tfrac12\operatorname{tr}(XY),\quad
\Delta_e^c=4\Delta_e^T,\quad
\kappa=\frac{2g_{\mathrm{YM}}^2}{a},\quad
b=\frac1{2g_{\mathrm{YM}}^2a},\quad
\xi=\frac1{4g_{\mathrm{YM}}^4}.
\tag{1.4}
\]

The exact Wilson scalar \(2bM\) is retained in (1.3).

Let \(\psi_\xi>0\) be the normalized positive vacuum and \(E_\xi\) its
physical ground energy. Compactness and ellipticity give a smooth ground
eigenfunction and compact resolvent; the bounded smooth potential
preserves the free operator domain. The ground minimizer can be taken
nonnegative by the form inequality \(|\nabla|f||\le|\nabla f|\), with
zeros handled by \(\sqrt{|f|^2+\varepsilon^2}\). The strong maximum
principle makes it strictly positive. Integration by parts against its
eigen-equation gives

\[
\mathfrak q_\xi(\psi_\xi F,\psi_\xi G)
=\kappa\sum_{e,a}\int\psi_\xi^2\,
\overline{X_{e,a}F}X_{e,a}G\,dU,
\tag{1.5}
\]

where \(\mathfrak q_\xi\) is the closed form of \(H_{\kappa,b}-E_\xi\).
A second ground state divided by \(\psi_\xi\) would make every
derivative in (1.5) zero, hence would be constant. This proves
uniqueness. Gauge transformations act by
\(U_e\mapsto g_{s(e)}U_eg_{t(e)}^{-1}\); link Casimirs and Wilson traces
intertwine this action. Uniqueness makes \(\psi_\xi\) invariant. The
original convention \(U_e\mapsto h_{s(e)}^{-1}U_eh_{t(e)}\) is exactly
this action with \(g_v=h_v^{-1}\).

Define smooth invariant, vacuum-orthogonal vectors and finite real
matrices by

\[
\gamma_{e,\xi}=\langle\psi_\xi,E_e\psi_\xi\rangle,\qquad
v_{e,\xi}=(E_e-\gamma_{e,\xi})\psi_\xi,
\] \[
C_{ef}(\xi)=\langle v_{e,\xi},v_{f,\xi}\rangle,\qquad
\mathcal N_{ef}(\xi)=\mathfrak q_\xi(v_{e,\xi},v_{f,\xi}).
\tag{1.6}
\]

Both matrices are positive semidefinite and real symmetric by their Gram
definitions and reality of every operator and state. Smoothness ensures
all vectors belong to the physical form and operator domains. All
perturbation remainders below are on the fixed finite box and may depend
on \(L\); no volume-uniform estimate is asserted.

\subsection{2. True-vacuum expansion and the leading
matrices}\label{magnetic-part-4-true-vacuum-expansion-and-the-leading-matrices}

The free single-link spin-\(j\) Casimir is \(j(j+1)\). Its normalization
is fixed by \(-\sum_aT_a^2=3I/4\) on the fundamental representation; the
general value follows from the angular-momentum raising/lowering
matrices. The Peter--Weyl tensor decomposition gives a simple free
constant eigenvector and full-space first nonzero eigenvalue \(3/4\).
Each elementary face uses four distinct fundamental links, so

\[
E_eW_p=\tfrac34\,\mathbf1_{\{e\in\partial p\}}W_p,\qquad
H_0W_p=3W_p,\qquad
\int W_pW_q\,dU=\delta_{pq},\qquad \int W_p\,dU=0.
\tag{2.1}
\]

For the Haar identities, a singly occurring link can be multiplied by
\(-I\) to reverse the sign of its fundamental entries. Distinct
elementary faces have distinct four-edge supports, so one integrates an
edge in only one of the supports. Equal faces reduce to
\(\int(\operatorname{tr}U)^2dU=1\), since the scalar component of a
uniformly distributed unit quaternion has second moment \(1/4\). A
square's two directions and coordinate minima can be recovered from its
edge support, justifying the distinct-support assertion.

The potential norm is exactly \(\|\mathcal V\|_\infty=2M\), by
\(|W_p|\le2\) and the all-identity configuration. The free resolvent on
\(|z|=3/8\) has norm at most \(8/3\), giving the explicit
spectral-projection Neumann disk

\[
|\xi|<\frac3{16M}.
\tag{2.2}
\]

Its contour projection has rank one analytically near zero and
throughout that disk. For real \(\xi\) in it, the lowest eigenvalue of
\(H_0-\xi\mathcal V\) lies in \([-2M|\xi|,0]\), while the next is at
least \(3/4-2M|\xi|>3/8\), so it is the actual ground projection. A
local analytic Haar-mean-one eigenvector \(\phi_\xi\) is obtained by
projecting 1 and dividing by its nonzero Haar mean. Differentiating its
eigen-equation and then normalizing in \(L^2\) gives

\[
\psi_\xi=1+\xi A+\xi^2D+O_L(\xi^3),\qquad
A=\tfrac13\mathcal V,\qquad
\int D\,dU=-M/18.
\tag{2.3}
\]

Here the auxiliary dimensionless eigenvalue has \(e(0)=e'(0)=0\), and
\(E_\xi=\kappa e(\xi)+2bM\) exactly. Analyticity and its remainders hold
in every fixed Sobolev space: from \[
\phi_\xi=(H_0+1)^{-1}[(e(\xi)+1+\xi\mathcal V)\phi_\xi]
\] one gains two derivatives at a time, because multiplication by the
fixed smooth potential is bounded on those spaces. Thus applying any
fixed number of the differential operators below is justified on the
actual analytic vacuum.

Let \(r_e=|\{p:e\in\partial p\}|\), and \[
\mathsf I_{pe}=\mathbf1_{\{e\in\partial p\}},\qquad
t_{ef}=\sum_p\mathsf I_{pe}\mathsf I_{pf}.
\tag{2.4}
\] Thus \(t_{ef}\) is the number of elementary faces containing both
physical edges; \(t_{ee}=r_e\). Equation (2.1) gives

\[
E_eA=\tfrac14\sum_{p\ni e}W_p,\qquad
\gamma_{e,\xi}=\frac{\xi^2}{12}r_e+O_L(\xi^3),\qquad
v_{e,\xi}=\frac{\xi}{4}\sum_{p\ni e}W_p+O_L(\xi^2).
\tag{2.5}
\]

For the expectation coefficient, all terms pairing an \(E_e\)-derivative
against the constant vanish by Haar integration, and
\(\langle A,E_eA\rangle=(3/4)r_e/9=r_e/12\). Equation (2.5) holds in
every fixed sufficiently high Sobolev norm, not merely weakly.

The exact parity that improves all scalar remainders is the central link
cochain

\[
z_i(n)=(-1)^{\sum_{k<i}n_k}I,\qquad
(\mathcal Zf)(U)=f((z_e^{-1}U_e)_e).
\tag{2.6}
\]

For a face \(i<j\), the factor \(z_j(n+\mathbf e_i)\) has the opposite
sign from \(z_j(n)\), while \(z_i(n+\mathbf e_j)=z_i(n)\). Thus the
product of its four center factors is \(-1\), on every face in all three
planes. Therefore \[
\mathcal Z\mathcal V=-\mathcal V,\quad
\mathcal Z(H_0-\xi\mathcal V)\mathcal Z^{-1}=H_0+\xi\mathcal V.
\] The real unitary involution \(\mathcal Z\) commutes with every
\(E_e\). Uniqueness of the positive vacuum gives \[
\psi_{-\xi}=\mathcal Z\psi_\xi,\quad e(-\xi)=e(\xi),\quad
\gamma_{e,-\xi}=\gamma_{e,\xi},\quad
v_{e,-\xi}=\mathcal Zv_{e,\xi}.
\tag{2.7}
\] Negative \(\xi\) is used only in this finite analytic argument. The
physical coupling is positive. Both matrix entries in (1.6) are even
real-analytic functions of \(\xi\), because the excitation operator
\(\kappa(H_0-\xi\mathcal V-e(\xi))\) is conjugated to its
negative-\(\xi\) version by \(\mathcal Z\).

Using (2.1), (2.5), and this evenness now proves, with all constants,

\[
\boxed{
C_{ef}(\xi)=\frac{\xi^2}{16}\,t_{ef}+O_L(\xi^4),\qquad
\mathcal N_{ef}(\xi)=\frac{3\kappa\xi^2}{16}\,t_{ef}
+O_L(\kappa\xi^4).}
\tag{2.8}
\]

The expectation in (2.5) likewise has remainder \(O_L(\xi^4)\). There
are finitely many entries, so the matrix remainder bounds also hold in
any fixed matrix norm, with a box-dependent constant. In particular
\(\mathcal N=3\kappa C+O_L(\kappa\xi^4)\), but this does not justify
taking their quotient along a weight on which the leading matrix
vanishes.

For completeness the numerator has the exact, not perturbative,
finite-range identity

\[
\mathcal N_{ef}
=\tfrac12\langle\psi_\xi,[E_e,[H_{\kappa,b},E_f]]\psi_\xi\rangle
=\frac b2\sum_{p:\,e,f\in\partial p}
\langle\psi_\xi,[E_e,[2-W_p,E_f]]\psi_\xi\rangle.
\tag{2.9}
\]

Expand the double commutator into four terms, use the vacuum equation on
the two endpoint \(H\)'s, and use commutation of \(E_e,E_f\). Reality
and symmetry turn the result into twice (1.6). The electric Hamiltonian
commutes with each \(E_f\). If \(f\notin p\), its derivatives commute
with multiplication by \(W_p\); if \(f\in p\) but \(e\notin p\), the
resulting differential coefficients involve only \(p\)'s variables, and
commute with \(E_e\). This proves the support in (2.9). No comparable
all-coupling finite-range claim is made for \(C\).

\subsection{3. Entire kernel of the unsigned edge-to-face
map}\label{magnetic-part-4-entire-kernel-of-the-unsigned-edge-to-face-map}

For real or complex weights \(w=(w_e)\), define

\[
(\mathsf Iw)_{n;i,j}
=w_i(n)+w_j(n+\mathbf e_i)+w_i(n+\mathbf e_j)+w_j(n).
\tag{3.1}
\]

The leading covariance kernel in (2.8) is exactly \(\ker\mathsf I\),
because \(\sum_{e,f}\overline{w_e}t_{ef}w_f=\sum_p|(\mathsf Iw)_p|^2\).
We now give all such weights, with unique finite parameters.

Let \(\ell=-L\), and put \(\varepsilon(n)=(-1)^{n_1+n_2+n_3}\). Write \[
w_i(n)=\varepsilon(n)a_i(n),\qquad
D_iF(n)=F(n+\mathbf e_i)-F(n).
\tag{3.2}
\] Substituting into (3.1) gives \[
\mathsf Iw=0\quad\Longleftrightarrow\quad
D_ja_i+D_ia_j=0\quad (i\ne j)
\tag{3.3}
\] on every contained face. These are equations on the actual edge
arrays; their differences are taken only where all displayed points
exist.

Choose freely:

\begin{itemize}
\tightlist
\item
  one function \(f_i(x)\) for each \(i\), with \(\ell\le x\le L-1\);
\item
  one function \(\varphi_{ij}(x,y)\) for each \(i<j\), on
  \(\ell\le x,y\le L\), with the two exact anchor conditions \[
  \varphi_{ij}(\ell,y)=0,\qquad \varphi_{ij}(x,\ell)=0.
  \tag{3.4}
  \]
\end{itemize}

Then every kernel weight is represented uniquely by

\[
\boxed{
a_i(n)=f_i(n_i)
+\sum_{j>i}D_i\varphi_{ij}(n_i,n_j)
-\sum_{j<i}D_i\varphi_{ji}(n_j,n_i),\qquad
w_i(n)=\varepsilon(n)a_i(n).}
\tag{3.5}
\]

In the last sum, \(D_i\) differentiates the second argument, namely the
coordinate \(n_i\); in the first it differentiates the first argument.
All parameters have their stated edge or vertex domains, so (3.5) is
defined at every physical link.

First, (3.5) satisfies (3.3). For a pair \(i<j\), all unrelated summands
have zero cross derivative; its two relevant terms are
\(D_jD_i\varphi_{ij}\) and \(-D_iD_j\varphi_{ij}\), which cancel. Thus
every choice of parameters gives a kernel weight.

To prove completeness, take any solution of (3.3). For three distinct
indices \(i,j,k\), the equations and commuting finite differences on
every contained cube give \[
D_kD_ja_i=-D_kD_ia_j
=D_iD_ja_k=-D_jD_ka_i.
\] Therefore \(D_kD_ja_i=0\). For each fixed \(n_i\), the array \(a_i\)
has zero mixed difference in its two transverse coordinates. Summing
that mixed-difference equation across a transverse rectangle gives \[
a_i(n)=f_i(n_i)+u_{ij}(n_i,n_j)+u_{ik}(n_i,n_k),
\tag{3.6}
\] where \[
f_i(x)=a_i(x\mathbf e_i+\ell\mathbf e_j+\ell\mathbf e_k),\qquad
u_{ij}(x,y)=a_i(x\mathbf e_i+y\mathbf e_j+\ell\mathbf e_k)-f_i(x).
\] In particular \(u_{ij}(x,\ell)=0\). The original pair equation
reduces to \[
D_ju_{ij}+D_iu_{ji}=0.
\tag{3.7}
\] For \(i<j\), define the anchored potential explicitly by the finite
sum \[
\varphi_{ij}(x,y)=\sum_{r=\ell}^{x-1}u_{ij}(r,y).
\tag{3.8}
\] The empty sum is zero. This satisfies both (3.4) anchors and
\(D_i\varphi_{ij}=u_{ij}\). Summing (3.7) in the first coordinate gives
\(D_j\varphi_{ij}=-u_{ji}\), because \(u_{ji}(y,\ell)=0\). Equations
(3.6) and (3.8) now give exactly (3.5). They also reconstruct all its
parameters uniquely from \(a\), proving uniqueness.

There are \(3m\) free \(f_i\) values and \(3m^2\) free anchored
potential values. Thus

\[
\boxed{\dim\ker\mathsf I=3m(m+1)=N-M,\qquad
\operatorname{rank}\mathsf I=M.}
\tag{3.9}
\]

In particular this unsigned incidence map is onto the full face array
space. This conclusion follows from the proved parametrization and
rank-nullity, not from a numerical rank test.

A useful explicit family in the kernel is

\[
w_1(n)=(-1)^{n_2+n_3}F(n_1),\qquad w_2(n)=w_3(n)=0
\tag{3.10}
\] for any prescribed function \(F\) on the direction-1 edge coordinate.
Its opposite direction-1 edges cancel on every \(12\) and \(13\) face,
while every \(23\) face has zero weight. In (3.5) it corresponds to
\(f_1(n_1)=(-1)^{n_1}F(n_1)\), other \(f_i=0\), and all potentials zero.

The actual state for a real weight is \[
\Gamma_w=\sum_ew_eE_e,\quad
\gamma_{w,\xi}=\langle\psi_\xi,\Gamma_w\psi_\xi\rangle,\quad
v_{w,\xi}=(\Gamma_w-\gamma_{w,\xi})\psi_\xi=\sum_ew_ev_{e,\xi}.
\tag{3.11}
\] Equations (2.5) and (3.1) give \[
v_{w,\xi}=\frac{\xi}{4}\sum_p(\mathsf Iw)_pW_p+O_{L,w}(\xi^2).
\tag{3.12}
\] For \(w\in\ker\mathsf I\), this merely removes its order-\(\xi\)
state. The next sections compute, rather than assume, what remains.

\subsection{4. Exact adjacent-face modes and the complete second vacuum
coefficient}\label{magnetic-part-4-exact-adjacent-face-modes-and-the-complete-second-vacuum-coefficient}

Let \(\mathscr P_2\) be the set of unordered pairs of distinct
elementary faces sharing an edge. Distinct faces in this open cubical
box share at most one edge. For \(\{p,q\}\in\mathscr P_2\), denote that
edge by \(e=e(p,q)\). Reverse one face traversal if necessary and
cyclically base both words at an endpoint of \(e\), so they have the
form

\[
P_*=U_eA,\qquad Q_*=U_e^{-1}B,\qquad Q'_*=BU_e^{-1}.
\tag{4.1}
\]

This does not change either \(SU(2)\) trace: trace is cyclic and
\(\operatorname{tr}U^{-1}=\operatorname{tr}U\). The outer words \(A,B\)
use three physical links each, disjoint from one another and from \(e\),
with their actual inherited orientations. Define

\[
R_{pq}=\operatorname{tr}(AB),\qquad
D_{pq,0}=W_pW_q,\qquad
D_{pq,1}=D_{pq,0}-\tfrac12R_{pq}.
\tag{4.2}
\]

The first is the six-edge outer boundary Wilson loop. Its trace is
independent of the choice of starting endpoint or total reversal, as is
the definition of \(D_{pq,1}\).

The normalization (1.3) gives the Pauli completeness contraction

\[
\sum_a\operatorname{tr}(T_aC)\operatorname{tr}(T_aD)
=-\tfrac12\operatorname{tr}(CD)+\tfrac14\operatorname{tr}C\operatorname{tr}D.
\tag{4.3}
\]

This follows by multiplying the elementary matrix identity
\(\sum_a(\sigma_a)_{ij}(\sigma_a)_{kl}
=2\delta_{il}\delta_{jk}-\delta_{ij}\delta_{kl}\) by the two matrices'
entries and summing. On the shared edge the two derivatives are
respectively \(\operatorname{tr}(T_aP_*)\) and
\(-\operatorname{tr}(T_aQ'_*)\), since the second occurrence is inverse.
Therefore

\[
G_{pq}:=\sum_a(X_{e,a}W_p)(X_{e,a}W_q)
=\tfrac12R_{pq}-\tfrac14D_{pq,0}
=\tfrac38R_{pq}-\tfrac14D_{pq,1}.
\tag{4.4}
\]

Each outer edge acts with Casimir \(3/4\) on both functions in (4.2),
because its fundamental or inverse fundamental matrix entries occur
linearly. On \(e\), \(R_{pq}\) is constant. The product rule and (4.4)
give \[
\Delta_e^TD_{pq,0}=-\tfrac32D_{pq,0}+2G_{pq}
=R_{pq}-2D_{pq,0}=-2D_{pq,1}.
\] Consequently the joint edge-Casimir labels and full electric
eigenvalues are

\[
\begin{array}{c|ccc|c}
 &\text{six outer edges}&\text{shared edge }e&\text{all other edges}&H_0\\ \hline
R_{pq}&3/4&0&0&9/2\\
D_{pq,1}&3/4&2&0&13/2.
\end{array}
\tag{4.5}
\]

This retains every additional box link as spin zero, rather than
deleting it from the Hamiltonian.

For their norms, condition on \(U_e\). The disjoint outer words \(A,B\)
are independent Haar matrices by left/right Haar invariance and
independence of their link variables. Thus \(P_*,Q'_*\) are independent
Haar matrices. Write them as \(p_0I+i\sum_ap_a\sigma_a\) and
\(q_0I+i\sum_aq_a\sigma_a\). The two coefficient vectors are independent
uniform points on the unit sphere in \(\mathbb R^4\), with second
moments \(\mathbb E p_\alpha p_\beta=\delta_{\alpha\beta}/4\). Matrix
multiplication gives \[
R_{pq}=2(x-y),\qquad D_{pq,1}=3x+y,\qquad
x=p_0q_0,\quad y=\sum_{a=1}^3p_aq_a.
\] Here \(\mathbb E x^2=1/16\), \(\mathbb E y^2=3/16\), and
\(\mathbb E xy=0\). It follows that

\[
\|R_{pq}\|_0^2=1,\qquad
\|D_{pq,1}\|_0^2=3/4,\qquad
\langle R_{pq},D_{pq,1}\rangle_0=0.
\tag{4.6}
\]

All norms are full-box Haar norms; every extra link contributes the
integral of 1.

The modes belonging to different unordered pairs are orthogonal as well.
We give the geometry needed for that claim. The outer boundary of two
coplanar adjacent unit faces is the edge set of a \(2\)-by-\(1\)
rectangle. Its plane, bounding rectangle, and its unique two unit
squares can be read from that edge set. For two perpendicular adjacent
faces, their bounding box is a unit cube. Their outer six-edge boundary
contains exactly six of its vertices; the two absent vertices form the
edge opposite the shared edge. The missing edge and the cube determine
the shared edge and its two incident faces uniquely. The coplanar and
perpendicular cases have respectively two and three nonzero coordinate
spans and cannot coincide. Therefore distinct unordered adjacent pairs
have distinct six-edge outer supports.

Two \(R\)-modes with different outer supports have different eigenvalues
for at least one \(E_f\) (fundamental versus zero), so self-adjointness
makes them orthogonal. An \(R\)-mode and any \(D_1\)-mode are orthogonal
because the latter has shared-edge eigenvalue 2 while the former has
only eigenvalues 0 or \(3/4\). Two \(D_1\)-modes with different shared
edges differ in the eigenvalue-2 location; if the shared edge is equal,
the distinct outer supports give a \(3/4\)-versus-zero location. They
are therefore also orthogonal. This proves full pairwise orthogonality
of the displayed family on the whole box.

We can now write the actual complete second vacuum coefficient \(D\) in
(2.3). The Haar-mean-one vacuum equation gives \[
e(\xi)=-\frac M3\xi^2+O_L(\xi^3),\qquad
H_0\phi_2=\tfrac13(\mathcal V^2-M),\qquad \int\phi_2\,dU=0.
\tag{4.7}
\] The first equality follows by pairing the second-order equation with
1 and using \(\langle\mathcal V,A\rangle=M/3\). The normalized
coefficient is \(D=\phi_2-M/18\).

For a repeated face, \(W_p^2-1\) is the adjoint character: in
eigenvalues \(e^{it},e^{-it}\), it equals \(e^{2it}+1+e^{-2it}\), the
spin-1 character. On each of its four links its Casimir is 2, so its
full eigenvalue is 8. For two faces without a common physical edge, the
product has eight distinct fundamental link factors and full eigenvalue
6, even if the faces share a vertex. For adjacent faces use (4.2),
(4.5). Inverting \(H_0\) on its mean-zero eigenfunctions in (4.7) gives

\[
\boxed{\begin{split}
D={}&-\frac M{18}+\frac1{24}\sum_p(W_p^2-1)\\
&+\frac19\sum_{\substack{\{p,q\}\\p\ne q,\ \partial p\cap\partial q=\varnothing}}W_pW_q\\
&+\sum_{\{p,q\}\in\mathscr P_2}
\left(\frac2{27}R_{pq}+\frac4{39}D_{pq,1}\right).
\end{split}}
\tag{4.8}
\]

The second and third sums are over unordered face pairs; the factor 2
for each distinct pair in \(\mathcal V^2\) is included in their
coefficients. All cases of the full square \((\sum_pW_p)^2\) occur in
(4.8). It is an exact Taylor coefficient of the true vacuum, not a
two-face model for that vacuum.

\subsection{5. First surviving state for every leading-null real
weight}\label{magnetic-part-4-first-surviving-state-for-every-leading-null-real-weight}

Take \(w\in\ker\mathsf I\), real. Then \(\Gamma_w1=0\),
\(\Gamma_wW_p=(3/4)(\mathsf Iw)_pW_p=0\), and hence \(\Gamma_wA=0\). It
follows from the normalized analytic expansion that \[
\gamma_{w,\xi}=O_{L,w}(\xi^4).
\tag{5.1}
\] Indeed every term paired against a constant vanishes by
self-adjointness and \(\Gamma_w1=0\). The order-\(\xi^2\) expectation is
\(\langle A,\Gamma_wA\rangle=0\). The expectation is even by (2.7), so
the cubic coefficient also vanishes.

Every repeated-face term in (4.8) is annihilated by \(\Gamma_w\): all
its used edges have eigenvalue 2, giving \(2(\mathsf Iw)_p=0\). A pair
without a shared edge is annihilated because its weighted Casimir is
\((3/4)[(\mathsf Iw)_p+(\mathsf Iw)_q]=0\).

For an adjacent pair with shared edge \(e\), the sum of its six outer
weights is \[
\sum_{f\text{ outer}}w_f=(\mathsf Iw)_p+(\mathsf Iw)_q-2w_e=-2w_e.
\] The joint Casimirs (4.5) therefore give exactly

\[
\Gamma_wR_{pq}=-\tfrac32w_eR_{pq},\qquad
\Gamma_wD_{pq,1}=\tfrac12w_eD_{pq,1}.
\tag{5.2}
\]

Apply these identities to the full expression (4.8), and then use (5.1)
in the centered state (3.11). The actual first surviving term is

\[
\boxed{
v_{w,\xi}=\xi^2S_w+O_{L,w}(\xi^3),\qquad
S_w=\sum_{\{p,q\}\in\mathscr P_2}
w_{e(p,q)}\left(-\frac19R_{pq}+\frac2{39}D_{pq,1}\right).}
\tag{5.3}
\]

The error holds in every fixed Sobolev norm after applying \(\Gamma_w\),
because the vacuum remainder was established in every sufficiently high
Sobolev norm. Centering affects no term in (5.3).

Let

\[
\mathcal K_L(w,z)=\sum_{e\in\mathsf E_L}\binom{r_e}{2}w_ez_e,\qquad
\mathcal K_L(w)=\mathcal K_L(w,w).
\tag{5.4}
\]

Each unordered adjacent pair has a unique shared edge. Thus summing its
shared-edge weight product is exactly (5.4). The pairwise mode
orthogonality and norms (4.5)--(4.6) give, for real
\(w,z\in\ker\mathsf I\),

\[
\begin{split}
\langle S_w,S_z\rangle_0
&=\left(\frac1{81}+\frac34\frac4{1521}\right)\mathcal K_L(w,z)
=\frac{196}{13689}\mathcal K_L(w,z),\\
\langle S_w,H_0S_z\rangle_0
&=\left(\frac{9/2}{81}+\frac{13}2\frac34\frac4{1521}\right)\mathcal K_L(w,z)
=\frac8{117}\mathcal K_L(w,z).
\end{split}
\tag{5.5}
\]

Every box edge has at least two incident faces: if \(e\) is direction
\(i\), each of its two transverse coordinate directions supplies at
least one face extending inward from the box boundary. Explicitly, \[
r_{(n,i)}=d(n_j)+d(n_k),\qquad
d(x)=\begin{cases}1,&x=-L\text{ or }L,\\2,&-L<x<L,\end{cases}
\quad \{i,j,k\}=\{1,2,3\}.
\tag{5.6}
\] Thus \(2\le r_e\le4\) and \(\mathcal K_L(w)>0\) for every nonzero
real weight. In particular \(S_w\ne0\): all leading-null nonzero weights
have a genuinely surviving second-order state.

Evenness (2.7) applies to every fixed real linear combination \(v_w\)
and to its polarized norm and energy forms. The leading state in (5.3)
is order \(\xi^2\); its scalar products therefore have leading order
\(\xi^4\), and their possible \(\xi^5\) terms vanish. The exact
excitation operator is \(\kappa(H_0-\xi\mathcal V-e(\xi))\). Equations
(5.3)--(5.5) prove the complete restricted leading forms

\[
\boxed{\begin{split}
\langle v_{w,\xi},v_{z,\xi}\rangle
&=\frac{196}{13689}\xi^4\mathcal K_L(w,z)+O_{L,w,z}(\xi^6),\\
\mathfrak q_\xi(v_{w,\xi},v_{z,\xi})
&=\frac{8\kappa}{117}\xi^4\mathcal K_L(w,z)+O_{L,w,z}(\kappa\xi^6).
\end{split}}
\tag{5.7}
\]

For each fixed nonzero \(w\) in the kernel, its squared norm is positive
for all sufficiently small positive \(\xi\), and division yields

\[
\boxed{
\frac{\mathfrak q_\xi(v_{w,\xi},v_{w,\xi})}{\|v_{w,\xi}\|_2^2}
=\frac{234}{49}\kappa+O_{L,w}(\kappa\xi^2).}
\tag{5.8}
\]

This quotient is the energy of the actual smooth invariant,
vacuum-orthogonal vector in (3.11), not of a substituted free vacuum.
Its leading value is between the two displayed free energies
\(9\kappa/2\) and \(13\kappa/2\). The leading incidence kernel raises
the first surviving order and removes the one-face energy \(3\kappa\);
it does not make the normalized energy vanish.

One can make the distinction from an exact zero-state kernel stronger on
each fixed box. Decompose the real edge space as
\((\ker\mathsf I)^\perp\oplus\ker\mathsf I\). On the first summand (2.8)
has a strictly positive \(\xi^2\) leading form, because \(\mathsf I\) is
injective there. On the second, (5.7) has a strictly positive \(\xi^4\)
leading form. Mixed matrix entries have zero \(\xi^2\) term and, by
evenness, zero \(\xi^3\) term, hence are \(O_L(\xi^4)\). All are fixed
finite matrices. If \(u,z\) belong respectively to these summands, their
mixed contribution is bounded by \(B_L\xi^4\|u\|\|z\|\). The elementary
inequality \(2ab\le a^2+b^2\), scaled by the positive \(\xi^2\) lower
bound on the first summand, absorbs it into a portion of that bound plus
\(O_L(\xi^6)\|z\|^2\). For sufficiently small positive \(\xi\), the
latter is smaller than a portion of the positive \(\xi^4\) bound on the
second summand. Thus the entire actual covariance matrix \(C(\xi)\) is
positive definite on the real edge space for all sufficiently small
positive \(\xi\), with a box-dependent interval. No assertion about
arbitrary \(b/\kappa\) is inferred from this local result.

\subsection{6. Explicit alternating family and the original native
weight}\label{magnetic-part-4-explicit-alternating-family-and-the-original-native-weight}

For (3.10), the first surviving term is given explicitly by (5.3). Take
first \(F(n_1)=1\). Only pairs whose shared edge is in direction 1
contribute, with the signed coefficient \((-1)^{n_2+n_3}\). The squared
weight of each is 1. To count them exactly, hold \(n_1\) fixed. If both
transverse coordinates are interior there are \((m-1)^2\) such edges
with \(r_e=4\); if exactly one is on the boundary there are \(4(m-1)\)
with \(r_e=3\); if both are on the boundary there are 4 with \(r_e=2\).
Therefore

\[
\sum_{n_2,n_3}\binom{r_{(n,1)}}2
=6(m-1)^2+12(m-1)+4=6m^2-2.
\] There are \(m\) possible \(n_1\), giving

\[
\boxed{\mathcal K_L(w)=m(6m^2-2)=48L^3-4L.}
\tag{6.1}
\]

For the complete family with arbitrary real \(F\), \[
\mathcal K_L(w)=(6m^2-2)\sum_{n_1=-L}^{L-1}F(n_1)^2.
\tag{6.2}
\] If \(F\ne0\), this is strictly positive. Substitution into (5.3),
(5.7), and (5.8) gives the exact first surviving actual-vacuum state,
its norm and energy coefficients, and its leading quotient for that
specified family, with all faces retained.

The original native weights are instead

\[
w^{\mathrm{nat}}_{(n,1)}=n_2^2,\qquad
w^{\mathrm{nat}}_{(n,2)}=w^{\mathrm{nat}}_{(n,3)}=0.
\tag{6.3}
\]

Their unsigned sums, with \(t=n_2\), are \[
(\mathsf Iw^{\mathrm{nat}})_{n;1,2}=t^2+(t+1)^2=2(t^2+t+\tfrac12),
\] \[
(\mathsf Iw^{\mathrm{nat}})_{n;1,3}=2t^2,\qquad
(\mathsf Iw^{\mathrm{nat}})_{n;2,3}=0.
\tag{6.4}
\] They are not in the leading kernel: in particular the \(12\) sum is
strictly positive for every integer \(t\). Keeping each face
multiplicity gives

\[
\sum_p|(\mathsf Iw^{\mathrm{nat}})_p|^2=4\mathcal A_L,
\] \[
\begin{split}
\mathcal A_L
&=2L(2L+1)\sum_{t=-L}^{L-1}(t^2+t+\tfrac12)^2+
4L^2\sum_{t=-L}^{L}t^4\\
&=\frac{L^2(2L+1)}{15}(24L^4+24L^3+8L^2-4L+3).
\end{split}
\tag{6.5}
\]

The finite sum formulas are
\(\sum_{t=-L}^{L-1}(t^2+t+\tfrac12)^2=L(4L^4+1)/10\) and
\(\sum_{t=-L}^{L}t^4=L(L+1)(2L+1)(3L^2+3L-1)/15\). They follow by
expansion, or by checking \(L=1\) and the two added endpoint terms under
\(L\mapsto L+1\). Substitution gives the polynomial in (6.5).

Equations (2.8), (3.12), and (6.5) therefore yield, for the actual
native state,

\[
\|v_{w^{\mathrm{nat}},\xi}\|^2=\frac{\xi^2}{4}\mathcal A_L+O_L(\xi^4),
\qquad
\mathfrak q_\xi(v_{w^{\mathrm{nat}},\xi},v_{w^{\mathrm{nat}},\xi})
=\frac{3\kappa\xi^2}{4}\mathcal A_L+O_L(\kappa\xi^4),
\tag{6.6}
\] and quotient \(3\kappa+O_L(\kappa\xi^2)\) on this fixed box. This
agrees with the native projected magnetic-translation coefficient: its
fixed-state small-angle leading vector is
\(-\tfrac23\theta^2v_{w^{\mathrm{nat}},\xi}\), so the \(\xi^2\theta^4\)
squared-norm coefficient is \((4/9)(\mathcal A_L/4)=\mathcal A_L/9\).

The exact comparison is therefore: native nonnegative \(n_2^2\) weights
retain the one-face modes at order \(\xi\), whereas transverse
alternating weights cancel those modes and reveal two-face modes at
order \(\xi^2\). This is a proved state-space distinction and does not
assert a favorable infinite-volume scaling of either actual quotient.

\subsection{7. The minimum over all real electric weights, including
coupling-dependent
ones}\label{magnetic-part-4-the-minimum-over-all-real-electric-weights-including-coupling-dependent-ones}

The fixed-weight conclusions do not by themselves control weights
approaching the leading kernel as \(\xi\to0\). We now close that
possibility by inequalities valid simultaneously for every real edge
weight on the fixed box, without changing any physical state,
coefficient, or denominator.

Write the original real edge space as the orthogonal direct sum \[
\mathbb R^{\mathsf E_L}=\mathscr R_L\oplus\mathscr Z_L,\qquad
\mathscr Z_L=\ker\mathsf I,\quad \mathscr R_L=(\ker\mathsf I)^\perp,
\qquad w=u+z.
\tag{7.1}
\] Orthogonality here is the ordinary coordinate inner product on the
finite list of physical edge weights. It is used to state estimates, not
as a replacement for the actual covariance denominator \(w^TC(\xi)w\).

Because \(\mathsf I\) is injective on \(\mathscr R_L\),
finite-dimensional compactness gives the strictly positive constant \[
\lambda_L=\min_{\substack{u\in\mathscr R_L\\\|u\|=1}}\|\mathsf Iu\|^2>0.
\tag{7.2}
\] Set \(r_0=\lambda_L/16\) and \(d_0=196/13689\). The proved matrix
remainders in (2.8) and (5.7), together with evenness of the mixed
block, give on some fixed-box interval \(0<\xi<\epsilon_L\) \[
u^TCu\ge\tfrac12r_0\xi^2\|u\|^2,\qquad
z^TCz\ge\tfrac12d_0\xi^4\|z\|^2,\qquad
|u^TCz|\le B_L\xi^4\|u\|\|z\|,
\tag{7.3}
\] where \(B_L<\infty\). The second inequality uses
\(\mathcal K_L(z)\ge\|z\|^2\) from \(r_e\ge2\). Every bound is between
fixed finite matrix blocks and therefore holds for all their vectors
with one common box-dependent constant.

The elementary inequality \[
2B_L\xi^4\|u\|\|z\|
\le\tfrac14r_0\xi^2\|u\|^2+
\frac{4B_L^2}{r_0}\xi^6\|z\|^2
\tag{7.4}
\] follows by expanding the square of
\(\sqrt{r_0}\xi\|u\|/2-(2B_L/\sqrt{r_0})\xi^3\|z\|\). Shrink the already
nonempty interval so \((4B_L^2/r_0)\xi^2\le d_0/4\); if \(B_L=0\) this
step is unnecessary. Adding the three blocks in (7.3) then proves the
simultaneous physical denominator bound \[
\boxed{w^TC(\xi)w\ge c_L\xi^2\|u\|^2+
\tfrac14d_0\xi^4\|z\|^2,\qquad c_L=r_0/4=\lambda_L/64>0.}
\tag{7.5}
\] In particular the actual centered state is nonzero for every nonzero
real \(w\) on this interval. No component has been scaled by a power of
\(\xi\) in this argument.

Now define the actual matrix difference \[
\mathcal B(\xi)=\kappa^{-1}\mathcal N(\xi)-3C(\xi).
\tag{7.6}
\] Its \(\mathscr R_L\)-block and mixed block are \(O_L(\xi^4)\) by
(2.8). Its restriction to \(\mathscr Z_L\) has the strictly positive
leading form \[
z^T\mathcal Bz=\delta_0\xi^4\mathcal K_L(z)+O_L(\xi^6)\|z\|^2,
\qquad
\delta_0=\frac8{117}-3\frac{196}{13689}
=\frac{348}{13689}=\frac{116}{4563}>0.
\tag{7.7}
\] Hence, on a possibly smaller nonempty interval, there are finite
\(A_L,B'_L\) such that simultaneously \[
u^T\mathcal Bu\ge-A_L\xi^4\|u\|^2,\quad
|u^T\mathcal Bz|\le B'_L\xi^4\|u\|\|z\|,\quad
z^T\mathcal Bz\ge\tfrac12\delta_0\xi^4\|z\|^2.
\tag{7.8}
\] Complete the square in the last two terms, or use \[
\tfrac12\delta_0\|z\|^2-2B'_L\|u\|\|z\|
\ge-\frac{2(B'_L)^2}{\delta_0}\|u\|^2.
\] It follows that \[
w^T\mathcal Bw\ge-K_L\xi^4\|u\|^2,\qquad
K_L=A_L+\frac{2(B'_L)^2}{\delta_0}.
\tag{7.9}
\] The retained physical denominator (7.5) now gives the
uniform-in-weight lower estimate \[
\boxed{
\frac{w^T\mathcal N(\xi)w}{w^TC(\xi)w}
=3\kappa+\kappa\frac{w^T\mathcal B(\xi)w}{w^TC(\xi)w}
\ge3\kappa-\frac{K_L}{c_L}\kappa\xi^2
\quad\text{for every real }w\ne0.}
\tag{7.10}
\] This includes arbitrary \(\xi\)-dependent weights, including ones
whose component ratio \(\|u\|/\|z\|\) tends to zero or infinity.

For the matching upper estimate on the minimum, choose one fixed face
\(p_0\) in the box and its original edge-indicator weight \[
(u_0)_e=\mathbf1_{\{e\in\partial p_0\}}.
\tag{7.11}
\] For every \(z\in\mathscr Z_L\),
\(\langle u_0,z\rangle=(\mathsf Iz)_{p_0}=0\), so
\(u_0\in\mathscr R_L\). Its actual four entries equal 1 and
\(\|u_0\|^2=4\). The absolute \(\mathscr R_L\)-block bound on
\(\mathcal B\), together with (7.5), gives \[
\left|\frac{u_0^T\mathcal N u_0}{u_0^TCu_0}-3\kappa\right|
\le\frac{A_L}{c_L}\kappa\xi^2,
\tag{7.12}
\] enlarging \(A_L\) to bound that entire block in absolute value if
necessary. Both numerator and denominator in (7.12) are the original
state forms; no rescaled state norm has replaced either.

Since \(C(\xi)\) is positive definite, this continuous homogeneous
quotient attains its minimum on the nonzero real weight rays:
equivalently one applies finite-dimensional compactness on a coordinate
unit sphere while keeping the denominator \(w^TCw\) unchanged. Combining
(7.10) and (7.12) proves \[
\boxed{
\min_{w\in\mathbb R^{\mathsf E_L}\setminus\{0\}}
\frac{\mathfrak q_\xi(v_{w,\xi},v_{w,\xi})}{\|v_{w,\xi}\|^2}
=3\kappa+O_L(\kappa\xi^2),\qquad \xi\downarrow0.}
\tag{7.13}
\] This is uniform over all weights on the fixed box, not uniform in
\(L\). It uses the actual positive fourth-order form on the leading
kernel; deleting that form would leave the degenerating denominator
uncontrolled.

\subsection{8. Verification and limits of the
result}\label{magnetic-part-4-verification-and-limits-of-the-result}

The companion check\_electric\_covariance\_incidence\_20260908.py
enumerates all links and all elementary faces for \(L=2,3,4\). It
verifies edge and face counts, the incident-face numbers \(r_e\),
uniqueness of the six-edge support for every unordered adjacent pair,
the full anchored kernel construction and reconstruction, and every
individual parameter basis vector for the entire \(L=2\) box. These
boxes have respectively kernel dimensions 60, 126, 216 and adjacent-pair
counts 1128, 3852, 9168. It verifies both the alternating family and the
native \(n_2^2\) coefficient, along with exact rational arithmetic for
the two-face vacuum coefficients, the norm coefficient \(196/13689\),
the energy coefficient \(8/117\), and quotient \(234/49\).

The rank and kernel theorem itself is proved for every \(L\) in Section
3; it is not inferred from those finite checks. Likewise the checks do
not replace the actual-vacuum analytic and domain arguments in Sections
1, 2, 4, and 5. No floating-point vacuum approximation, numerical Haar
sampling, continuum replacement, remote action, PDF generation, new
agent, or heavy job is used.

The parent finite-range matrix construction is in
magnetic\_translation\_true\_vacuum.md, Section 8; the standalone
full-vertex projection and native angular coefficient are in
MAGNETIC\_TRANSLATION\_PLAQUETTE.md; the adjacent-mode normalization is
consistent with BULK\_PLAQUETTE\_EXTENSION.md. All constants and the
specific proofs needed for the present calculation are given above.
Every conclusion concerns a fixed finite, full interacting box with a
controlled local expansion in \(b/\kappa\). The guaranteed disk
\(3/(16M)\) and all stated remainder constants are retained; no
uniform-volume, fixed-coupling joint-limit, or four-dimensional
continuum gap conclusion is claimed.

For historical convention context, see J. Kogut and L. Susskind,
``Hamiltonian formulation of Wilson's lattice gauge theories,''
\emph{Physical Review D} \textbf{11}, 395--408 (1975),
\href{https://doi.org/10.1103/PhysRevD.11.395}{doi:10.1103/PhysRevD.11.395}.
This is a convention reference, not a substitution of its coupling
normalization for (1.3)--(1.4). The full incidence kernel, its mode
coefficients, and the all-weight estimates used here are proved in this
source.

\clearpage
% Complete next live proof; its original equation tags are retained.
\section{The actual first physical excitation band of the open SU(2)
box}\label{magnetic-part-5-the-actual-first-physical-excitation-band-of-the-open-su2-box}

8 September 2026. The calculation below concerns the eigenvalues of the
full interacting Hamiltonian, not only the energy of a chosen trial
state. It retains all spatial links and plaquettes, the two shared-link
spin channels, the actual vacuum energy, and the original physical
coefficients. Its finite-box error bound is explicit. The large-box
limit computed here is a limit of a Taylor coefficient, not an
interchange of the volume and coupling limits.

\subsection{1. Hamiltonian and the
result}\label{magnetic-part-5-hamiltonian-and-the-result}

Fix an integer \(L\geq2\). The vertices are
\(\mathsf V_L=\{-L,\ldots,L\}^3\). Let \(\mathsf E_L\) contain each
positively oriented nearest-neighbor link \(e=(n,i)\) whose endpoints
belong to this set, and let \(\mathsf P_L\) contain every elementary
square \(p=(n;i,j)\), \(i<j\), including the boundary squares. Write \[
 m=2L,\qquad N=3m(m+1)^2,\qquad M=3m^2(m+1).
\tag{1.1}
\] For \(U_e\in SU(2)\), inverse traversal means \(U_e^{-1}\), and \[
 W_p(U)=\operatorname{tr}\!left(
 U_i(n)U_j(n+\mathbf e_i)U_i(n+\mathbf e_j)^{-1}U_j(n)^{-1}\right).
\tag{1.2}
\] The trace is real. Use product Haar probability measure on
\(SU(2)^N\), \(T_a=-i\sigma_a/2\), and \[
 X_{e,a}f=\left.\frac{d}{dt}f(\ldots,e^{tT_a}U_e,\ldots)\right|_{t=0},
 \qquad E_e=-\sum_{a=1}^3X_{e,a}^2.
\tag{1.3}
\] The spin-\(j\) eigenvalue of \(E_e\) is \(j(j+1)\). In particular the
fundamental value is \(3/4\). The original Lie metric
\(c(A,B)=-\operatorname{tr}(AB)/2\) has orthonormal basis \(2T_a\), so
\(-\Delta_e^c=4E_e\); this factor is retained in \[
 \begin{split}
 H&=\kappa H_0+b\sum_p(2-W_p)
    =\kappa K_\xi+2bM I,\qquad H_0=\sum_eE_e,\\
 K_\xi&=H_0-\xi S,\qquad S=\sum_pW_p,\\
 \kappa&=\frac{2g_{\rm YM}^2}{a},\qquad
 b=\frac{1}{2g_{\rm YM}^2a},\qquad
 \xi=\frac b\kappa=\frac{1}{4g_{\rm YM}^4}.
 \end{split}
\tag{1.4}
\] Here \(a>0\) is spatial lattice spacing. There are three spatial
dimensions and Hamiltonian time. No spatial dimension, link, or
interaction is removed.

The physical Hilbert space \(\mathcal H_{\rm phys}\) is the closed
subspace of \(L^2(SU(2)^N,dU)\) invariant under
\(U_e\mapsto g_{s(e)}U_eg_{t(e)}^{-1}\) at every vertex. Both \(H_0\)
and \(S\) preserve it. The full operator has Sobolev domain \(H^2\) and
form domain \(H^1\); on the physical subspace use their intersections
with \(\mathcal H_{\rm phys}\). Bounded smooth multiplication by \(S\)
preserves self-adjointness and compact resolvent. The heat kernel of the
link Laplacian is strictly positive on the connected compact
configuration space, and the bounded real potential preserves positivity
improvement by its heat-kernel product formula. Consequently the lowest
eigenfunction can be chosen strictly positive and is unique up to a
scalar. Alternatively, its modulus minimizes the form, elliptic
regularity and the maximum principle make it strictly positive, and
integration by parts gives \[
 \langle\psi f,(H-\mathcal E)\psi f\rangle
   =\kappa\sum_{e,a}\int|X_{e,a}f|^2\psi^2dU.
\tag{1.5}
\] Any other lowest eigenfunction divided by \(\psi\) is therefore
constant. Gauge transformations fix this positive unit vector. Thus the
full and physical bottoms agree. Write
\(\mathcal E(\xi)=2bM+\kappa e(\xi)\), and let \(\operatorname{gap}_L\)
be the least strictly positive physical eigenvalue of
\(H-\mathcal E(\xi)I\).

Define the finite face graph as follows. Two distinct faces are adjacent
exactly when they share a physical edge. Its adjacency matrix is
\(A_L\), its face degrees are \(d_p\), and \[
 D_L^{\rm face}=\operatorname{diag}(d_p),\qquad
 Q_L=D_L^{\rm face}+A_L,\qquad q_L=\lambda_{\max}(Q_L),\qquad
 \nu_L=\frac7{15}-\frac{q_L}{21}.
\tag{1.6}
\] The superscript distinguishes this degree matrix from any original
cusp determinant. No cusp coordinate or determinant has been identified
with it.

The result is \[
 \boxed{\operatorname{gap}_L
   =\kappa\left(3+\nu_L\xi^2+\mathcal R_L(\xi)\right),\qquad
 |\mathcal R_L(\xi)|\leq
 \frac73\left(\frac{64M}{3}\right)^4|\xi|^4,
 \quad |\xi|\leq\frac{3}{128M}.}
\tag{1.7}
\] The error bound concerns the actual physical eigenvalue. Moreover \[
 \boxed{0<\nu_L+\frac{71}{105}
       \leq\frac{4(6L+1)}{21L(2L+1)},\qquad
       \lim_{L\to\infty}\nu_L=-\frac{71}{105}.}
\tag{1.8}
\] In particular \(\nu_L<0\) for every \(L\geq2\). The original
coefficients in (1.7) are equivalently \[
 \operatorname{gap}_L=
 \frac{6g_{\rm YM}^2}{a}
 +\frac{\nu_L}{8g_{\rm YM}^6a}
 +\frac{2g_{\rm YM}^2}{a}\mathcal R_L\!\left(\frac1{4g_{\rm YM}^4}\right).
\tag{1.9}
\] Thus no change of energy units or coupling is hidden in the negative
second coefficient. Sections 2--6 prove all parts of these statements.

\subsection{2. The free physical eigenspace and its
isolation}\label{magnetic-part-5-the-free-physical-eigenspace-and-its-isolation}

Peter--Weyl decomposition on each link gives an orthogonal decomposition
by spin assignments \((j_e)\). Gauge invariance is imposed by taking the
invariant tensor space at each vertex, with dual representations on
oppositely oriented indices. This is an orthogonal projection because it
is the integral of the unitary vertex action. Tensor products of the
resulting invariant contractions span the physical subspace, by applying
that bounded projection to the dense product Peter--Weyl span. Their
free energies are exactly \(\sum_ej_e(j_e+1)\).

The support of a nonconstant such contraction is the graph of links with
\(j_e>0\). It cannot have a vertex of degree one: its tensor factor
there would be a single nontrivial irreducible representation, which has
no invariant vector. Every nonempty finite graph of minimum degree at
least two contains a cycle, by continuing an edge path until a vertex
repeats. The cubic lattice graph is simple and bipartite, so every cycle
has length at least four. Each support edge contributes at least
\(3/4\). Hence every nonzero physical energy is at least \(3\).

If the energy is strictly below \(9/2\), there are at most five support
edges. A cycle in this bipartite graph must then have length four. A
fifth edge cannot form a separate component of minimum degree two,
cannot have an endpoint outside the four-cycle without producing degree
one, and cannot be a chord: the possible new chords join vertices in the
same bipartition class. Therefore the support is just that four-cycle.
Every four-cycle in the cubic nearest-neighbor graph is one elementary
square: closing a four-step simple path requires one positive and one
negative step in each of two distinct coordinate directions.

At a degree-two vertex, Schur orthogonality implies that an invariant
contraction exists only when the two incident spins agree; it is then
one-dimensional. Thus a square has the same spin \(j\) on all four
edges, with energy \(4j(j+1)\). Below \(9/2\), only \(j=1/2\) occurs,
giving energy \(3\) and the fundamental Wilson trace \(W_p\). The next
possible physical energy is exactly \(9/2\), realized by the fundamental
trace on a six-edge \(1\)-by-\(2\) rectangle in the box.

Single-link Haar integration gives \[
 \int W_p\,dU=0,\qquad \langle W_p,W_q\rangle=\delta_{pq},\qquad
 H_0W_p=3W_p.
\tag{2.1}
\] For distinct squares a link in only one boundary makes the integral
zero under multiplication of that link by \(-I\). For equal squares
their holonomy is Haar distributed and the fundamental character has
squared norm one. This last identity also follows by writing a Haar
element as \(u_0I+i\sum_au_a\sigma_a\) on the unit three-sphere and
using \(\int u_0^2=1/4\). Consequently \[
 \ker H_0=\mathbb C1,\qquad
 \ker(H_0-3)=\mathcal P:=\operatorname{span}\{W_p:p\in\mathsf P_L\},
 \qquad\operatorname{spec}(H_0)\cap(3,9/2)=\varnothing.
\tag{2.2}
\] In this section spectra and kernels are physical ones. Let \(P\) be
the orthogonal projection onto \(\mathcal P\), let \(Q=I-P\) on the
physical space, and set \[
 R_3=Q\bigl(Q(H_0-3)Q\bigr)^{-1}Q.
\tag{2.3}
\] The inverse is bounded, since zero is excluded by (2.2). It acts by
\(-1/3\) on constants. Omitting this negative denominator would give an
incorrect excitation matrix.

\subsection{3. Every virtual channel, with its exact Haar
map}\label{magnetic-part-5-every-virtual-channel-with-its-exact-haar-map}

The elementary central link map \[
 U_{(n,i)}\longmapsto z_i(n)U_{(n,i)},\qquad
 z_i(n)=(-1)^{\sum_{j<i}n_j}I
\tag{3.1}
\] induces a Haar unitary involution \(\mathcal Z\). Around an \(ij\)
face with \(i<j\), the exponent in \(z_j(n+\mathbf e_i)\) differs by one
from that in \(z_j(n)\), while the other pair agrees. Thus every \(W_p\)
changes sign. The map commutes with every link Casimir and every vertex
gauge action, fixes the constant function, and is \(-I\) on
\(\mathcal P\). It follows in particular that \(PSP=0\).

Resolve the product of two face traces as follows. For a repeated face
put \(Y_p=W_p^2-1\). The character identity
\(\chi_{1/2}^2=\chi_0+\chi_1\) gives \[
 H_0Y_p=8Y_p,\qquad\|Y_p\|^2=1.
\tag{3.2}
\] For two distinct edge-disjoint faces, \(W_pW_q\) has energy \(6\) and
squared norm one, even if their boundaries meet at a vertex. Each of the
eight links has fundamental spin; integration of one independent link of
each face establishes the norm.

If distinct \(p,q\) share one edge \(e\), define the exact orthogonal
Haar projection and its complement \[
 Z_{pq,0}=\int_{SU(2)}W_pW_q\,dU_e,\qquad
 Z_{pq,1}=W_pW_q-Z_{pq,0}.
\tag{3.3}
\] Reverse one face traversal if needed, which preserves its real SU(2)
trace, and cyclically write the traces as \(\operatorname{tr}(U_eA)\)
and \(\operatorname{tr}(U_e^{-1}B)\). The six outer link variables in
A,B are distinct. The identity
\(\int U_{ij}\overline U_{kl}\,dU=\delta_{ik}\delta_{jl}/2\) gives \[
 Z_{pq,0}=\tfrac12\operatorname{tr}(AB).
\tag{3.4}
\] This is the fundamental trace on their six-edge boundary with its
actual factor \(1/2\), not a trace rescaled to unit norm. The
shared-edge tensor product decomposes into spin zero and spin one. Hence
\[
 \begin{array}{c|cc}
 &Z_{pq,0}&Z_{pq,1}\\ \hline
 H_0\text{ eigenvalue}&9/2&13/2\\
 \text{squared norm}&1/4&3/4\\
 R_3\text{ multiplier}&2/3&2/7
 \end{array}
\tag{3.5}
\] The first norm follows from (3.4). The total product norm is one by
integrating the outer independent Haar words A and B first, and the two
channels are orthogonal because (3.3) is an orthogonal projection.

There are no missing cross terms between different unordered pairs.
Indeed a product associated with \(\{p,q\}\) has odd central-link parity
precisely on \(\partial p\mathbin\triangle\partial q\), in either
channel. If two different unordered pairs had the same boundary parity,
the symmetric difference of their face sets would be a nonempty closed
mod-two collection of at most four squares. Any square in it would need
another square on each of its four edges. Distinct elementary squares
share at most one edge, so four distinct other squares would be
necessary, contradicting the maximum of four squares in the collection.
Thus some link has unequal parity and its Haar integration kills every
mixed inner product. The repeated-face vectors have even parity
everywhere, so are orthogonal to these pairs; different repeated-face
vectors are orthogonal by their different joint link-spin assignments.
All these orthogonalities remain true after applying \(H_0\) or \(R_3\).

For \(v=\sum_pa_pW_p\), the entire first-order dressing of the spectral
subspace is consequently \[
 \begin{split}
 R_3Sv={}&-\frac13\sum_pa_p
 +\frac15\sum_pa_pY_p\\
 &+\frac13\sum_{\{p,q\}\ {m edge\!\!-\!\!disjoint}}
                      (a_p+a_q)W_pW_q\\
 &+\sum_{\{p,q\}\ {\rm adjacent}}(a_p+a_q)
                  \left(\frac23Z_{pq,0}+\frac27Z_{pq,1}\right).
 \end{split}
\tag{3.6}
\] The sums over distinct pairs are unordered. This is an exact map from
each free physical face vector into all of its first-order virtual
states.

\subsection{4. Degenerate perturbation and the vacuum
subtraction}\label{magnetic-part-5-degenerate-perturbation-and-the-vacuum-subtraction}

The second effective coefficient before subtracting the bottom energy is
\[
 K^{(2)}=-PSR_3SP.
\tag{4.1}
\] For completeness, this follows directly by writing the analytically
transported vector as \(v+\xi v_1+\xi^2v_2+\cdots\) with \(Pv_1=0\). The
first coefficient of the eigen-equation or spectral intertwining
equation gives \((H_0-3)v_1=QSv\), hence \(v_1=R_3Sv\). Projecting the
next coefficient onto \(\mathcal P\) gives (4.1). This works for the
whole degenerate subspace; no individual face is assumed to remain an
eigenvector. The exact analytic subspace map used here is constructed in
Section 5.

The vacuum intermediate state contributes \(+1/3\) to every matrix entry
in the \(W_p\) basis. A distinct edge-disjoint pair contributes \(-1/3\)
to its two diagonal entries and its two mutual off-diagonal entries. An
adjacent pair contributes the same pattern with value \[
 -\left(\frac{1/4}{9/2-3}+\frac{3/4}{13/2-3}\right)
       =-\left(\frac16+\frac3{14}\right)=-\frac8{21}.
\tag{4.2}
\] A repeated-face channel contributes \(-1/5\) to its diagonal.
Therefore \[
 \begin{split}
 K^{(2)}_{pp}
  &=\frac13-\frac15-\frac{M-1-d_p}{3}-\frac8{21}d_p
    =\frac7{15}-\frac M3-\frac{d_p}{21},\\
 K^{(2)}_{pq}
  &=\begin{cases}-1/21,&p,q\text{ adjacent},\\0,&p\ne q\text{ not adjacent}.
    \end{cases}
 \end{split}
\tag{4.3}
\] In particular the nonlocal rank-one vacuum contribution cancels the
edge-disjoint off-diagonal contribution, rather than either one being
omitted.

For the actual positive unit vacuum, coefficient comparison yields \[
 \psi_\xi=1+\frac\xi3S+O_{H^2,L}(\xi^2),\qquad
 e(\xi)=-\frac M3\xi^2+O_L(\xi^4).
\tag{4.4}
\] The first equation follows because \(H_0S=3S\). Pairing the
second-order equation with constants gives
\(-\langle S,S\rangle/3=-M/3\). The error is even by (3.1), as proved
below. Thus the physical excitation coefficient is \[
 \boxed{F_L=K^{(2)}+\frac M3I
          =\frac7{15}I-\frac1{21}(D_L^{\rm face}+A_L).}
\tag{4.5}
\] The first term in (3.6) also explicitly checks the
vacuum-orthogonality sign. Pairing \(v+\xi R_3Sv\) with (4.4) gives at
order \(\xi\) the two cancelling scalars \(-\sum_pa_p/3\) and
\(+\sum_pa_p/3\). The exact transported subspace is orthogonal to the
exact vacuum at real coupling, since their spectral contours are
disjoint.

\subsection{5. Exact spectral transport, evenness, and an explicit
remainder}\label{magnetic-part-5-exact-spectral-transport-evenness-and-an-explicit-remainder}

This section supplies the operator error bound in (1.7). Work throughout
in \(\mathcal H_{\rm phys}\). Put \[
 r=\frac{3}{64M},\qquad \mathcal C_3=\{z:|z-3|=3/4\}.
\tag{5.1}
\] On \(\mathcal C_3\), equation (2.2) gives
\(\|(z-H_0)^{-1}\|\leq4/3\). Since \(\|S\|=2M\), the Neumann factor is
at most \(1/8\) for complex \(|\xi|\leq r\). Equality in \(\|S\|=2M\)
follows from the all-identity link configuration and its
positive-measure neighborhoods. Consequently the analytic resolvent
\(R_\xi(z)=(z-K_\xi)^{-1}\) satisfies \[
 \|R_\xi(z)\|\leq\frac{32}{21},\qquad
 P_\xi=\frac1{2\pi i}\int_{\mathcal C_3}R_\xi(z)\,dz,
 \qquad \|P_\xi-P\|\leq\frac17.
\tag{5.2}
\] For the last bound, integrate the resolvent difference, using
\((3/4)(4/3)^2(2Mr)/(1-1/8)=1/7\). The contour integral is an idempotent
of rank M, by its continuous deformation from P.

On \(\mathcal P\) set \(T_\xi=PP_\xi P|_{\mathcal P}\). The convergent
binomial series at I defines \(T_\xi^{-1/2}\) analytically, since
\(\|T_\xi-I\|\leq1/7\), with \(\|T_\xi^{-1/2}\|\leq\sqrt{7/6}\). Define
exact maps \[
 V_\xi=P_\xi P T_\xi^{-1/2}:\mathcal P\longrightarrow
                    \operatorname{ran}P_\xi,\qquad
 U_\xi=T_\xi^{-1/2}PP_\xi:\operatorname{ran}P_\xi\longrightarrow
                    \mathcal P.
\tag{5.3}
\] Idempotence gives \(U_\xi V_\xi=I\); equal finite dimensions give the
reverse identity on \(\operatorname{ran}P_\xi\). For real \(\xi\),
\(P_\xi\) is orthogonal, \(U_\xi=V_\xi^*\), and \(V_\xi\) is an
isometry. Formula (5.3) retains the whole Gram factor; it is not a
removal of a physical field amplitude or scale. Its derivative at zero
has \(PV'_0=0\), as differentiating \(P_\xi^2=P_\xi\) gives
\(PP'_0P=0\). Differentiating its spectral intertwining equation then
gives exactly \(V'_0=R_3SP\), as used in (3.6) and (4.1).

The exact finite matrix representing the cluster is \[
 B_\xi=U_\xi K_\xi V_\xi
      =T_\xi^{-1/2}P(K_\xi P_\xi)P T_\xi^{-1/2}.
\tag{5.4}
\] The products are bounded because \(K_\xi P_\xi\) is its contour
integral \((2\pi i)^{-1}\int zR_\xi(z)dz\). In particular \[
 \begin{split}
 \|B_\xi-3I\|
 &\leq\frac76\left\|\frac1{2\pi i}
       \int_{\mathcal C_3}(z-3)R_\xi(z)\,dz\right\|\\
 &\leq\frac76\left(\frac34\right)^2\frac{32}{21}=1.
 \end{split}
\tag{5.5}
\] The corresponding contour \(|z|=3/4\) around the physical free bottom
has the same resolvent bounds and rank one. Its unique eigenvalue e(ξ)
is analytic and lies within that circle, so \(|e(\xi)|\leq3/4\).

Equation (3.1) gives \(K_{-\xi}=\mathcal ZK_\xi\mathcal Z\),
\(P_{-\xi}=\mathcal ZP_\xi\mathcal Z\), and \(T_{-\xi}=T_\xi\), since
\(\mathcal ZP=-P\). Thus \(V_{-\xi}=-\mathcal ZV_\xi\),
\(U_{-\xi}=-U_\xi\mathcal Z\), and \[
 B_{-\xi}=B_\xi,\qquad e(-\xi)=e(\xi).
\tag{5.6}
\] These are analytic identities, not a choice of labels at a degenerate
eigenvalue. Therefore the matrix \(G_\xi=B_\xi-e(\xi)I-3I\) is even
analytic on \(|\xi|\leq r\), has \(G_0=0\), has second coefficient F\_L
by (4.5), and has norm at most \(7/4\). Cauchy's coefficient estimate
followed by the geometric series gives, for \(|\xi|\leq r/2\), \[
 \|G_\xi-\xi^2F_L\|
 \leq\frac74\sum_{k\geq2}\left(\frac{|\xi|}{r}\right)^{2k}
 \leq\frac73r^{-4}|\xi|^4.
\tag{5.7}
\] There is a neighborhood beyond the closed radius r where the contours
and series above exist, since each displayed Neumann and Gram bound is
strict at its analytic threshold, so Cauchy's estimate at r is
legitimate.

For real \(|\xi|\leq r/2\), the perturbation norm is at most
\(2M|\xi|\leq3/64\). The min--max principle moves each ordered physical
eigenvalue of H\_0 by no more than this number. There is one bottom, M
eigenvalues near 3, and all remaining eigenvalues are at least
\(9/2-3/64\); hence the first positive physical excitation is precisely
the lowest eigenvalue of \(\kappa(B_\xi-e(\xi)I)\). For Hermitian
matrices the min--max principle bounds the change of every ordered
eigenvalue by the perturbation norm. Applying it to (5.7), and observing
that \(\lambda_{\min}(F_L)=7/15-q_L/21\), proves (1.7), including its
absolute error bound.

The constant \(2bM\) has not disappeared by fiat: both the ground energy
and every physical cluster eigenvalue of H are their K\_ξ counterparts
multiplied by κ and then increased by \(2bM\). Their difference is
(1.7).

\subsection{6. The unchanged open-box graph and its coefficient
limit}\label{magnetic-part-5-the-unchanged-open-box-graph-and-its-coefficient-limit}

For a link in direction i, let j,k be the other directions and put \[
 a_j(n_j)=\begin{cases}1,&n_j=\pm L,\\2,&-L<n_j<L.\end{cases}
 \qquad r_e=a_j(n_j)+a_k(n_k).
\tag{6.1}
\] This is exactly the number of faces incident to e. The face degree is
\(d_p=\sum_{e\in\partial p}(r_e-1)\), since two different squares share
at most one edge. Hence \(d_p\leq12\), with strict inequality for some
boundary faces. The number \(\mathcal A_L\) of unordered adjacent face
pairs is \[
 \mathcal A_L=\sum_e\binom{r_e}{2}=144L^3-12L.
\tag{6.2}
\] Here are the finite sums establishing the coefficient. For either
transverse coordinate, \(\sum a_j=2m\) and \(\sum a_j^2=4m-2\). For one
direction i, summing \(r_e(r_e-1)\) over the two transverse coordinates
gives \[
 2(m+1)(4m-2)+2(2m)^2-2(m+1)(2m)=12m^2-4.
\] There are m positions along i, three choices of i, and the factor
\(1/2\) in (6.2), giving \(18m^3-6m=144L^3-12L\). In particular
\(\sum_pd_p=2\mathcal A_L\).

The constant face vector gives the Rayleigh lower bound \[
 q_L\geq\frac{4\mathcal A_L}{M}
       =\frac{48L^2-4}{L(2L+1)},\qquad q_L\leq24.
\tag{6.3}
\] The upper bound follows from the maximum absolute row sum of Q\_L,
equal to \(2\max_pd_p\leq24\). It is strict. The face graph is
connected: within a coordinate plane the square tiling is connected by
edge adjacency; neighboring parallel layers connect through a square in
a transverse plane, and that same adjacency joins the three plane
orientations. A maximizing eigenvector of the real symmetric nonnegative
Q\_L can be chosen nonnegative, because replacing its entries by their
absolute values does not decrease its quadratic form. Connectedness and
the eigen-equation make all entries positive. If its eigenvalue were 24,
an entry of maximum size would force its degree to equal 12 and every
neighboring entry to have the same size. Propagation through the
connected graph would force every face degree to equal 12, contradicting
the boundary faces. Thus \(q_L<24\).

Subtracting (6.3) from 24 gives \[
 0<24-q_L\leq\frac{24L+4}{L(2L+1)}.
\tag{6.4}
\] Equations (1.6) and (6.4) prove (1.8). Also
\(24-(24L+4)/(L(2L+1))\geq24-14/L\geq17\) for \(L\geq2\). Since
\(17>49/5\), (6.3) gives \(\nu_L<7/15-(49/5)/21=0\). This proof used the
actual open incidence sets, not periodic boundaries.

\subsection{7. Meaning of the coefficient and relation to previous
states}\label{magnetic-part-5-meaning-of-the-coefficient-and-relation-to-previous-states}

At each fixed box, the interacting first physical eigenvalue decreases
from \(3\kappa\) at second order in ξ. This statement is stronger than a
change in one proposed Rayleigh quotient: Sections 2 and 5 isolate the
entire first physical cluster, and Section 4 diagonalizes its full
second-order operator. Its neighbor coefficient \(-1/21\) and the degree
contribution are exact effects of both retained spin channels.

The earlier electric and magnetic-translation trial families have
leading energy \(3\kappa\) in their non-kernel directions. The exact map
from that leading face space to the actual first spectral subspace is
V\_ξ in (5.3), with the first correction (3.6). It includes the vacuum
subtraction, repeated-face spin-one terms, all edge-disjoint pairs, and
all hinged or coplanar adjacent pairs. Thus this calculation supplies
their full first spectral dressing rather than claiming those trial
families and the physical spectrum are unrelated.

No zero of a truncated polynomial is asserted to be a zero of the
physical gap. The guaranteed interval in (1.7) shrinks as \(M^{-1}\),
while its explicit fourth-order bound grows as \(M^4\). Equation (1.8)
is a coefficient limit and does not prove a fixed-positive-coupling
infinite-volume limit of (1.7). The state-space and spectral transport
are exact at each finite box; extending the spectral estimates to the
original interacting spatial continuum remains a mathematical
calculation not supplied by this coefficient alone.

\subsection{8. Primary literature and attribution of this
calculation}\label{magnetic-part-5-primary-literature-and-attribution-of-this-calculation}

The Hamiltonian lattice framework is that of J. Kogut and L. Susskind,
\emph{Hamiltonian formulation of Wilson's lattice gauge theories},
Physical Review D \textbf{11} (1975), 395--408,
\url{https://doi.org/10.1103/PhysRevD.11.395}. Equations (1.3)--(1.4)
give the complete conventions actually used here, rather than borrowing
a coefficient from a differently parametrized Hamiltonian.

Strong-coupling scalar-gap series in three spatial dimensions already
have an extensive literature, including A. C. Irving, T. E. Preece, and
C. J. Hamer, \emph{Cluster expansion approach to non-abelian lattice
gauge theory in (3+1)D (I). SU(2)}, Nuclear Physics B \textbf{270}
(1986), 536--552, \url{https://doi.org/10.1016/0550-3213(86)90567-5}.
The publisher's abstract identifies its vacuum-energy, string-tension
and scalar-gap expansions. The present derivation is self-contained and
does not claim historical novelty for a strong-coupling coefficient; no
unread coefficient table from that article is used as a proof or as a
claimed numerical match.

The interacting four-dimensional continuum target is the construction
and spectral assertion in A. Jaffe and E. Witten, \emph{Quantum
Yang--Mills Theory},
\url{https://www.claymath.org/wp-content/uploads/2022/06/yangmills.pdf}.
Neither that continuum construction nor a counterexample to its mass-gap
assertion is claimed by the finite-box result (1.7).

\clearpage
% Complete next live proof; its original equation tags are retained.
\section{The full open-box face graph, its bulk symbol, and the
first-band coefficient
limit}\label{magnetic-part-6-the-full-open-box-face-graph-its-bulk-symbol-and-the-first-band-coefficient-limit}

8 September 2026. This calculation retains every boundary face of the
original open spatial box. The infinite and periodic graphs are defined
independently, and their relation to the open graph is an explicit
zero-extension identity. The interior sine function is a trial vector in
the actual finite open graph, not a change of boundary conditions.

\subsection{1. Original open box, physical orientations, and the face
operator}\label{magnetic-part-6-original-open-box-physical-orientations-and-the-face-operator}

Let \(L\ge2\) be an integer, put \(m=2L\), and take vertices in
\(\{-L,\ldots,L\}^3\). A physical edge \(e=(n,i)\) runs from \(n\) to
\(n+\mathbf e_i\), whenever both endpoints belong to the box. The face
\(p=(n;i,j)\), \(i<j\), has the canonical oriented edge word

\[
(n,i)^+\,(n+\mathbf e_i,j)^+\,
(n+\mathbf e_j,i)^-\,(n,j)^-.
\tag{1.1}
\]

Its underlying set of four physical edges is denoted \(\partial p\);
signs in (1.1) do not create additional edges. All contained elementary
faces are included in \(\mathsf P_L\). The number of faces is

\[
M_L=|\mathsf P_L|=3m^2(m+1)=12L^2(2L+1).
\tag{1.2}
\]

Indeed each of the three face planes has \(m\) base choices in each
plane direction and \(m+1\) in its normal direction.

Two distinct faces are adjacent exactly when their physical edge sets
intersect. Two distinct elementary squares in the open box share at most
one physical edge: coplanar unit squares can meet along only one side,
while perpendicular coordinate squares intersect in at most a single
unit coordinate segment. Thus the adjacency matrix \(\mathsf A_L\) is a
simple unweighted symmetric matrix: \[
(\mathsf A_L)_{pq}=\mathbf1_{\{p\ne q,\ \partial p\cap\partial q\ne\varnothing\}}.
\] Let \(d_p=\sum_q(\mathsf A_L)_{pq}\), and define the actual open
signless face Laplacian

\[
\boxed{\mathsf Q_L=\operatorname{diag}(d_p)+\mathsf A_L.}
\tag{1.3}
\]

For every complex face vector \(x\), \[
\langle x,\mathsf Q_Lx\rangle
=\sum_{\{p,q\}\text{ adjacent}}|x_p+x_q|^2.
\tag{1.4}
\] This follows by expanding each square: each diagonal \(|x_p|^2\)
occurs \(d_p\) times, and every cross term agrees with the adjacency
quadratic form. In particular \(\mathsf Q_L\) is positive semidefinite.
This is the signless operator, not
\(\operatorname{diag}(d)-\mathsf A_L\).

If \(\mathsf I_{pe}=\mathbf1_{\{e\in\partial p\}}\) is the original
unsigned face-edge incidence matrix, the same one-edge intersection fact
gives \[
\mathsf I\mathsf I^T=4I+\mathsf A_L,\qquad
\mathsf Q_L=\mathsf I\mathsf I^T+\operatorname{diag}(d_p-4).
\tag{1.5}
\] The diagonal 4 in this incidence product is not the adjacency degree,
which varies at the boundary.

\subsection{2. Exact degrees, all boundary multiplicities, and
connectedness}\label{magnetic-part-6-exact-degrees-all-boundary-multiplicities-and-connectedness}

For an edge \(e=(n,i)\), let \(r_e\) be the number of contained
elementary faces incident on it. If \(j,k\) are its two transverse
directions, then

\[
r_e=\delta(n_j)+\delta(n_k),\qquad
\delta(t)=
\begin{cases}
1,&t=-L\text{ or }L,\\
2,&-L<t<L.
\end{cases}
\tag{2.1}
\]

In each transverse direction a face can extend on both sides in the
interior and only inward at the boundary. Thus \(2\le r_e\le4\). Every
neighbor of a fixed face is counted at its unique shared edge, so \[
d_p=\sum_{e\in\partial p}(r_e-1).
\tag{2.2}
\]

For a face \(p=(n;i,j)\), write \(k\) for its normal direction and
define, on the plane base-coordinate range \(-L\le t\le L-1\), \[
\beta(t)=\mathbf1_{\{t=-L\}}+\mathbf1_{\{t=L-1\}}.
\tag{2.3}
\] Since \(\delta(t)+\delta(t+1)=4-\beta(t)\), summing (2.1) over the
four edges in (2.2) gives the exact boundary formula

\[
\boxed{
d_{n;i,j}=
\begin{cases}
12-\beta(n_i)-\beta(n_j),&-L<n_k<L,\\
8-\beta(n_i)-\beta(n_j),&n_k=-L\text{ or }L.
\end{cases}}
\tag{2.4}
\]

This produces all six degree classes, with no excluded faces:

{\def\LTcaptype{none} % do not increment counter
\begin{longtable}[]{@{}rr@{}}
\toprule\noalign{}
Degree & Exact number of faces \\
\midrule\noalign{}
\endhead
\bottomrule\noalign{}
\endlastfoot
\(12\) & \(3(m-1)(m-2)^2\) \\
\(11\) & \(12(m-1)(m-2)\) \\
\(10\) & \(12(m-1)\) \\
\(8\) & \(6(m-2)^2\) \\
\(7\) & \(24(m-2)\) \\
\(6\) & \(24\) \\
\end{longtable}
}

To verify the table, there are \(m-1\) interior normal-coordinate values
or two boundary values. In either plane direction there are \(m-2\)
bases with \(\beta=0\) and two with \(\beta=1\). Their products give
respectively \((m-2)^2\), \(4(m-2)\), and 4 plane-coordinate pairs with
sum of the two \(\beta\)'s equal to 0, 1, or 2. Multiply by the three
planes and the appropriate normal-coordinate count. Summing the table
gives (1.2).

The exact number \(a_L\) of unordered adjacent pairs can be obtained
either by summing this degree table and dividing by 2, or directly by
the incident edge stars. For each edge direction, hold its \(m\)
longitudinal bases fixed. There are \((m-1)^2\) transverse interior
positions with \(r_e=4\), \(4(m-1)\) positions with exactly one boundary
coordinate and \(r_e=3\), and 4 positions with two boundary coordinates
and \(r_e=2\). Each contributes \(\binom{r_e}{2}\), with no pair counted
twice. Therefore

\[
\boxed{a_L=3m[6(m-1)^2+12(m-1)+4]
=18m^3-6m=144L^3-12L.}
\tag{2.5}
\]

In particular \[
\sum_pd_p=2a_L,\qquad
\sum_p(12-d_p)=12m(3m+1).
\tag{2.6}
\]

The whole face graph is connected. First, the \(12\)-faces in any fixed
\(n_3\) layer form a connected rectangular face grid, by coplanar edge
adjacency. Consecutive \(n_3\) layers are joined through a \(13\)-face:
its bottom and top direction-1 edges belong to \(12\)-faces in the two
layers. For instance choose fixed permissible \(n_1,n_2\) and the
\(13\)-face between the two layers; both neighboring \(12\)-faces exist.
Hence all \(12\)-layers lie in one connected component. Every
\(13\)-face has a direction-1 edge which belongs to a \(12\)-face,
choosing the latter's direction-2 extension inward if that edge is on a
boundary. Every \(23\)-face likewise has a direction-2 edge belonging to
a \(12\)-face, with direction-1 extension chosen inward. Thus every face
belongs to that one component.

\subsection{3. Full open-box bounds and strict inequality at finite
volume}\label{magnetic-part-6-full-open-box-bounds-and-strict-inequality-at-finite-volume}

Let \(\Lambda_L=\lambda_{\max}(\mathsf Q_L)\). The unchanged constant
face vector has squared norm \(M_L\) and, by (1.4), quadratic form
\(4a_L\). Hence

\[
\boxed{
\frac{48L^2-4}{L(2L+1)}
=\frac{4a_L}{M_L}\le\Lambda_L.}
\tag{3.1}
\]

The boundary formula gives \(d_p\le12\). A particularly useful exact
deficit identity is

\[
\boxed{
\langle x,(24I-\mathsf Q_L)x\rangle
=\sum_{\{p,q\}\text{ adjacent}}|x_p-x_q|^2
+2\sum_p(12-d_p)|x_p|^2.}
\tag{3.2}
\]

Indeed the first sum is the quadratic form of
\(\operatorname{diag}(d)-\mathsf A_L\), and adding the second term gives
\(24I-\operatorname{diag}(d)-\mathsf A_L\). Both terms are nonnegative.
Equality would force \(x_p=x_q\) on every adjacency and \(x_p=0\) on
every degree-deficient face. Connectedness makes \(x\) constant, and the
nonempty degree-6 boundary class forces that constant to be zero.
Therefore, in this finite-dimensional space,

\[
\boxed{\frac{48L^2-4}{L(2L+1)}\le\Lambda_L<24.}
\tag{3.3}
\]

This already proves \(\Lambda_L\to24\). It also records why the constant
vector is not an eigenvector at the open boundary: its row values are
\(2d_p\), which are not constant.

\subsection{4. Independently defined bulk and periodic graphs, with full
translations}\label{magnetic-part-6-independently-defined-bulk-and-periodic-graphs-with-full-translations}

The infinite graph has vertices \((n;\alpha)\), where
\(n\in\mathbb Z^3\) and the plane type \(\alpha\) is one of
\(12,13,23\). Its physical edges and oriented faces use (1.1) on the
entire infinite cubic lattice. Two faces are adjacent by the same
shared-physical-edge rule. Every edge has four incident faces, every
face has twelve distinct neighbors, and its signless operator on
\(\ell^2(\mathbb Z^3;\mathbb C^3)\) is \[
\mathsf Q_\infty=12I+\mathsf A_\infty.
\tag{4.1}
\]

Here is its entire translation list. Within plane \(ij\), the four
neighbors of \((n;ij)\) are \((n\pm\mathbf e_i;ij)\) and
\((n\pm\mathbf e_j;ij)\). Between different planes, for
\(a,b\in\{0,1\}\),

{\def\LTcaptype{none} % do not increment counter
\begin{longtable}[]{@{}llrl@{}}
\toprule\noalign{}
Source plane & Target plane & Shared direction & Target base \\
\midrule\noalign{}
\endhead
\bottomrule\noalign{}
\endlastfoot
\(12\) & \(13\) & \(1\) & \(n+a\mathbf e_2-b\mathbf e_3\) \\
\(12\) & \(23\) & \(2\) & \(n+a\mathbf e_1-b\mathbf e_3\) \\
\(13\) & \(23\) & \(3\) & \(n+a\mathbf e_1-b\mathbf e_2\) \\
\end{longtable}
}

Each row has four distinct offsets, including the zero offset, and the
reverse neighbors have their negatives.

To prove the list with orientations retained, the two direction-\(i\)
edges of a face \(ij\) lie at \(n\) and \(n+\mathbf e_j\) with signs
\(+,-\); its two direction-\(j\) edges lie at \(n\) and
\(n+\mathbf e_i\) with signs \(-,+\). A common physical edge between two
planes is obtained by equating one of their two bases in the shared
direction. For the first row this equates \(n+a\mathbf e_2\) with the
target base plus \(b\mathbf e_3\), yielding the displayed offset. The
other two rows follow by equating the direction-2 or direction-3 bases.
The traversal signs in the three rows are respectively
\(((-1)^a,(-1)^b)\), \((-(-1)^a,(-1)^b)\), and \((-(-1)^a,-(-1)^b)\).
The graph is unsigned and so includes all four sign combinations; none
is discarded or counted twice. This proves both the exact translations
and the multiplicity 12.

A separate finite periodic graph can be defined by taking base
coordinates in \[
\mathbb Z/N_1\mathbb Z\times\mathbb Z/N_2\mathbb Z\times\mathbb Z/N_3\mathbb Z,
\qquad N_i\ge3,
\tag{4.2}
\] and using the same physical edge labels and face words modulo those
periods. The condition \(N_i\ge3\) prevents the \(+\mathbf e_i\) and
\(-\mathbf e_i\) neighbors from coinciding and prevents two elementary
faces sharing multiple physical edges. The neighbor list above is then
valid modulo the periods, still with twelve distinct neighbors per face.
This definition is independent of the original open box; it is not
imposed on \(\mathsf P_L\).

\subsection{5. Exact three-plane Fourier symbol and the mode attaining
24}\label{magnetic-part-6-exact-three-plane-fourier-symbol-and-the-mode-attaining-24}

For a finitely supported infinite face field, define \[
\widehat f_\alpha(k)=\sum_{n\in\mathbb Z^3}f_\alpha(n)e^{-ik\cdot n},
\qquad k\in[-\pi,\pi]^3,
\tag{5.1}
\] with inverse measure \(dk/(2\pi)^3\). Orthogonality of the
exponentials proves the isometry on finite sums, which extends by
density to a unitary map from \(\ell^2(\mathbb Z^3;\mathbb C^3)\) onto
the corresponding \(L^2\) space. The translation \(f_\beta(n+\delta)\)
gives the multiplier \(e^{ik\cdot\delta}\). In the order \(12,13,23\),
the neighbor list therefore gives the exact Hermitian symbol

\[
\boxed{
\widehat{\mathsf Q}(k)=12I+
\begin{pmatrix}
2\cos k_1+2\cos k_2&
(1+e^{ik_2})(1+e^{-ik_3})&
(1+e^{ik_1})(1+e^{-ik_3})\\
(1+e^{-ik_2})(1+e^{ik_3})&
2\cos k_1+2\cos k_3&
(1+e^{ik_1})(1+e^{-ik_2})\\
(1+e^{-ik_1})(1+e^{ik_3})&
(1+e^{-ik_1})(1+e^{ik_2})&
2\cos k_2+2\cos k_3
\end{pmatrix}.}
\tag{5.2}
\]

No common translation phases have been suppressed. If one instead uses
each face center as its Fourier origin, define the explicit unitary
fiber change \[
\mathsf D(k)=\operatorname{diag}
\big(e^{i(k_1+k_2)/2},e^{i(k_1+k_3)/2},e^{i(k_2+k_3)/2}\big),
\qquad \widehat f=\mathsf D(k)\,u.
\tag{5.3}
\] Writing \(c_i=\cos(k_i/2)\), one obtains the real centered symbol

\[
\boxed{
\mathsf D(k)^*\widehat{\mathsf Q}(k)\mathsf D(k)
=
\begin{pmatrix}
12+2\cos k_1+2\cos k_2&4c_2c_3&4c_1c_3\\
4c_2c_3&12+2\cos k_1+2\cos k_3&4c_1c_2\\
4c_1c_3&4c_1c_2&12+2\cos k_2+2\cos k_3
\end{pmatrix}.}
\tag{5.4}
\]

For example the phase of the \(12,13\) entry in (5.2) is
\(e^{i(k_2-k_3)/2}\), which cancels exactly with the two center phases
in (5.3). Thus (5.4) is a unitary change of fiber coordinates, not a
replacement of physical translation phases.

At zero momentum, \[
\widehat{\mathsf Q}(0)=12I+4\mathbf1\mathbf1^T,
\tag{5.5}
\] with eigenvalues 24 on the equal-plane vector \((1,1,1)\), and 12
twice on its orthogonal complement.

This is the only Bloch mode at eigenvalue 24, modulo reciprocal-lattice
momenta. To prove it, for an original-base amplitude
\(\eta\in\mathbb C^3\), the exact quadratic form of
\(24I-\widehat{\mathsf Q}(k)=12I-\widehat{\mathsf A}(k)\) is \[
\sum_{\alpha=ij}\sum_{r\in\{i,j\}}
|\eta_\alpha-e^{ik_r}\eta_\alpha|^2
+\sum_{\alpha<\beta}\sum_{\delta\in\mathsf T_{\alpha\beta}}
|\eta_\alpha-e^{ik\cdot\delta}\eta_\beta|^2,
\tag{5.6}
\] where \(\mathsf T_{\alpha\beta}\) is the four-offset set in Section
4. Expanding this expression recovers (5.2): each component receives
four units of same-plane diagonal contribution and eight from the
cross-plane rows. It is nonnegative. Each cross-plane offset set
includes zero, so equality forces all three amplitudes to agree. If they
are nonzero, the same-plane terms force
\(e^{ik_1}=e^{ik_2}=e^{ik_3}=1\). Conversely that momentum and equal
amplitudes give (5.5). Hence \[
\lambda_{\max}(\widehat{\mathsf Q}(k))\le24,
\] with equality exactly at \(k\in(2\pi\mathbb Z)^3\) and the stated
equal-plane eigenvector.

On the independently defined periodic graph (4.2), the same proof uses
the discrete momenta \(k_i=2\pi r_i/N_i\). Its maximum eigenvalue is
exactly 24, achieved only by the constant equal-plane field. On the
infinite graph, the constant Bloch field is not in \(\ell^2\). The
symbol implies operator norm 24: its continuous eigenvalue near zero
gives fields supported in arbitrarily small momentum neighborhoods with
Rayleigh values arbitrarily close to 24. There is no nonzero \(\ell^2\)
eigenvector at 24, because equality in (5.6) would require its Fourier
transform to be supported at the measure-zero momentum \(k=0\).

\subsection{6. Exact map from the infinite graph to the unchanged open
graph}\label{magnetic-part-6-exact-map-from-the-infinite-graph-to-the-unchanged-open-graph}

Let \(J_L:\mathbb C^{\mathsf P_L}\to\ell^2(\mathbb Z^3;\mathbb C^3)\)
extend a face vector by zero on every face outside the open box. This is
an isometry in the original counting inner products. The open face graph
is the induced subgraph of the infinite one: two contained faces share a
physical edge in the finite box exactly when they do so in the infinite
lattice. Consequently

\[
\boxed{
J_L^*\mathsf Q_\infty J_L=12I+\mathsf A_L,\qquad
\mathsf Q_L=J_L^*\mathsf Q_\infty J_L
-\operatorname{diag}(12-d_p).}
\tag{6.1}
\]

This is the exact comparison, including the boundary correction.
Replacing the open operator by its bulk compression would add the
nonzero boundary diagonal \(\operatorname{diag}(12-d_p)\); replacing it
by a periodic graph would also add different adjacency edges. Neither
operation is performed here.

The following stronger bound is a trial calculation in \(\mathsf Q_L\)
itself. Put \[
\ell=2L-2,\qquad t=\frac{\pi}{2L-1}=\frac{\pi}{\ell+1},\qquad c=\cos t,
\] and define on the integer line \[
f_r=
\begin{cases}
\sin(rt),&1\le r\le\ell,\\
0,&\text{otherwise}.
\end{cases}
\tag{6.2}
\] For a face of any of the three plane types with base \(n\), set \[
x_{n;\alpha}=g(n),\qquad
g(n)=\prod_{i=1}^3f_{n_i+L}.
\tag{6.3}
\] This is a vector on every face of the original open box, with zero
entries whenever any base coordinate is outside \([-L+1,L-2]\). Every
face in its support has degree exactly 12 by (2.4). Thus the boundary
diagonal in (6.1) vanishes on the support, and \[
\langle x,\mathsf Q_Lx\rangle
=\langle J_Lx,\mathsf Q_\infty J_Lx\rangle.
\tag{6.4}
\] No outer face or boundary interaction is deleted from the finite
operator: its trial amplitude is simply zero. Equation (6.4) proves the
required map of this particular interior trial to the actual open
signless graph.

To calculate its quotient exactly, let \[
S=\sum_{r=1}^{\ell}f_r^2>0,\qquad
T=\sum_{r=1}^{\ell-1}f_rf_{r+1}.
\] The discrete sine equation \(f_{r-1}+f_{r+1}=2cf_r\), with
\(f_0=f_{\ell+1}=0\), multiplied by \(f_r\) and summed, gives \(T=cS\).
Direct product factorization gives \[
\sum_ng(n)^2=S^3,\quad
\sum_ng(n)g(n+\mathbf e_i)=TS^2,\quad
\sum_ng(n)g(n+\mathbf e_i-\mathbf e_j)=T^2S\quad(i\ne j).
\tag{6.5}
\] All sums are over the zero-extended integer lattice and hence are
finite.

The squared norm of (6.3) is \(3S^3\). The diagonal contribution of
\(\mathsf Q_\infty\) is \(36S^3\). The same-plane adjacency
contribution, from two directions in each of three planes and both signs
of each shift, is \(12TS^2\). For each of three unordered plane pairs,
the four shifts in Section 4 contribute \(S^3+2TS^2+T^2S=S(S+T)^2\); the
reverse matrix entries double this, giving total \(6S(S+T)^2\). Thus
(6.4) yields the exact open-graph trial quotient

\[
\boxed{
\frac{\langle x,\mathsf Q_Lx\rangle}{\|x\|^2}
=\frac{36S^3+12TS^2+6S(S+T)^2}{3S^3}
=14+8c+2c^2.}
\tag{6.6}
\]

In particular the full open-box estimates can be sharpened to

\[
\boxed{
\max\left\{\frac{48L^2-4}{L(2L+1)},\
14+8\cos\frac{\pi}{2L-1}+2\cos^2\frac{\pi}{2L-1}\right\}
\le\Lambda_L<24.}
\tag{6.7}
\]

Since \(24-(14+8c+2c^2)=2(1-c)(5+c)\) and
\(1-\cos t=2\sin^2(t/2)\le t^2/2\), the deficit satisfies

\[
\boxed{
0<24-\Lambda_L
\le\min\left\{
\frac{24L+4}{L(2L+1)},\
2(1-c)(5+c)\right\}
\le\frac{6\pi^2}{(2L-1)^2}.}
\tag{6.8}
\]

The last bound is quantitative \(O(L^{-2})\), obtained without imposing
periodic conditions or dropping the actual boundary diagonal.

\subsection{7. Exact consequence for the specified second-order
coefficient
matrix}\label{magnetic-part-6-exact-consequence-for-the-specified-second-order-coefficient-matrix}

Define the finite matrix requested by the excitation calculation, with
its given constants, by

\[
\mathsf F_L=\frac7{15}I-\frac1{21}\mathsf Q_L.
\tag{7.1}
\]

This note proves the spectrum of this explicitly specified graph matrix.
The derivation identifying (7.1) with the physical true-vacuum
second-order excitation matrix is the separate parent band calculation;
no replacement of that physical calculation is assumed here.

Because \(\mathsf F_L\) is an affine polynomial in the real symmetric
\(\mathsf Q_L\), they have the same orthonormal eigenvectors, with the
exact eigenvalue map \(q\mapsto7/15-q/21\). Therefore

\[
\lambda_{\min}(\mathsf F_L)=\frac7{15}-\frac{\Lambda_L}{21}.
\tag{7.2}
\] Equations (6.8) and (7.2) give the strict and quantitative
coefficient bounds

\[
\boxed{
0<\lambda_{\min}(\mathsf F_L)+\frac{71}{105}
\le\frac1{21}\min\left\{
\frac{24L+4}{L(2L+1)},\
2(1-c)(5+c)\right\}
\le\frac{2\pi^2}{7(2L-1)^2}.}
\tag{7.3}
\] Here \(7/15-24/21=-71/105\) exactly. In particular \[
\boxed{\lim_{L\to\infty}\lambda_{\min}(\mathsf F_L)=-\frac{71}{105}.}
\tag{7.4}
\]

The constant equal-plane bulk mode supplies the limiting edge of this
coefficient matrix, while the finite open maximizer has the strictly
smaller value in (6.7). These statements concern the explicit finite
Taylor-coefficient matrices and their large-box limit. They do not
control any remainder of the physical perturbation series uniformly in
volume, and they do not establish an infinite-volume gap or a
four-dimensional continuum limit.

\subsection{8. Lightweight exact
verification}\label{magnetic-part-6-lightweight-exact-verification}

The companion check\_first\_excitation\_face\_graph.py enumerates every
face and adjacency of the full open boxes \(L=2,3,4,5\). It checks the
six exact degree classes, connectedness, all edge-star multiplicities,
the constant-vector quotient, and all boundary deficits. The four graph
sizes are respectively:

{\def\LTcaptype{none} % do not increment counter
\begin{longtable}[]{@{}rrr@{}}
\toprule\noalign{}
\(L\) & Faces & Unordered adjacencies \\
\midrule\noalign{}
\endhead
\bottomrule\noalign{}
\endlastfoot
2 & 240 & 1128 \\
3 & 756 & 3852 \\
4 & 1728 & 9168 \\
5 & 3300 & 17940 \\
\end{longtable}
}

The checker directly computes the quadratic form of an integer-valued
separable interior trial on the unchanged whole open graph and verifies
the exact factorization in (6.6) before the sine substitution. For the
sine itself it uses the exact polynomials \(f_r/\sin t=U_{r-1}(c)\),
where \(U_r\) are the Chebyshev polynomials and \(U_\ell(c)=0\). It
verifies \(T=cS\) and the full trial quotient by polynomial division
with zero remainder, not by a numerical trigonometric eigenvalue
estimate.

An independently generated periodic graph with periods \((3,4,5)\)
checks the full translation list and twelve-neighbor multiplicity. Exact
Laurent variables \(x_i=e^{ik_i/2}\) verify the Fourier symbol and
center-phase change; at \(x_i=1\) the exact characteristic polynomial is
\((\lambda-24)(\lambda-12)^2\). The rational coefficient \(-71/105\) is
checked exactly.

The all-\(L\) graph statements, the infinite/periodic symbol, and both
open variational bounds are proved above; finite checks are diagnostics,
not a substitute for those proofs. No numerical eigenvalue, spatial
continuum approximation, remote action, new agent, PDF, or heavy job was
used.

\clearpage
% Complete next live proof; its original equation tags are retained.
\section{The exact center-disorder endpoint of the native magnetic
translation}\label{magnetic-part-7-the-exact-center-disorder-endpoint-of-the-native-magnetic-translation}

8 September 2026. This note calculates the relation between the
continuous native magnetic translation, discrete center transformations,
Wilson loops, and actual quantum spectral moments. It uses the full
original open SU(2) Wilson box. In particular a discrete transformation
is not identified with a nonzero excitation merely because it is not the
identity on link variables. Its action on the actual vacuum and its
energy are calculated explicitly.

\subsection{1. Original operator and all center
cochains}\label{magnetic-part-7-original-operator-and-all-center-cochains}

Fix the open cubical complex with vertices \(\{-L,\ldots,L\}^3\),
\(L\geq2\), positive physical links \(\mathsf E_L\), and all elementary
faces \(\mathsf P_L\). Reverse traversal uses the inverse of the same
link. For \(p=(n;i,j)\), \(i<j\), set \[
 W_p=\operatorname{tr}\{U_i(n)U_j(n+\mathbf e_i)
                  U_i(n+\mathbf e_j)^{-1}U_j(n)^{-1}\}.
\] Work in product Haar probability \(L^2(SU(2)^{\mathsf E_L},dU)\),
with \(T_a=-i\sigma_a/2\), left derivatives \(X_{e,a}\), and
\(E_e=-\sum_aX_{e,a}^2\). Retain \[
 H=\kappa\sum_eE_e+b\sum_p(2-W_p),\qquad
 \kappa=\frac{2g^2}{a},\quad b=\frac1{2g^2a},\quad
 \xi=\frac b\kappa=\frac1{4g^4},\quad a,g>0.
\tag{1}
\] Its domain is the product Sobolev space \(H^2\), its form domain is
\(H^1\), and its potential is bounded and smooth. The positive unit
vacuum \(\psi\) exists and is unique: modulus decreases form energy,
elliptic regularity and the maximum principle give strict positivity,
and the identity \[
 q_{H-\mathcal E}(\psi F)=\kappa\sum_{e,a}\int\psi^2|X_{e,a}F|^2dU
\tag{2}
\] proves uniqueness by dividing another bottom eigenvector by \(\psi\).
The gauge action \(U_e\mapsto g_{s(e)}U_eg_{t(e)}^{-1}\) preserves the
operator and hence fixes \(\psi\). Put \(A=H-\mathcal E\).

Let \(z=(z_e)\), \(z_e\in\{+1,-1\}\), be any real center one-cochain on
the original physical edges; assign the same sign to the reverse edge,
which is its inverse in this group. Define \[
 (D_zF)(U)=F((z_eU_e)_e),\qquad
 (\delta z)_p=\prod_{e\in\partial p}z_e,
 \quad\mathcal T_z=\{p:(\delta z)_p=-1\},
 \quad S_z=\sum_{p\in\mathcal T_z}W_p.
\tag{3}
\] Here \(z_eU_e\) means multiplication by \(z_e I_2\). Haar invariance
gives \(D_z^*=D_z=D_z^{-1}\). These diffeomorphisms preserve all integer
Sobolev domains, commute with all link Casimirs, and commute with every
vertex gauge transformation because \(z_eI_2\) is central. They obey
\(D_zD_w=D_{zw}\).

Each face word gains exactly its four central factors, so, as operator
identities on \(H^2\), \[
 D_zW_pD_z=(\delta z)_pW_p,\qquad
 \boxed{D_zHD_z=H+2bS_z.}
\tag{4}
\] All faces outside \(\mathcal T_z\) remain present in both
Hamiltonians. The scalar \(2b|\mathsf P_L|\) is identical on both sides
and has not been suppressed. Equation (4) is a unitary relation to an
explicitly sign-twisted potential, not permission to erase that
potential change.

\subsection{2. Exact native magnetic endpoint and its Wilson
algebra}\label{magnetic-part-7-exact-native-magnetic-endpoint-and-its-wilson-algebra}

The native smooth field in the retained cusp coordinates has link
transports \[
 h_1(n)=\exp(-\theta n_2H_c),\qquad h_2=h_3=I,
 \qquad H_c=i\sigma_3=-2T_3,\qquad
 \theta=\frac{2\pi a^2}{D_T},\quad
 D_T=L_Tq_T+6m_T^2>0.
\tag{5}
\] Its translation is \((T_\theta F)(U)=F((h_e^{-1}U_e)_e)\), and its
physical translated state is obtained by vertex Haar projection \(\Pi\).
At the specified parameter value \(\theta=\pi\), \[
 h_1(n)=(-1)^{n_2}I,
 \quad z^{\rm nat}_{(n,1)}=(-1)^{n_2},
 \quad z^{\rm nat}_{(n,2)}=z^{\rm nat}_{(n,3)}=1,
 \quad T_\pi=D_{z^{\rm nat}}.
\tag{6}
\] Thus \(\Pi T_\pi\psi=T_\pi\psi\) exactly. Its face coboundary is \[
 (\delta z^{\rm nat})_{12}=-1,\qquad
 (\delta z^{\rm nat})_{13}=(\delta z^{\rm nat})_{23}=1.
\tag{7}
\] The two direction-1 signs on a 12 face are \((-1)^{n_2}\) and
\((-1)^{n_2+1}\), whereas they agree on a 13 face. There are
\((2L)^2(2L+1)\) twisted 12 faces, not all M faces. At \(\theta=2\pi\)
every h is the identity. Equations (5)--(7) keep both the continuous
path and its discrete endpoints. The endpoint \(\theta=\pi\) means
precisely \(D_T=2a^2\); it is not an assertion that the small-angle cusp
exhaustion reaches that endpoint with its original spacing and
determinant.

For any closed oriented edge word C, allowing repeated traversals, let
\(W_C\) be the fundamental trace of its ordered holonomy and put
\(z(C)=\prod_{e\text{ traversed in }C}z_e\), with multiplicities.
Centrality and inversion give \[
 \boxed{D_zW_C=z(C)W_CD_z.}
\tag{8}
\] If a mod-two face chain \(\Sigma\) has boundary C modulo two, every
internal edge occurs twice, so \[
 z(C)=\prod_{p\in\Sigma}(\delta z)_p
     =(-1)^{|\Sigma\cap\mathcal T_z|}.
\tag{9}
\] This proves the discrete Wilson--disorder commutation pairing,
including its multiplicities and sign, directly in the original Hilbert
space. For (6), a rectangular 12 loop of integer side lengths r,s has
\(z^{\rm nat}(C)=(-1)^{rs}\). Its sign does not involve a rescaled area:
the physical area is \(a^2rs\), and (5) at \(\theta=\pi\) gives the same
phase \(\exp(\pi rsH_c)=(-1)^{rs}I\).

\subsection{3. Vacuum overlap, both spectral moments, and the exact zero
branch}\label{magnetic-part-7-vacuum-overlap-both-spectral-moments-and-the-exact-zero-branch}

For every center cochain set \[
 c_z=\langle\psi,D_z\psi\rangle,\qquad
 \chi_z=D_z\psi-c_z\psi,\qquad d_z=\|\chi_z\|^2=1-c_z^2.
\tag{10}
\] Strict positivity of \(\psi\) implies \(0<c_z\leq1\). All vectors in
(10) are smooth and gauge invariant, and \(\chi_z\perp\psi\). Using (4),
without replacing the vacuum, gives \[
 \boxed{q_A(\chi_z)=2b\langle S_z\rangle_\psi,\qquad
        A\chi_z=-2bS_zD_z\psi,\qquad
        \|A\chi_z\|^2=4b^2\langle S_z^2\rangle_\psi.}
\tag{11}
\] The first formula follows from
\(\langle D_z\psi,HD_z\psi\rangle-mathcal E
=\langle\psi,(D_zHD_z-H)\psi\rangle\); centering has no effect on the
form since \(A\psi=0\). For the second,
\(H D_z\psi-D_zH\psi=(H-D_zHD_z)D_z\psi=-2bS_zD_z\psi\). Finally
\(D_zS_zD_z=-S_z\), hence \(D_zS_z^2D_z=S_z^2\), proving the last
formula after Haar substitution.

If \(\mathcal T_z\) is empty, (4) commutes with H and uniqueness and
positivity of its unit vacuum imply \(D_z\psi=\psi\). Hence \(c_z=1\)
and \(\chi_z=0\) exactly. Conversely, if \(\chi_z=0\), the norm and
positive overlap force \(D_z\psi=\psi\). Apply (4) to this vector and
its original eigen-equation to get \(S_z\psi=0\). Strict positivity
makes \(S_z=0\) pointwise. Distinct elementary \(W_p\)'s are orthonormal
under product Haar measure: integrate a link in only one boundary, and
use \(\int(\operatorname{tr}U)^2dU=1\) when the faces agree. Therefore
\(\|S_z\|_{L^2(dU)}^2=|\mathcal T_z|\), and the latter must be zero. We
have proved \[
 \boxed{\chi_z=0\quad\Longleftrightarrow\quad\mathcal T_z=\varnothing.}
\tag{12}
\] This also proves \(\langle S_z\rangle_\psi>0\) when
\(\mathcal T_z\ne\varnothing\): (11) is the strictly positive form of a
nonzero physical vacuum-orthogonal vector of a compact-resolvent
Hamiltonian with a unique bottom.

For a nonempty twist the full physical spectral probability is
\(\nu_z(B)=d_z^{-1}\langle\chi_z,\mathbf1_B(A)\chi_z\rangle\). It has no
zero atom, and (11) gives its exact moments \[
 \boxed{\int\omega\,d\nu_z=
   \frac{2b\langle S_z\rangle_\psi}{1-c_z^2},\qquad
 \int\omega^2\,d\nu_z=
   \frac{4b^2\langle S_z^2\rangle_\psi}{1-c_z^2}.}
\tag{13}
\] The second moment retains every mixed \(\langle W_pW_q\rangle_\psi\)
between twisted faces. No independence is assumed. The denominator is
the exact overlap loss, not an assumed unit excitation norm.

\subsection{4. Small-coupling coefficient in the same
vacuum}\label{magnetic-part-7-small-coupling-coefficient-in-the-same-vacuum}

At each fixed box, the actual analytic positive vacuum of
\(H=\kappa(H_0-\xi\sum_pW_p)+2bM\) has \[
 \psi_\xi=1+\frac\xi3\sum_pW_p+O_{H^2,L}(\xi^2).
\tag{14}
\] One proof is to integrate the resolvent on \(|z|=3/8\) in the full
Hilbert space: the first free eigenvalue is \(3/4\),
\(\|\sum W_p\|=2M\), and a Neumann series converges for
\(|\xi|<3/(16M)\). Differentiating the eigen-equation gives (14),
because \(H_0W_p=3W_p\). Thus \[
 (D_z-I)\psi_\xi=-\frac{2\xi}{3}S_z+O_{H^2,L,z}(\xi^2).
\] For an involution and a real unit vector,
\(1-c_z=\tfrac12\|(D_z-I)\psi_\xi\|^2\). Consequently \[
 \begin{split}
 1-c_z&=\frac{2|\mathcal T_z|}{9}\xi^2+O_L(\xi^4),\\
 d_z&=\frac{4|\mathcal T_z|}{9}\xi^2+O_L(\xi^4),\\
 q_A(\chi_z)&=\frac{4|\mathcal T_z|}{3}\kappa\xi^2
                         +O_L(\kappa\xi^4),\\
 \|A\chi_z\|^2&=4|\mathcal T_z|\kappa^2\xi^2
                         +O_L(\kappa^2\xi^4).
 \end{split}
\tag{15}
\] To verify the even remainders, use the independent central cochain
\(\widetilde z_i(n)=(-1)^{\sum_{j<i}n_j}\), which changes every
elementary face sign. Its involution Z conjugates \(H_0-\xi\sum W_p\) to
\(H_0+\xi\sum W_p\), commutes with \(D_z\), and sends \(S_z\) to
\(-S_z\). Uniqueness gives \(\psi_{-\xi}=Z\psi_\xi\). Thus c\_z and
\(\langle S_z^2\rangle\) are even, whereas \(\langle S_z\rangle\) is
odd. Its first derivative from (14) is \(2|\mathcal T_z|/3\), and its
value at zero is zero. These facts give precisely all errors in (15).

For a nonempty fixed twist, division of the displayed quantities proves
the leading first moment \(3\kappa+O_L(\kappa\xi^2)\) and second moment
\(9\kappa^2+O_L(\kappa^2\xi^2)\). Equivalently the spectral probability
in the energy variable \(\omega/\kappa\) tends to \(\delta_3\) as
\(\xi\to0\): its integral of \((\omega/\kappa-3)^2\) tends to zero by
these moments, and the resulting Chebyshev bound proves weak
convergence. This is a fixed-box limit. It supplies the complete quantum
dictionary for (6), not a large-volume gap assertion.

\subsection{5. The original integer line flux and its adjoint center
class}\label{magnetic-part-7-the-original-integer-line-flux-and-its-adjoint-center-class}

There is a second exact relation, now for the original smooth bundle
rather than the finite-link endpoint. On the marked two-torus let the
original Hermitian line bundle have transition phases
\(e^{i\phi_{ab}}\), with \(\phi_{ab}+\phi_{bc}+\phi_{ca}=2\pi n_{abc}\)
on triple overlaps. The canonical cusp connection on this line has
integer curvature period \((2\pi)^{-1}\int f=1\). Its embedded rank-two
SU(2) bundle is \(L\oplus L^{-1}\), with explicit transition matrices \[
 g_{ab}=\operatorname{diag}(e^{i\phi_{ab}},e^{-i\phi_{ab}})
        =\exp(\phi_{ab}H_c),\qquad g_{ab}g_{bc}g_{ca}=I.
\tag{16}
\] The product identity follows from the integer \(n_{abc}\);
determinant one is also exact. Under the adjoint representation,
\([!H_c,T_1!]=-2T_2\), \([!H_c,T_2!]=2T_1\), and H\_c fixes T\_3. Hence
\[
 \begin{split}
 \operatorname{Ad}_{e^{\phi H_c}}T_1
   &=\cos(2\phi)T_1-\sin(2\phi)T_2,\\
 \operatorname{Ad}_{e^{\phi H_c}}T_2
   &=\sin(2\phi)T_1+\cos(2\phi)T_2.
 \end{split}
\tag{17}
\] In the oriented real plane with complex coordinate \(v_1+iv_2\), this
is multiplication by \(e^{-2i\phi}\). Its winding/Chern number is
therefore \(-2c_1(L)\). More directly, the obstruction cocycle to
lifting the adjoint SO(3) transitions to SU(2) is identically \(+I\),
because (16) is already such a lift. Its \(\mathbb Z/2\) obstruction
class is zero. This proves the exact integer-to-center map for this
particular embedding, including the factor 2 and its orientation sign;
it does not erase the nonzero original integer curvature period.

The associated holonomy identities are \(\exp(2\pi H_c)=I\) and
\(\exp(\pi H_c)=-I\). Replacing (16) by a half-angle lift would change
its triple-overlap product to \((-1)^{n_{abc}}I\), precisely the center
obstruction. Such a change is not the original line bundle with its
original transition matrices. The continuous field, its integer line
data, the adjoint center class, and the finite center-disorder endpoint
are therefore related by the explicit maps (5)--(9), (16)--(17), rather
than by an assumed equality of their flux labels.

\subsection{6. Primary context and what the calculation
establishes}\label{magnetic-part-7-primary-context-and-what-the-calculation-establishes}

Wilson--disorder pairing and center flux are established gauge-theory
constructions. A primary reference is G. 't Hooft, \emph{On the phase
transition towards permanent quark confinement}, Nuclear Physics B
\textbf{138} (1978), 1--25, DOI
\url{https://doi.org/10.1016/0550-3213(78)90153-0}. The actual
compact-link operator and its pairing needed here are proved in
(3)--(9). For the SU(2) global-charge conventions, D. Tong, \emph{Gauge
Theory}, Sections 2.6.1--2.6.2,
\url{https://www.damtp.cam.ac.uk/user/tong/gaugetheory.html}, discusses
the factor between adjoint and fundamental charge quantization;
(16)--(17) give it directly in the original H\_c convention.

P. de Forcrand and L. von Smekal, \emph{'t Hooft Loops, Electric Flux
Sectors and Confinement in SU(2) Yang--Mills Theory}, Physical Review D
\textbf{66} (2002), 011504,
\url{https://doi.org/10.1103/PhysRevD.66.011504},
\url{https://arxiv.org/abs/hep-lat/0107018v2}, study flux-sector free
energies at finite temperature. Their finite-temperature free energy is
not substituted for the vacuum excitation moment (13). The operator
relation and both moments used here were calculated directly in the full
H.

The center endpoint gives a nonzero actual quantum state precisely for a
nonempty coboundary twist. A cocycle that commutes with H fixes its
unique finite-box vacuum and produces the zero centered vector. Nonempty
twists give the exact overlap-dependent moments (13), including every
correlation. These statements neither assume nor prove a continuum mass
gap. They identify the spectral quantities that the original magnetic
path actually reaches at its discrete endpoint, while retaining the
smooth integer-flux bundle and its explicit adjoint relation.

\clearpage
% Complete next live proof; its original equation tags are retained.
\section{The unchanged native magnetic vector and its open-face spectral
measure}\label{magnetic-part-8-the-unchanged-native-magnetic-vector-and-its-open-face-spectral-measure}

8 September 2026. This note proves an all-volume spectral-probability
limit for the actual native magnetic leading face vector. The full
finite-volume amplitude, the original coordinate-dependent coefficients,
every open boundary, and the affine map to the first physical excitation
coefficient are retained. The calculation concerns a sequence of finite
Taylor-coefficient matrices; the coupling-dependent spectral transport
is a separate calculation.

\subsection{1. Original finite objects and the exact physical
coefficient
map}\label{magnetic-part-8-original-finite-objects-and-the-exact-physical-coefficient-map}

Fix an integer \(L\ge2\), put \(m=2L\), and use the original vertices
\(\{-L,\ldots,L\}^3\). A positively directed physical edge \(e=(n,i)\)
exists when both \(n\) and \(n+\mathbf e_i\) lie in this set. A face
\(p=(n;i,j)\), \(i<j\), has oriented edge word \[
(n,i)^+(n+\mathbf e_i,j)^+(n+\mathbf e_j,i)^-(n,j)^-.
\tag{1.1}
\] Thus its two plane base coordinates belong to
\([-L,L-1]\cap\mathbb Z\), and its normal base coordinate belongs to
\([-L,L]\cap\mathbb Z\). All \(M_L=3m^2(m+1)\) such faces are included.
Distinct faces are adjacent exactly when they share a physical edge.
Write \(\mathsf A_L\) for the unsigned adjacency matrix, \(d_p\) for its
degrees, and \[
\mathsf Q_L=\operatorname{diag}(d_p)+\mathsf A_L
  :\mathbb C^{\mathsf P_L}\longrightarrow\mathbb C^{\mathsf P_L}.
\tag{1.2}
\] The inner product is counting measure, conjugate-linear in its first
argument. In particular no orientation sign is inserted in (1.2).

Each lattice edge belongs to at most four faces, and two distinct
elementary faces share at most one edge, so \(d_p\le12\). Directly
expanding squares gives \[
\begin{split}
\langle x,\mathsf Q_Lx\rangle
 &=\sum_{\{p,q\}\text{ adjacent}}|x_p+x_q|^2,\\
\langle x,(24I-\mathsf Q_L)x\rangle
 &=\sum_{\{p,q\}\text{ adjacent}}|x_p-x_q|^2
   +2\sum_p(12-d_p)|x_p|^2.
\end{split}
\tag{1.3}
\] Hence \(\operatorname{spec}(\mathsf Q_L)\subset[0,24]\). This proof
uses the unchanged open graph, including its degree-deficient faces.

For the original link variables \(U_e\in SU(2)\), retain \[
\begin{gathered}
W_{n;ij}(U)=\operatorname{tr}
 \bigl(U_i(n)U_j(n+\mathbf e_i)
 U_i(n+\mathbf e_j)^{-1}U_j(n)^{-1}\bigr),\\
T_a=-i\sigma_a/2,\qquad
X_{e,a}f=\left.\frac{d}{dt}f(\ldots,e^{tT_a}U_e,\ldots)\right|_{t=0},
\qquad E_e=-\sum_{a=1}^3X_{e,a}^2,\\
H=\kappa H_0+b\sum_p(2-W_p)=\kappa K_\xi+2bM_LI,
\quad H_0=\sum_eE_e,\quad K_\xi=H_0-\xi\sum_pW_p,\\
\kappa=\frac{2g_{\rm YM}^2}{a},\qquad
b=\frac{1}{2g_{\rm YM}^2a},\qquad
\xi=\frac{b}{\kappa}=\frac1{4g_{\rm YM}^4},\qquad a>0,\ g_{\rm YM}>0.
\end{gathered}
\tag{1.4}
\] Here \(a\) without a face subscript is the physical lattice spacing;
the face coefficient vector below is denoted \(\boldsymbol a_L\). The
original Lie metric has \(-\Delta_e^c=4E_e\). No factor in (1.4) is
changed. The gauge-invariant physical space is the invariant closed
subspace of product-Haar \(L^2(SU(2)^N)\). The operator and form domains
are its intersections with \(H^2\) and \(H^1\), respectively.

Define the linear face-to-state map \[
\mathcal W_L:\mathbb C^{\mathsf P_L}\longrightarrow\mathcal H_{\rm phys},
\qquad \mathcal W_Lx=\sum_px_pW_p.
\tag{1.5}
\] It is an isometry onto the free energy-three face space: for unequal
faces a link appearing in only one face changes the sign of the
integrand under \(U_e\mapsto-U_e\), proving orthogonality; for equal
faces their holonomy is Haar distributed and the fundamental character
has squared Haar norm one. Every face trace has fundamental Casimir
\(3/4\) on its four links, giving \(H_0\mathcal W_L=3\mathcal W_L\).

The original native translation has \[
\Gamma=\sum_{e=(n,1)}n_2^2E_e,
\qquad h_1(n)=\exp(-\theta n_2H_c),\quad h_2=h_3=I,
\qquad H_c=i\sigma_3=-2T_3,
\tag{1.6}
\] with \(\theta=2\pi a^2/D\), \(D=L_Tq_T+6m_T^2>0\),
\(L_T=\operatorname{Im}\tau_T\), \(q_T=-\operatorname{Im}\beta_T\), and
\(m_T=\operatorname{Im}\mu_T\). These are the same physical coordinates
and magnetic scale as in magnetic\_translation\_true\_vacuum.md. In
first\_physical\_excitation\_band.md the positive unit vacuum is proved
to satisfy \(\psi_\xi=1+\xi\sum_pW_p/3+O_{H^2,L}(\xi^2)\). Applying the
specified \(\Gamma:H^2\to L^2\), at fixed \(L\), gives \[
\Gamma\psi_\xi=\xi\mathcal W_L\boldsymbol a_L+O_{L^2,L}(\xi^2),
\quad
\boxed{
a_{12}(n)=\frac{n_2^2+(n_2+1)^2}{4},\qquad
a_{13}(n)=\frac{n_2^2}{2},\qquad a_{23}(n)=0.}
\tag{1.7}
\] Indeed a 12-face has direction-1 edge bases at heights \(n_2,n_2+1\),
a 13-face has two at height \(n_2\), and a 23-face has none; multiplying
the sum of squared heights by \((3/4)/3\) gives exactly (1.7). The
centered vector \((\Gamma-\langle\psi_\xi,\Gamma\psi_\xi\rangle)
\psi_\xi\) has the same leading coefficient because
\(\langle1,\mathcal W_L\boldsymbol a_L\rangle=0\) and its centering
scalar is \(O_L(\xi^2)\).

The same band source proves the true-vacuum second-order matrix \[
\mathsf F_L=\frac7{15}I-\frac1{21}\mathsf Q_L.
\tag{1.8}
\] That derivation retains the whole interaction, the shared-edge
spin-zero and spin-one virtual channels, and the vacuum subtraction.
This note calculates the spectral measure of (1.8) on precisely (1.7).

\subsection{2. The unscaled spectral measure and the complete native
amplitude}\label{magnetic-part-8-the-unscaled-spectral-measure-and-the-complete-native-amplitude}

For a Borel set \(B\subset\mathbb R\), define \[
\sigma_L(B)=\langle\boldsymbol a_L,
 \mathbf1_B(\mathsf Q_L)\boldsymbol a_L\rangle,
\qquad A_L=\sigma_L(\mathbb R)=\|\boldsymbol a_L\|^2,
\qquad \mu_L=\frac{\sigma_L}{A_L}.
\tag{2.1}
\] The finite-dimensional spectral theorem defines the indicated
orthogonal projections. No unit-amplitude replacement of the native
state occurs: \(\sigma_L\) is kept, and only its explicitly named
probability \(\mu_L\) is divided by its actual nonzero mass.

Write \(I_L=[-L,L-1]\cap\mathbb Z\), and set \[
s_r=\sum_{t\in I_L}t^r,
\qquad f(t)=\frac{t^2+(t+1)^2}{4}
 =\frac{t^2}{2}+\frac t2+\frac14.
\tag{2.2}
\] Both displayed expressions define the same original coefficient. The
exact finite sums needed below are \[
\begin{gathered}
s_0=2L,\quad s_1=-L,\quad s_2=\frac{L(2L^2+1)}3,
\quad s_3=-L^3,\quad
s_4=\frac{L(6L^4+10L^2-1)}{15},\\
\sum_{t\in I_L}f(t)^2=\frac{L(4L^4+1)}{40},\qquad
\sum_{t\in I_L}\left(\frac t2+\frac14\right)^2
 =\frac{L(4L^2-1)}{24}.
\end{gathered}
\tag{2.3}
\] Here is an elementary verification of all sums. The odd sums pair
\(t\) with \(-t\) except for \(-L\), giving \(s_1,s_3\). For the even
sums, \(s_{2r}=L^{2r}+2\sum_{t=1}^{L-1}t^{2r}\). The polynomials
\(P_2(z)=z(z+1)(2z+1)/6\) and \(P_4(z)=z(z+1)(2z+1)(3z^2+3z-1)/30\) obey
\(P_2(0)=P_4(0)=0\) and \(P_r(z)-P_r(z-1)=z^r\) for \(r=2,4\), by
multiplication. Telescoping proves the two even sums. Expanding the
retained polynomial \(f^2=t^4/4+t^3/2+t^2/2+t/4+1/16\) and the last
square in (2.3), then substituting \(s_0,\ldots,s_4\), proves the
remaining identities.

There are \(m(m+1)\) 12-faces at each permitted \(n_2\), and \(m^2\)
13-faces at each height from \(-L\) through \(L\), inclusive. Therefore
\[
\begin{split}
A_L
 &=m(m+1)\sum_{t=-L}^{L-1}
       \left(\frac{t^2+(t+1)^2}{4}\right)^2
   +\frac{m^2}{4}\sum_{t=-L}^{L}t^4\\
 &=\boxed{\frac45L^7+\frac65L^6+\frac23L^5
                  +\frac1{30}L^3+\frac1{20}L^2}.
\end{split}
\tag{2.4}
\] In particular \(A_L>0\), \(A_L\ge4L^7/5\), and
\(A_L/L^7\longrightarrow4/5\). The boundary heights \(n_2=L\) and
\(n_3=L\) have been included in this mass.

\subsection{3. Exact comparison in the same finite Hilbert
space}\label{magnetic-part-8-exact-comparison-in-the-same-finite-hilbert-space}

Introduce a comparison vector solely for estimating (2.1). On the common
base cube \(C_L=I_L^3\), put \[
g_L(n)=\mathbf1_{C_L}(n)\frac{n_2^2}{2},
\quad
\boldsymbol b_L(n;12)=g_L(n),\quad
\boldsymbol b_L(n;13)=g_L(n),\quad
\boldsymbol b_L(n;23)=0.
\tag{3.1}
\] This defines an entry on every original face; it is zero on the
remaining faces. Set
\(\boldsymbol r_L=\boldsymbol a_L-\boldsymbol b_L\). Its complete
nonzero-entry list is \[
\begin{array}{c|c}
\text{original face and base}&r_L(n;ij)\\ \hline
n\in C_L,\ ij=12&n_2/2+1/4\\
ij=12,\ n_3=L,\ n_1,n_2\in I_L&f(n_2)\\
ij=13,\ n_2=L,\ n_1,n_3\in I_L&L^2/2.
\end{array}
\tag{3.2}
\] Every other entry is zero. The three displayed sets are disjoint. No
term in (3.2) is deleted from the native vector or its measure.

Using (2.3) on these exact supports gives \[
\begin{split}
B_L:=\|\boldsymbol b_L\|^2
 &=\frac{m^2s_4}{2}
   =\frac45L^7+\frac43L^5-\frac2{15}L^3,\\
R_L:=\|\boldsymbol r_L\|^2
 &=m^2\frac{L(4L^2-1)}{24}
   +m\frac{L(4L^4+1)}{40}+\frac{m^2L^4}{4}\\
 &=\frac65L^6+\frac23L^5-\frac16L^3+\frac1{20}L^2,\\
C_L^{\rm cross}:=\langle\boldsymbol b_L,\boldsymbol r_L\rangle
 &=m^2\left(\frac{s_3}{4}+\frac{s_2}{8}\right)
   =-\frac23L^5+\frac16L^3,\\
A_L&=B_L+2C_L^{\rm cross}+R_L.
\end{split}
\tag{3.3}
\] The negative cross term is retained. These formulas give
\(B_L/A_L\to1\) and \(R_L/A_L\to0\). They follow from a literal
decomposition in \(\mathbb C^{\mathsf P_L}\), not a change of
coordinates, of face sets, or of physical background.

In the order \((12,13,23)\), use the two real plane vectors \[
v_{24}=\left(\frac23,\frac23,\frac23\right),\qquad
v_{12}=\left(\frac13,\frac13,-\frac23\right),\qquad
(1,1,0)=v_{24}+v_{12}.
\tag{3.4}
\] Their dot product is zero, and their squared lengths are \(4/3\) and
\(2/3\). Let \(\boldsymbol u_{\lambda,L}(n;\alpha)
=g_L(n)(v_\lambda)_\alpha\) on every original face. All three face types
exist on \(C_L\), so this is a literal orthogonal decomposition: \[
\boldsymbol b_L=\boldsymbol u_{24,L}+\boldsymbol u_{12,L},\qquad
\langle\boldsymbol u_{24,L},\boldsymbol u_{12,L}\rangle=0,
\qquad
\|\boldsymbol u_{24,L}\|^2=\frac23B_L,\quad
\|\boldsymbol u_{12,L}\|^2=\frac13B_L.
\tag{3.5}
\] These are not asserted to be finite-open-graph eigenvectors. The next
section bounds their exact eigen-equation errors.

\subsection{4. Every shift and the finite-difference/open-boundary
estimates}\label{magnetic-part-8-every-shift-and-the-finite-differenceopen-boundary-estimates}

Independently define the infinite graph on
\(\mathbb Z^3\times\{12,13,23\}\) by the same physical-edge sharing
rule. Its operator on counting-measure
\(\ell^2(\mathbb Z^3;\mathbb C^3)\) is
\(\mathsf Q_\infty=12I+\mathsf A_\infty\). Within plane \(ij\), the
neighbors have bases \(n\pm\mathbf e_i,n\pm\mathbf e_j\). For the
distinct planes, the entire list is \[
\begin{array}{c|c|c}
\text{source}&\text{target}&\text{target base, }a,b\in\{0,1\}\\ \hline
12&13&n+a\mathbf e_2-b\mathbf e_3\\
12&23&n+a\mathbf e_1-b\mathbf e_3\\
13&23&n+a\mathbf e_1-b\mathbf e_2.
\end{array}
\tag{4.1}
\] The reverse list has the negative offsets. To prove this directly,
two faces sharing direction \(i\) have their two direction-\(i\) edges
at their respective base and base plus their other plane direction.
Equating these two possible physical-edge bases gives exactly the four
offsets in each row. The inverse traversals in (1.1) change no
physical-edge equality. There are four coplanar neighbors and eight
cross-plane neighbors, all distinct, yielding degree twelve.

Let \(J_L\) extend a face vector by zero to this infinite counting
space. Since the original graph is its induced subgraph, \[
J_L^*\mathsf Q_\infty J_L=12I+\mathsf A_L,
\qquad
\boxed{\mathsf Q_L=J_L^*\mathsf Q_\infty J_L-\mathsf D_L},
\quad \mathsf D_L=\operatorname{diag}(12-d_p).
\tag{4.2}
\] Both maps in (4.2) have the displayed original domains, and \(J_L\)
is an isometry. The nonzero boundary matrix \(\mathsf D_L\) is kept.

For the zero-extended scalar \(g_L\) in (3.1) define \[
\mathcal D_{i,L}=\sum_{n\in\mathbb Z^3}
 |g_L(n+\mathbf e_i)-g_L(n)|^2.
\tag{4.3}
\] Separation of the three finite coordinate sums gives exactly \[
\begin{split}
\mathcal D_{1,L}=\mathcal D_{3,L}&=\frac m2s_4,\\
\mathcal D_{2,L}
 &=m^2\left[\frac{L^4+(L-1)^4}{4}
       +\sum_{t=-L}^{L-2}\left(t+\frac12\right)^2\right],\\
\sum_{t=-L}^{L-2}\left(t+\frac12\right)^2
 &=\frac{8L^3-12L^2+10L-3}{12},\\
\mathcal D_L^{\rm sum}
 :=\mathcal D_{1,L}+\mathcal D_{2,L}+\mathcal D_{3,L}
 &=\frac{14}{5}L^6-\frac43L^5+\frac{10}{3}L^4
                         -\frac23L^3-\frac2{15}L^2.
\end{split}
\tag{4.4}
\] For directions 1 and 3 the only changes are the two boundary planes,
each contributing \(ms_4/4\). In direction 2 the lower and upper jumps
contribute \(L^4/4\) and \((L-1)^4/4\); each internal difference is
\(((t+1)^2-t^2)/2=t+1/2\). Expanding the last finite sum using (2.3)
with its upper endpoint removed proves its displayed polynomial and then
the last line. Thus the jump terms have not been replaced by interior
derivatives.

For a constant plane vector \(v\in\mathbb C^3\) and any eigenpair \[
(12I+4\mathbf1\mathbf1^T)v=\lambda v,
\tag{4.5}
\] the residual of \(g_Lv\) in the infinite graph has, in component
\(\alpha\) and at base \(n\), precisely \[
\bigl((\mathsf Q_\infty-\lambda)(g_Lv)\bigr)_\alpha(n)
 =\sum_{\beta,\delta\in\mathsf T_{\alpha\beta}}
 v_\beta\bigl(g_L(n+\delta)-g_L(n)\bigr).
\tag{4.6}
\] There are twelve terms, with \(\mathsf T_{\alpha\alpha}\) the four
coplanar shifts and other sets given by (4.1). Indeed the removed
zero-shift matrix is diagonal 16 and off-diagonal 4, precisely the
matrix in (4.5). The two vectors (3.4) satisfy (4.5) at 24 and 12: their
component sums are 2 and 0, respectively.

Translation invariance of counting norm and
\(\|x+y\|^2\le2\|x\|^2+2\|y\|^2\) imply \[
\sum_n|g_L(n+\mathbf e_i-\mathbf e_j)-g_L(n)|^2
 \le2(\mathcal D_{i,L}+\mathcal D_{j,L}).
\tag{4.7}
\] Split the change into its two successive coordinate shifts to prove
this, and use translation isometry for each summand. The same bound
holds with either sign or with the entire offset reversed.

Apply Cauchy--Schwarz to the twelve terms in (4.6), sum over
\(n,\alpha\), and collect each fixed target type \(\beta=ij\), with
normal direction \(k\). The four same-plane shifts contribute exactly
\(2(\mathcal D_i+\mathcal D_j)\). The four shifts to one other plane
contribute at most \(3(\mathcal D_i+\mathcal D_k)\) and those to the
remaining plane at most \(3(\mathcal D_j+\mathcal D_k)\): their terms
are a zero shift, two single shifts, and one double shift bounded by
(4.7). Consequently the coefficient for \(|v_\beta|^2\) is at most \[
12[5(\mathcal D_i+\mathcal D_j)+6\mathcal D_k]
 \le72\mathcal D_L^{\rm sum}.
\] We have proved the explicit infinite residual estimate \[
\| (\mathsf Q_\infty-\lambda)(g_Lv)\|^2
 \le72\|v\|^2\mathcal D_L^{\rm sum}.
\tag{4.8}
\]

Here is the complete boundary correction on the original support. Define
\(\beta_L(t)=\mathbf1_{\{t=-L\}}+\mathbf1_{\{t=L-1\}}\) for
\(t\in I_L\). A face \(ij\) of base \(n\in C_L\) and normal direction
\(k\) has \[
12-d_{n;ij}=\beta_L(n_i)+\beta_L(n_j)+4\mathbf1_{\{n_k=-L\}}.
\tag{4.9}
\] For proof, each of its four edges has two transverse extension
directions, each with two incident faces except that a transverse
coordinate equal to \(-L\) or \(L\) permits only one. Summing the three
other incident faces over the four edges gives degree twelve, minus one
for either extreme plane base and minus four when the normal coordinate
is an outer boundary. On the common cube the only outer normal boundary
is \(-L\), giving (4.9). On the rest of the original face set the same
formula uses \(4\mathbf1_{\{n_k\in\{-L,L\}\}}\); our comparison vector
is zero there, while its native residual is already retained in (3.2).

For any \(v\), the exact squared norm of this boundary error is \[
\mathcal E_L(v):=\|\mathsf D_L(g_Lv)\|^2
 =\sum_{ij}|v_{ij}|^2\sum_{n\in C_L}\frac{n_2^4}{4}
  [\beta_L(n_i)+\beta_L(n_j)+4\mathbf1_{\{n_k=-L\}}]^2.
\tag{4.10}
\] In each plane type the nonzero summands lie in at most five
coordinate planes of \(m^2\) points each. The bracket is at most 6 and
\(n_2^4/4\le L^4/4\). Counting overlaps more than once is an upper
bound, so \[
\mathcal E_L(v)\le45m^2L^4\|v\|^2.
\tag{4.11}
\] Equations (4.2), (4.8), and the triangle inequality, with every
quantity in the same original counting norm, prove \[
\begin{gathered}
\varepsilon_{\lambda,L}
 :=\| (\mathsf Q_L-\lambda I)\boldsymbol u_{\lambda,L}\|,\\
\boxed{
\varepsilon_{\lambda,L}
 \le\sqrt{72\|v_\lambda\|^2\mathcal D_L^{\rm sum}}
          +\sqrt{\mathcal E_L(v_\lambda)}
 \le\|v_\lambda\|\mathcal K_L,}\qquad
\mathcal K_L=\sqrt{72\mathcal D_L^{\rm sum}}+\sqrt{45m^2L^4},\\
\varepsilon_{24,L}+\varepsilon_{12,L}
 \le\frac{2+\sqrt2}{\sqrt3}\mathcal K_L.
\end{gathered}
\tag{4.12}
\] Since (4.4) gives \(\mathcal D_L^{\rm sum}=O(L^6)\), all these
residual bounds are \(O(L^3)\); their squared norms are small relative
to \(A_L\), which has order \(L^7\). The explicit formulas, rather than
this order notation, are the finite-volume estimates used below.

\subsection{5. Quantitative spectral probability limit for the actual
native
vector}\label{magnetic-part-8-quantitative-spectral-probability-limit-for-the-actual-native-vector}

For a complex Lipschitz function \(h:[0,24]\to\mathbb C\), write
\(\|h\|_\infty\) and \(\operatorname{Lip}(h)\) for its supremum norm and
Lipschitz constant. Define the explicit finite numbers \[
\begin{split}
\mathcal N_L
 &=\frac{(\sqrt{A_L}+\sqrt{B_L})\sqrt{R_L}+|B_L-A_L|}{A_L},\\
\mathcal T_L
 &=\frac{\sqrt{B_L}}{A_L}\,
           \frac{2+\sqrt2}{\sqrt3}\mathcal K_L.
\end{split}
\tag{5.1}
\] All symbols on the right sides were calculated above without
asymptotic replacement. We claim, for every \(L\ge2\), \[
\boxed{
\left|\int h(q)\,d\mu_L(q)
       -\frac23h(24)-\frac13h(12)\right|
 \le\mathcal N_L\|h\|_\infty
       +\mathcal T_L\operatorname{Lip}(h).}
\tag{5.2}
\]

First, by (3.2), functional calculus, and Cauchy--Schwarz, \[
\left|\langle\boldsymbol a_L,h(\mathsf Q_L)\boldsymbol a_L\rangle
       -\langle\boldsymbol b_L,h(\mathsf Q_L)\boldsymbol b_L\rangle\right|
 \le\|h\|_\infty(\sqrt{A_L}+\sqrt{B_L})\sqrt{R_L}.
\tag{5.3}
\] For example expand the difference as
\(\langle\boldsymbol r_L,h(\mathsf Q_L)\boldsymbol a_L\rangle
+\langle\boldsymbol b_L,h(\mathsf Q_L)\boldsymbol r_L\rangle\), and
apply the operator-norm bound \(\|h(\mathsf Q_L)\|\le\|h\|_\infty\).
Second, the finite spectral theorem and the scalar Lipschitz inequality
give, for either \(\lambda=24,12\), \[
\|[h(\mathsf Q_L)-h(\lambda)I]\boldsymbol u_{\lambda,L}\|^2
 =\int|h(q)-h(\lambda)|^2
       d\langle\boldsymbol u_{\lambda,L},
         \mathbf1_{dq}(\mathsf Q_L)\boldsymbol u_{\lambda,L}\rangle
 \le\operatorname{Lip}(h)^2\varepsilon_{\lambda,L}^2.
\tag{5.4}
\] Use \(\boldsymbol b_L=\boldsymbol u_{24,L}+\boldsymbol u_{12,L}\) and
their exact orthogonality to obtain \[
\begin{split}
&\left|\langle\boldsymbol b_L,h(\mathsf Q_L)\boldsymbol b_L\rangle
       -B_L\left[\frac23h(24)+\frac13h(12)\right]\right|\\
&\hspace{15mm}\le\sqrt{B_L}\operatorname{Lip}(h)
   (\varepsilon_{24,L}+\varepsilon_{12,L}).
\end{split}
\tag{5.5}
\] Indeed the difference equals the inner product of \(\boldsymbol b_L\)
with the sum of the two error vectors in (5.4). This step bounds the
actual spectral cross terms; it does not assume that distinct
approximate eigenvectors have orthogonal spectral projections. Combine
(5.3), (5.5), (4.12), and
\(|\frac23h(24)+\frac13h(12)|\le\|h\|_\infty\), then divide by the
actual \(A_L\), proving (5.2).

Equations (2.4), (3.3), and (4.4) show directly
\(\mathcal N_L=O(L^{-1/2})\) and \(\mathcal T_L=O(L^{-1/2})\). In
particular the right side of (5.2) tends to zero for each Lipschitz
\(h\). To pass to all continuous \(h\), partition the compact interval
\([0,24]\) into sufficiently short intervals and use linear
interpolation of its endpoint values. Uniform continuity makes the
resulting Lipschitz function \(h_\eta\) satisfy
\(\|h-h_\eta\|_\infty\le\eta\). Both measures in (5.2) have total mass
one, so the substitution error is at most \(2\eta\). First let
\(L\to\infty\), then \(\eta\to0\). This proves \[
\boxed{\mu_L\ \Longrightarrow\ \frac23\delta_{24}+\frac13\delta_{12}.}
\tag{5.6}
\] The unscaled statement is simultaneously retained: (5.3)--(5.5) prove
\[
\begin{split}
\left|\int h\,d\sigma_L
       -A_L\left[\frac23h(24)+\frac13h(12)\right]\right|
 &\le A_L\bigl(\mathcal N_L\|h\|_\infty
                +\mathcal T_L\operatorname{Lip}(h)\bigr),\\
\boxed{L^{-7}\sigma_L\ \Longrightarrow\
              \frac8{15}\delta_{24}+\frac4{15}\delta_{12}.}
\end{split}
\tag{5.7}
\] The second convergence is of finite measures, whose total masses tend
to \(4/5\); it follows by multiplying (5.6) by \(A_L/L^7\). The original
\(\sigma_L\) itself has diverging total mass \(A_L\), which remains
exactly (2.4).

There is also a direct concentration estimate that does not merely test
fixed functions. Put \(p(q)=(q-12)(q-24)\). On \([0,24]\),
\(|q-12|\le12\), \(|q-24|\le24\), and \(|p(q)|\le288\) (the quadratic
has its interior extremum \(-36\) at 18 and endpoint values 288 and 0).
By the decomposition (3.2),(3.5), \[
\|p(\mathsf Q_L)\boldsymbol a_L\|
 \le12\varepsilon_{24,L}+24\varepsilon_{12,L}+288\sqrt{R_L}.
\tag{5.8}
\] For \(0<\delta\le6\), any point of \([0,24]\) at distance at least
\(\delta\) from both 12 and 24 has \(|p(q)|\ge\delta(12-\delta)\): below
12 the lower bound is the larger \(\delta(12+\delta)\), and between the
two roots the product \((q-12)(24-q)\) attains its minimum on the
permitted interval at either endpoint. Spectral integration of (5.8)
therefore yields \[
\boxed{
\mu_L\{q:\min(|q-12|,|q-24|)\ge\delta\}
 \le\frac{[12\varepsilon_{24,L}+24\varepsilon_{12,L}
                    +288\sqrt{R_L}]^2}
             {A_L\delta^2(12-\delta)^2}.}
\tag{5.9}
\] The residuals on the right may be replaced by their explicit bounds
(4.12). For fixed \(\delta\), the displayed upper bound is
\(O(L^{-1})\). The limit weights follow from (5.2), not from
concentration alone.

\subsection{6. Exact pushforward to the band coefficient and original
energy
units}\label{magnetic-part-8-exact-pushforward-to-the-band-coefficient-and-original-energy-units}

Let \(\phi(q)=7/15-q/21\). Since (1.8) is exactly the affine functional
calculus of (1.2), their spectral projections obey \[
\mathbf1_B(\mathsf F_L)=\mathbf1_{\phi^{-1}(B)}(\mathsf Q_L).
\tag{6.1}
\] One proof is to choose an orthonormal eigenbasis of the real
symmetric \(\mathsf Q_L\); each vector with eigenvalue \(q\) has the
same vector with eigenvalue \(\phi(q)\) for \(\mathsf F_L\). The two
indicator operators therefore agree on a basis. The exact unscaled and
probability measures of \(\mathsf F_L\) on the unchanged
\(\boldsymbol a_L\) are \[
\sigma_L^F=\phi_*\sigma_L,\qquad
\mu_L^F=\phi_*\mu_L,\qquad \sigma_L^F(\mathbb R)=A_L.
\tag{6.2}
\] In particular \[
\boxed{\begin{split}
\mu_L^F&\Longrightarrow
 \frac23\delta_{-71/105}+\frac13\delta_{-11/105},\\
L^{-7}\sigma_L^F&\Longrightarrow
 \frac8{15}\delta_{-71/105}+\frac4{15}\delta_{-11/105}.
\end{split}}
\tag{6.3}
\] For a Lipschitz test function on \([-71/105,7/15]\), (5.2) applies to
its composition with \(\phi\), with Lipschitz constant divided by 21.
Thus (6.3) includes an explicit finite-volume error bound.

The support bound (1.3) permits polynomial test functions, giving \[
\begin{gathered}
\lim_L\frac{\langle\boldsymbol a_L,\mathsf Q_L\boldsymbol a_L\rangle}{A_L}
 =\frac23\,24+\frac13\,12=20,\\
\lim_L\frac{\|\mathsf Q_L\boldsymbol a_L\|^2}{A_L}
 =\frac23\,24^2+\frac13\,12^2=432,\\
\lim_L\operatorname{Var}_{\mu_L}(q)=432-20^2=32,\\
\lim_L\frac{\langle\boldsymbol a_L,\mathsf F_L\boldsymbol a_L\rangle}{A_L}
 =\frac7{15}-\frac{20}{21}=-\frac{17}{35},\qquad
\lim_L\operatorname{Var}_{\mu_L^F}(F)=\frac{32}{441}.
\end{gathered}
\tag{6.4}
\] The nonzero limiting variance proves that the native vector's
coefficient measure does not concentrate at the lowest coefficient
alone.

For the original energy coefficients, a coefficient \(q\) corresponds to
\[
\kappa\bigl(3+\xi^2\phi(q)\bigr)
 =\frac{6g_{\rm YM}^2}{a}
       +\frac{\phi(q)}{8g_{\rm YM}^6a}.
\tag{6.5}
\] This identity preserves all energy scales. Its two limiting
coefficient values insert \(-71/105\) and \(-11/105\), respectively, in
(6.5). Equation (6.5) is the Taylor polynomial, not the full interacting
energy. The band source supplies the exact transported finite matrix \[
V_\xi^*(H-\mathcal E(\xi)I)V_\xi
 =\kappa\left(3I+\xi^2\mathsf F_L+\mathcal R_{L,\xi}\right),
\quad
\|\mathcal R_{L,\xi}\|
 \le\frac73\left(\frac{64M_L}{3}\right)^4|\xi|^4,
\quad |\xi|\le\frac3{128M_L}.
\tag{6.6}
\] Here \(V_\xi\) is the exact isometry from the free face eigenspace to
the true interacting first cluster, and the left side acts on that free
space using (1.5). Its full vacuum energy \(\mathcal E(\xi)\) includes
the original scalar \(2bM_L\); subtracting the same scalar from both
energies explains its cancellation in (6.6). The norm and coupling
restriction in (6.6) are stated, not replaced by a volume-uniform error.
All nonlinear Wilson terms remain in \(H\) and in its exact spectral
transport; none is estimated by setting a plaquette to its classical
value. The native centered magnetic vector and its projection through
\(V_\xi\) have their own coupling-dependent dressing, treated in the
companion transport calculation.

The leading-vector amplitude also transfers exactly under (1.5): \[
\|\xi\mathcal W_L\boldsymbol a_L\|^2=\xi^2A_L,
\qquad
\left\| -\frac23\theta^2\xi\mathcal W_L\boldsymbol a_L\right\|^2
 =\frac49\theta^4\xi^2A_L
 =\frac{4\pi^4a^8}{9D^4g_{\rm YM}^8}\,A_L.
\tag{6.7}
\] The scalar on the right uses the original \(\theta=2\pi a^2/D\) and
\(\xi=1/(4g_{\rm YM}^4)\) without changing either. The unscaled graph
and coefficient measures of these two specified leading vectors are
respectively \(\xi^2\sigma_L\), \(\xi^2\sigma_L^F\), and
\(4\theta^4\xi^2/9\) times those original \(\sigma_L,\sigma_L^F\). The
associated probabilities, when these scalars are nonzero, are exactly
\(\mu_L,\mu_L^F\). For the actual centered translated state, the source
identity is \(\chi_\theta=-2\theta^2
(\Gamma-\langle\psi_\xi,\Gamma\psi_\xi\rangle)\psi_\xi/3
+O_{H^1,L,\xi}(\theta^4)\). Together with (1.7) this gives \[
\chi_\theta=-\frac23\theta^2
 [\xi\mathcal W_L\boldsymbol a_L+O_{L^2,L}(\xi^2)]
 +O_{L^2,L,\xi}(\theta^4).
\tag{6.8}
\] The first remainder is the fixed-volume small-coupling remainder for
the centered generator vector; the second is the fixed-volume,
fixed-coupling small-angle remainder. Their dependencies are retained.
Equations (6.7)--(6.8) specify exactly which leading coefficient the
measure calculation addresses and retain the errors for the actual
state.

Thus (6.1)--(6.6) give an exact map from this face-measure limit to the
specified physical band coefficient, while retaining the finite-volume
remainder. Neither (5.6) nor (6.3) takes a fixed positive coupling to
infinite volume, takes the lattice spacing to zero, or constructs the
interacting physical continuum spectral measure.

\subsection{7. Reproducibility and calculation
record}\label{magnetic-part-8-reproducibility-and-calculation-record}

The companion check\_native\_magnetic\_band\_measure.py uses only the
Python standard library and exact integer/rational arithmetic. It
independently enumerates every original face and every physical-edge
star for \(L=2,3,4,5,6\), constructs the open adjacency, and verifies
(1.7), (2.3)--(2.4), every entry of (3.2), all norms and the negative
cross term in (3.3), the orthogonal weights (3.5), every finite
difference (4.4), the boundary identity (4.9), the exact boundary sum
(4.10), and rational squared versions of the residual bounds. It checks
the native graph moments directly and verifies the rational maps and
limiting moments in (6.3)--(6.4). Polynomial identities are additionally
checked by exact coefficient operations or by their exact finite-sum
difference identities; no floating eigenvalue calculation is used as a
proof of the limit.

The checker was executed successfully with
\texttt{python\ agent\_work/ym\_gap\_primary\_20260908/wilson\_variation/check\_native\_magnetic\_band\_measure.py}.
Its verified native amplitudes were
\(201,11151/4,18708,333355/4,285129\) for \(L=2,3,4,5,6\), respectively;
the corresponding full face counts were \(240,756,1728,3300,5616\). All
polynomial, physical-edge, finite-difference, boundary, and residual
identities passed.

Task continuity: the existing root logbook and user-input provenance
files were found before writing; this owned calculation had no
preexisting native\_magnetic\_band\_measure.md or checker to overwrite.
The three source notes named in the task were inspected, and all
original native coefficients were retained. The auxiliary common-cube
vector was introduced only with its full residual and cross term. The
candidate two-atom probability was then proved from the finite open
graph, rather than assumed from its bulk symbol. No new agent, remote
write, PDF, Lean process, or spectral numerical diagonalization is part
of this bounded calculation.

\clearpage
% Complete next live proof; its original equation tags are retained.
\section{Native magnetic state: exact physical band transport and
spectral
weights}\label{magnetic-part-9-native-magnetic-state-exact-physical-band-transport-and-spectral-weights}

8 September 2026. This proof uses the actual interacting vacuum of the
full open three-dimensional spatial SU(2) box. The angular limit is
taken first, at each fixed box and coupling; the coupling limit is taken
afterwards. All statements about spectral measures retain their
unnormalized mass. Time is Hamiltonian time. The original spatial
spacing, magnetic angle, electric coefficient, and both shared-edge spin
channels remain explicit.

\subsection{1. Configuration, domains, and the actual native
state}\label{magnetic-part-9-configuration-domains-and-the-actual-native-state}

Fix an integer \(L\geq2\). The vertex set is
\(\mathsf V_L=\{-L,\ldots,L\}^3\). Include every positive link
\(e=(n,i)\) whose two endpoints lie in the box and every elementary face
\(p=(n;i,j)\), \(i<j\), whose four vertices lie there. Denote these sets
by \(\mathsf E_L,\mathsf P_L\). Their cardinalities are \[
 N=3(2L)(2L+1)^2,\qquad M=3(2L)^2(2L+1).
\tag{1.1}
\] Negative traversal uses the inverse of the same link variable. On the
compact connected configuration manifold \(\mathsf Q_L=SU(2)^N\), use
product Haar probability measure and an inner product conjugate-linear
in its first entry. Define \[
\begin{split}
 W_{n;i,j}(U)&=\operatorname{tr}\bigl[
 U_i(n)U_j(n+\mathbf e_i)U_i(n+\mathbf e_j)^{-1}U_j(n)^{-1}\bigr],\\
 T_a&=-i\sigma_a/2,\qquad
 X_{e,a}f=\left.\frac{d}{dt}f(\ldots,e^{tT_a}U_e,\ldots)\right|_{t=0},\\
 E_e&=-\sum_{a=1}^3X_{e,a}^2,\qquad H_0=\sum_eE_e,\qquad S=\sum_pW_p,\\
 K_\xi&=H_0-\xi S,\qquad H=\kappa K_\xi+2bM I,\\
 \kappa&=\frac{2g_{\rm YM}^2}{a},\qquad
 b=\frac1{2g_{\rm YM}^2a},\qquad
 \xi=\frac b\kappa=\frac1{4g_{\rm YM}^4}.
\end{split}
\tag{1.2}
\] Here \(a>0\), \(g_{\rm YM}>0\). In particular the Lie metric
\(c(A,B)=-\operatorname{tr}(AB)/2\) has orthonormal frame \(2T_a\), so
\(-\Delta_e^c=4E_e\). This factor is included in \(\kappa\). The traces
are real on SU(2).

Let \(\Pi\) be the Haar average of the vertex gauge action
\(U_e\mapsto g_{s(e)}U_eg_{t(e)}^{-1}\). Haar invariance proves
\(\Pi^* =\Pi^2=\Pi\); its range \(\mathcal H_{\rm phys}\) is the
physical Hilbert space. Each \(E_e\) is bi-invariant and every \(W_p\)
is gauge invariant, so these operators preserve this subspace. All
subsequent spectral projections are on \(\mathcal H_{\rm phys}\).

The elliptic product Laplacian has operator domain \(H^2\), form domain
\(H^1\), and compact resolvent. We use the intersections with
\(\mathcal H_{\rm phys}\). Multiplication by the bounded smooth real
function \(S\) leaves these domains and self-adjointness unchanged. A
ground minimizer may be taken nonnegative because taking absolute value
does not increase the gradient form; this follows by differentiating
\(\sqrt{|f|^2+\varepsilon^2}\) and passing to the limit. Elliptic
regularity and the strong maximum principle make the minimizer smooth
and strictly positive. For such a unit vector \(\psi_\xi\), with
\(K_\xi\psi_\xi=e(\xi)\psi_\xi\), integration by parts gives \[
 \langle\psi_\xi F,(H-\mathcal E_\xi)\psi_\xi F\rangle
 =\kappa\sum_{e,a}\int\psi_\xi^2|X_{e,a}F|^2dU,
 \qquad\mathcal E_\xi=\kappa e(\xi)+2bM.
\tag{1.3}
\] The positive smooth function and its reciprocal are bounded on this
compact manifold. The identity therefore extends to the form domain by
density. A second ground eigenvector divided by \(\psi_\xi\) would have
all derivatives zero and hence be constant. Thus the positive unit
vacuum is unique, and gauge transformations fix it. The full and
physical bottoms coincide. Write the exact physical excitation operator
as \[
 \mathscr A_\xi=H-\mathcal E_\xi I
       =\kappa(K_\xi-e(\xi)I).
\tag{1.4}
\] This equality subtracts the same \(2bM\) present in both energies.

The native nonwrapping links and their translation are \[
 h_1(n)=\exp(-\theta n_2H_c),\quad h_2(n)=h_3(n)=I,
 \quad H_c=i\sigma_3=-2T_3,\quad
 (\mathcal T_\theta f)(U)=f((h_e^{-1}U_e)_e).
\tag{1.5}
\] The original geometric dictionary is \(\theta=2\pi a^2/D\), with the
original positive determinant \(D\). No integer coordinate is wrapped,
translated to a different origin, or rescaled here. Define \[
 c_{\xi,\theta}=\langle\psi_\xi,\mathcal T_\theta\psi_\xi\rangle,
 \qquad
 \chi_{\xi,\theta}=\Pi\mathcal T_\theta\psi_\xi
                      -c_{\xi,\theta}\psi_\xi.
\tag{1.6}
\] This is a smooth physical vector orthogonal to the same actual
vacuum: \(\langle\psi_\xi,\Pi\mathcal T_\theta\psi_\xi\rangle
=c_{\xi,\theta}\) since \(\Pi\psi_\xi=\psi_\xi\).

For completeness the exact full-representation projection in (1.6) is as
follows. There is only one nonidentity translated edge at a given
source. Gauge conjugation conjugates its left translation by the gauge
variable at that source. Independent source averages therefore give, on
invariant inputs, \[
 \Pi\mathcal T_\theta\Pi=\mathcal C_\theta\Pi,
 \qquad \mathcal C_\theta=\prod_{e=(n,1)}C_{e,\theta n_2}.
\tag{1.7}
\] The operator \(C_{e,s}\) averages left translation by conjugates of
\(e^{sH_c}\). On an edge-spin-\(j\) matrix coefficient Schur's lemma and
its trace give exactly \[
 c_j(s)=\frac1{2j+1}\sum_{m=-j}^j e^{2ims}.
\tag{1.8}
\] Indeed the averaged representation matrix commutes with all matrices
of the irreducible representation and its trace is the character. The
finite sum defines it at every angle, including zeros of the alternative
sine denominator. All link-spin sectors and all invariant multiplicities
are retained. This proves that \(\mathcal C_\theta\) is even,
self-adjoint, a contraction, and commutes with every link Casimir, with
\(\Pi\), and with Sobolev powers of \(1+H_0\). Conjugacy under
\(i\sigma_1\) sends \(s\) to \(-s\), which also directly proves evenness
and self-adjointness.

Put \[
 w_{(n,1)}=n_2^2,\qquad w_{(n,2)}=w_{(n,3)}=0,
 \quad \Gamma=\sum_e w_eE_e,
 \quad\gamma_\xi=\langle\psi_\xi,\Gamma\psi_\xi\rangle,
 \quad v_\xi=(\Gamma-\gamma_\xi)\psi_\xi.
\tag{1.9}
\] In joint Peter--Weyl coordinates, \(\Gamma\) is the nonnegative
diagonal self-adjoint operator with eigenvalues \(\sum_ew_ej_e(j_e+1)\).
Its domain consists precisely of coefficient families square summable
after multiplication by these eigenvalues. It contains \(H^2\); equality
is not needed because some weights vanish. Since \(0\leq w_e\leq L^2\),
the diagonal comparison \(\Gamma\leq L^2H_0\) proves boundedness
\(\Gamma:H^{k+2}\to H^k\). In particular all applications to smooth
vacua, and all products below on such vectors, have their stated
domains.

Taylor's formula for the smooth translation flow, averaged over the
compact gauge group, is valid as a map \(H^{k+4}\to H^k\). The Haar
average of an adjoint-vector pair is its scalar inner product times
\(\delta_{ab}/3\), since rotational invariance makes the average scalar
and its trace fixes that scalar. The native generator at an edge has
coefficient \(2n_2T_3\). Distinct translated sources average
independently. Its averaged second derivative is therefore
\(-4\Gamma/3\). Evenness removes the odd derivatives. Consequently \[
\begin{split}
 \mathcal C_\theta\psi_\xi
  &=\psi_\xi-\frac23\theta^2\Gamma\psi_\xi+O_{H^k,L,\xi}(\theta^4),\\
 \chi_{\xi,\theta}
  &=-\frac23\theta^2v_\xi+O_{H^k,L,\xi}(\theta^4).
\end{split}
\tag{1.10}
\] This is a Sobolev vector expansion; it does not regard \(\Gamma\) as
a bounded operator on \(L^2\). We will first use (1.10) at fixed
\(\xi\).

The vector \(v_\xi\) is nonzero for every \(b>0\). If
\(\Gamma\psi_\xi=\gamma_\xi\psi_\xi\), Haar integration gives
\(0=\gamma_\xi\int\psi_\xi\), so \(\gamma_\xi=0\). Its quadratic form
then implies that \(\psi_\xi\) is independent of every direction-1 edge
with \(n_2\ne0\). Choose \(e=((0,1,0),1)\), which has four incident
faces for \(L\ge2\). Fix all other link variables at the identity. The
kinetic terms and \(\psi_\xi\) are independent of \(U_e\), whereas the
potential in the ground equation is \(4b(2-\operatorname{tr}U_e)\).
Smoothness permits this evaluation and positivity makes the fixed value
of \(\psi_\xi\) nonzero. The equation is impossible for \(b>0\). This
also proves that (1.10) has a strictly positive leading squared norm at
every fixed positive coupling.

\subsection{2. The exact first-band
isometry}\label{magnetic-part-9-the-exact-first-band-isometry}

Here are the spectral isolation and transport facts with their proof.
Peter--Weyl decomposition followed by the bounded gauge projection gives
a complete physical decomposition into link spins and invariant vertex
contractions. A nonempty support graph of such a contraction has no
degree-one vertex: a single nontrivial irreducible representation has no
invariant vector. Thus it contains a cycle. The cubic nearest-neighbor
graph is bipartite, so such a cycle has at least four edges, each
costing at least \(3/4\). If total energy is below \(9/2\), there are at
most five support edges. A four-cycle must then occur. A fifth edge
could neither have a new endpoint without a degree-one vertex nor be a
chord, since the potential chords have endpoints of the same bipartition
class. The support is a four-cycle, which is one elementary square. At
its degree-two vertices the two spins must agree and the invariant
contraction is unique by Schur's lemma. Thus its energy is \(4j(j+1)\).
Below \(9/2\) only \(j=1/2\) occurs. A six-edge rectangular loop
realizes \(9/2\). Single-link Haar integration also gives \[
 \langle W_p,W_q\rangle=\delta_{pq},\quad
 H_0W_p=3W_p,\quad
 \ker H_0=\mathbb C1,\quad
 \ker(H_0-3)=\mathcal P=\operatorname{span}\{W_p\},
 \quad \operatorname{spec}(H_0)\cap(3,9/2)=\varnothing.
\tag{2.1}
\] For unequal faces a link occurring in just one kills the integral by
central multiplication by \(-I\). For equal faces their holonomy is Haar
and the fundamental character has norm one: writing
\(U=u_0I+i\sum_au_a\sigma_a\) on the unit three-sphere gives
\(\int(2u_0)^2=1\). The same parity proves \(\int W_p=0\).

Let \(P\) be the projection onto \(\mathcal P\), and define \[
 R_3=(I-P)\bigl((H_0-3)|_{\operatorname{ran}(I-P)}\bigr)^{-1}(I-P).
\tag{2.2}
\] It acts by \(-1/3\) on constants; that channel is retained. Set \[
 r_L=\frac3{64M},\qquad \mathcal C_3=\{z:|z-3|=3/4\},\qquad
 P_\xi=\frac1{2\pi i}\int_{\mathcal C_3}(z-K_\xi)^{-1}\,dz.
\tag{2.3}
\] The identity \(\|S\|=2M\) follows from its pointwise bound and the
identity configuration with positive-measure neighborhoods. On
\(\mathcal C_3\), the free resolvent has norm at most \(4/3\). For
complex \(|\xi|\le r_L\) the Neumann factor is at most \(1/8\), and
hence \[
 \|(z-K_\xi)^{-1}\|\le\frac{32}{21},\qquad
 \|P_\xi-P\|\le\frac17.
\tag{2.4}
\] The latter follows by integrating the resolvent difference:
\((3/4)(4/3)^2(2Mr_L)/(1-1/8)=1/7\). These analytic idempotents have
rank \(M\) by continuous deformation from \(P\). On \(\mathcal P\), let
\[
 T_\xi=PP_\xi P|_{\mathcal P},\qquad
 V_\xi=P_\xi P T_\xi^{-1/2}:\mathcal P\longrightarrow\operatorname{ran}P_\xi.
\tag{2.5}
\] The convergent binomial series at \(I\) defines the inverse square
root since \(\|T_\xi-I\|\le1/7\). With \(U_\xi=T_\xi^{-1/2}PP_\xi\) on
\(\operatorname{ran}P_\xi\), idempotence gives \(U_\xi V_\xi=I\), and
equal dimensions give the reverse identity. For real \(\xi\), \(P_\xi\)
is orthogonal and \(U_\xi=V_\xi^*\), so \(V_\xi\) is an isometry onto
the actual spectral subspace. The full Gram factor is part of this exact
coordinate map. Its range is smooth by elliptic regularity and finite
rank. The contour also makes \(K_\xi P_\xi\) bounded via the integral
with integrand \(z(z-K_\xi)^{-1}\); all intertwining products are thus
defined.

For real \(|\xi|\le r_L/2\), min--max moves each ordered physical
eigenvalue by at most \(2M|\xi|\le3/64\). There is one bottom, then
exactly \(M\) eigenvalues near 3, and every other eigenvalue is at least
\(9/2-3/64\). Thus \(P_\xi\) is precisely the first physical band, and
\(P_\xi\psi_\xi=0\) exactly. Define its exact matrix \[
 B_\xi=U_\xi K_\xi V_\xi,
 \qquad \mathscr A_\xi V_\xi
      =V_\xi\kappa(B_\xi-e(\xi)I).
\tag{2.6}
\] For every bounded Borel function \(f\), functional calculus gives \[
 f(\mathscr A_\xi)P_\xi
 =V_\xi f\bigl(\kappa(B_\xi-e(\xi)I)\bigr)V_\xi^*.
\tag{2.7}
\] Indeed an orthonormal eigenbasis of the finite Hermitian matrix is
mapped by \(V_\xi\) to an orthonormal eigenbasis of its invariant
subspace; both sides act by the same scalar on those vectors and vanish
on its orthogonal complement. This proves the exact observable map as
well as the exact state-space map.

The bottom projection near zero has the analogous isolated contour, so
\(e(\xi)\) and the positively chosen unit \(\psi_\xi\) are analytic near
zero. The equation \[
 \psi_\xi=(H_0+1)^{-1}[(1+e(\xi)+\xi S)\psi_\xi]
\tag{2.8}
\] and boundedness of smooth multiplication in each Sobolev space
upgrade analyticity by two derivatives at a time. Thus every fixed
Sobolev Taylor expansion used below is justified at fixed \(L\).

\subsection{3. Every second-order channel and the coefficient
energies}\label{magnetic-part-9-every-second-order-channel-and-the-coefficient-energies}

For a repeated face and for an unordered adjacent pair with shared edge
\(e\), put \[
 Y_p=W_p^2-1,\qquad
 Z_{pq,0}=\int W_pW_q\,dU_e,\qquad Z_{pq,1}=W_pW_q-Z_{pq,0}.
\tag{3.1}
\] An adjacent pair means two distinct faces sharing one physical edge,
including coplanar and hinged pairs. All other unordered distinct pairs
are called edge-disjoint, even when they share a vertex. Reverse one
face traversal if necessary and cyclically write its pair of traces as
\(\operatorname{tr}(U_eA)\operatorname{tr}(U_e^{-1}B)\). The identity
\(\int U_{ij}\overline U_{kl}\,dU
=\delta_{ik}\delta_{jl}/2\) gives \(Z_{pq,0}=\operatorname{tr}(AB)/2\).
The six outer links are distinct. The shared-edge tensor product has
spins 0 and 1, while every outer edge has spin \(1/2\). Repeated-face
character multiplication gives spin 1 on its four edges. Consequently
the complete table is \[
\begin{array}{c|ccccc}
 \text{vector}&1&Y_p&W_pW_q\ (\text{edge-disjoint})&Z_{pq,0}&Z_{pq,1}\\\hline
 H_0\text{ eigenvalue}&0&8&6&9/2&13/2\\
 \text{squared norm}&1&1&1&1/4&3/4.
\end{array}
\tag{3.2}
\] For the last two norms the Haar projection in (3.1) is orthogonal,
its image is half a norm-one character, and the original product has
squared norm one by integrating the independent outer Haar words first.
The other norms follow from the same character integration.

Vectors belonging to different unordered distinct pairs are orthogonal.
Their odd link parity is exactly
\(\partial p\mathbin\triangle\partial q\), in either shared-edge
channel. If different pairs had the same parity, their symmetric
difference would be a nonempty closed mod-two set of at most four
squares. Any square in that set needs another square on each of its four
edges. Distinct elementary squares share at most one edge, so at least
four \emph{other} squares would be necessary, a contradiction.
Integrating a link of unequal parity proves orthogonality. Repeated-face
vectors have even parity everywhere and so are orthogonal to distinct
pairs; distinct \(Y_p\)'s have different link-spin assignments.
Constants are orthogonal to all these nonconstant vectors. All
orthogonalities remain valid after applying \(H_0\), \(R_3\), or
\(\Gamma\), since the displayed vectors have definite joint link spins.

Comparing coefficients in the vacuum equation and its unit norm gives \[
\begin{split}
 \psi_\xi&=1+\xi A+\xi^2 B^{\rm vac}+O_{H^k,L}(\xi^3),
   \qquad A=S/3,\\
 H_0B^{\rm vac}&=(S^2-M)/3,
   \qquad \int B^{\rm vac}=-M/18,\\
 B^{\rm vac}&=-\frac M{18}+\sum_p\frac{Y_p}{24}
 +\sum_{\{p,q\}\,\text{edge-disjoint}}\frac{W_pW_q}{9}
 +\sum_{\{p,q\}\,\text{adjacent}}
       \left(\frac4{27}Z_{pq,0}+\frac4{39}Z_{pq,1}\right),\\
 e(\xi)&=-\frac M3\xi^2+O_L(\xi^4).
\end{split}
\tag{3.3}
\] The third line follows by expanding \(S^2-M\) into the repeated-face
terms and twice every distinct unordered pair, and dividing by the
energies in (3.2). The bottom coefficient is
\(-\langle S,A\rangle=-M/3\). The even remainder follows from the parity
proved below.

Let \(\mathscr D=R_3SP:\mathcal P\to\mathcal P^\perp\). For any
\(x=\sum_px_pW_p\), table (3.2) gives its entire value: \[
\begin{split}
 \mathscr D x={}&-\frac13\sum_px_p
 +\frac15\sum_px_pY_p
 +\frac13\sum_{\{p,q\}\,\text{edge-disjoint}}(x_p+x_q)W_pW_q\\
 &+\sum_{\{p,q\}\,\text{adjacent}}(x_p+x_q)
           \left(\frac23Z_{pq,0}+\frac27Z_{pq,1}\right).
\end{split}
\tag{3.4}
\] Differentiating \(P_\xi^2=P_\xi\) gives \(PP'_0P=0\). Differentiating
(2.5) therefore gives \(PV'_0=0\). The derivative of the intertwining
equation \(K_\xi V_\xi=V_\xi B_\xi\), projected onto
\(\mathcal P^\perp\), gives \(V'_0=\mathscr D\). Its second-order
projection onto \(\mathcal P\) gives \[
 B_\xi=3I-\xi^2PSR_3SP+O_L(\xi^3).
\tag{3.5}
\] The first-order projected term vanishes: the central link cochain
\(z_i(n)=(-1)^{\sum_{j<i}n_j}I\) changes the sign of every face. It
defines a Haar unitary involution \(\mathcal Z\) commuting with all
\(E_e\), gauge transformations, and native translations. It fixes 1 and
is \(-I\) on \(\mathcal P\); hence \(PSP=0\). The same cochain proves \[
\begin{gathered}
 K_{-\xi}=\mathcal ZK_\xi\mathcal Z,\quad
 \psi_{-\xi}=\mathcal Z\psi_\xi,\quad e(-\xi)=e(\xi),\quad
 P_{-\xi}=\mathcal ZP_\xi\mathcal Z,\\
 T_{-\xi}=T_\xi,\qquad V_{-\xi}=-\mathcal ZV_\xi,
 \qquad B_{-\xi}=B_\xi.
\end{gathered}
\tag{3.6}
\] Thus the odd matrix coefficients in (3.5) vanish without choosing
labels inside a degenerate eigenspace.

Define the actual open-box face adjacency matrix \(A_L^{\rm face}\) by
shared-edge adjacency, its degrees \(d_p\), and \[
 Q_L=\operatorname{diag}(d_p)+A_L^{\rm face},\qquad
 F_L=\frac7{15}I-\frac1{21}Q_L.
\tag{3.7}
\] We identify matrices on \(\ell^2(\mathsf P_L)\) with matrices on
\(\mathcal P\) by the exact unitary \(\mathcal Jx=\sum_px_pW_p\),
justified by (2.1). The constant channel in \(-PSR_3SP\) contributes
\(+1/3\) to every matrix entry. A repeated face contributes \(-1/5\) to
its diagonal. Each edge-disjoint pair contributes \(-1/3\) to its two
diagonals and two mutual off-diagonals. Each adjacent pair contributes
to that same four-entry pattern \[
 -\frac{1/4}{9/2-3}-\frac{3/4}{13/2-3}
 =-\frac16-\frac3{14}=-\frac8{21}.
\tag{3.8}
\] Hence its diagonal is \(1/3-1/5-(M-1-d_p)/3-8d_p/21
=7/15-M/3-d_p/21\); its off-diagonal is \(-1/21\) for adjacent faces and
zero otherwise. Adding the actual vacuum subtraction \(M/3\) in (3.3)
proves \[
 B_\xi-e(\xi)I=3I+\xi^2F_L+\mathcal R_\xi,
 \qquad \|\mathcal R_\xi\|\le\frac73 r_L^{-4}|\xi|^4
 \quad (|\xi|\le r_L/2).
\tag{3.9}
\] Here is also the bound's proof. The binomial series in (2.5) bounds
\(\|T_\xi^{-1/2}\|\le\sqrt{7/6}\); the contour for \(B_\xi-3I\) gives
norm at most \((7/6)(3/4)^2(32/21)=1\). The bottom contour \(|z|=3/4\)
gives \(|e(\xi)|\le3/4\). Thus the even analytic matrix
\(B_\xi-e(\xi)I-3I\) has norm at most \(7/4\) on \(|\xi|=r_L\). Cauchy's
estimate and a geometric series bound its terms of degree at least four
by \((7/4)\sum_{k\ge2}(|\xi|/r_L)^{2k}
\le(7/3)r_L^{-4}|\xi|^4\). The contour and binomial inequalities are
strict at their singularity thresholds, so analyticity extends to a
neighborhood of that circle and the Cauchy estimate is legitimate.

For every face vector \(x\),
\(\langle x,Q_Lx\rangle=\sum_{\{p,q\}\,\text{adjacent}}|x_p+x_q|^2\).
Every edge has at most four incident faces, so \(d_p\le12\). Therefore
\(0\le Q_L\le24I\) as quadratic forms, by
\(|x_p+x_q|^2\le2|x_p|^2+2|x_q|^2\). In particular
\(\operatorname{spec}(F_L)\subset[-71/105,7/15]\), with each exact
coefficient energy \(7/15-q/21\) corresponding to a spectral value \(q\)
of the same finite \(Q_L\).

\subsection{4. Transport of the actual magnetic variance
vector}\label{magnetic-part-9-transport-of-the-actual-magnetic-variance-vector}

Define the original face weights \[
 a_p=\frac14\sum_{e\in\partial p}w_e,
 \qquad \mathbf a=\sum_pa_pW_p=\mathcal J(a_p)_p.
\tag{4.1}
\] In all three planes their exact values are \[
 a_{n;1,2}=\frac{n_2^2+(n_2+1)^2}{4},\qquad
 a_{n;1,3}=\frac{n_2^2}{2},\qquad a_{n;2,3}=0.
\tag{4.2}
\] Since each boundary link of \(W_p\) has Casimir \(3/4\),
\(\Gamma W_p=3a_pW_p\), so \(\Gamma A=\mathbf a\). The vacuum scalar in
(1.9) is \[
 \gamma_\xi=\frac{\xi^2}{3}\sum_pa_p+O_L(\xi^4),
\tag{4.3}
\] because \(\Gamma1=0\), \(\langle A,\Gamma A\rangle=\sum_pa_p/3\), and
(3.6) makes the scalar even. Let
\(\mathbf t=\Gamma B^{\rm vac}-(\sum_pa_p)/3\). Then \[
 v_\xi=\xi\mathbf a+\xi^2\mathbf t+O_{H^k,L}(\xi^3).
\tag{4.4}
\] To spell out every coefficient, put \(s_{pq}=a_p+a_q\) for a pair and
let \(w_e\) retain its original value on its shared edge. Its weighted
electric eigenvalues are \[
\begin{array}{c|cccc}
 \text{vector}&Y_p&W_pW_q\ (\text{edge-disjoint})&Z_{pq,0}&Z_{pq,1}\\\hline
 \Gamma\text{ eigenvalue}&8a_p&3s_{pq}&3s_{pq}-3w_e/2&3s_{pq}+w_e/2.
\end{array}
\tag{4.5}
\] For the shared edge, \(3s_{pq}\) counts its two fundamental values as
\(3w_e/2\). Replacing these by spin zero subtracts \(3w_e/2\); replacing
them by spin one adds \(w_e/2\). This proves the last two entries,
including direction-2 and direction-3 shared edges where \(w_e=0\).
Inserting (4.5) into (3.3) proves \[
\begin{split}
 \mathbf t={}&-\frac13\sum_pa_p+\sum_p\frac{a_p}{3}Y_p
 +\sum_{\{p,q\}\,\text{edge-disjoint}}\frac{s_{pq}}3 W_pW_q\\
 &+\sum_{\{p,q\}\,\text{adjacent}}
 \left[\frac{4s_{pq}-2w_e}{9}Z_{pq,0}
       +\frac{12s_{pq}+2w_e}{39}Z_{pq,1}\right].
\end{split}
\tag{4.6}
\]

The exact transported magnetic coordinates and exact remainder are \[
 y_\xi=V_\xi^*v_\xi=V_\xi^*\Gamma\psi_\xi\in\mathcal P,
 \qquad \rho_\xi=(I-P_\xi)v_\xi\in\mathcal H_{\rm phys},
 \qquad v_\xi=V_\xi y_\xi+\rho_\xi.
\tag{4.7}
\] The equality dropping the scalar in \(y_\xi\) follows from the exact
identity \(V_\xi^*\psi_\xi=0\); no free-vacuum orthogonality is
substituted. Both summands are orthogonal to \(\psi_\xi\), and the two
summands are mutually orthogonal. The first lies in the actual band, the
second in its orthogonal complement with its ground-state component
exactly zero. These maps are defined on all smooth vectors here, and
\(V_\xi^*\) also extends boundedly to the whole Hilbert space.

Equations (3.4), (4.4), and \(P\mathbf t=0\) give \[
\begin{split}
 y_\xi&=\xi\mathbf a+O_L(\xi^3),\\
 P_\xi v_\xi&=\xi\mathbf a+\xi^2\mathscr D\mathbf a
                             +O_{H^k,L}(\xi^3),\\
 \rho_\xi&=\xi^2\mathbf r+O_{H^k,L}(\xi^3),
 \qquad \mathbf r=\mathbf t-\mathscr D\mathbf a.
\end{split}
\tag{4.8}
\] For the first line, \(V_\xi=I+\xi\mathscr D+O(\xi^2)\),
\(\mathscr D\mathbf a\perp\mathcal P\), and \(P\mathbf t=0\) remove the
quadratic term. Alternatively (3.6), \(v_{-\xi}=\mathcal Z v_\xi\), and
\(V_{-\xi}=-\mathcal ZV_\xi\) give the exact odd identity
\(y_{-\xi}=-y_\xi\). Every other remainder assertion in (4.8) is in the
fixed Sobolev norm indicated, using (2.8) and finite-rank elliptic
regularity for \(V_\xi\).

Subtracting the full expressions (3.4) and (4.6) gives \[
\boxed{
 \mathbf r=\sum_p\frac{2a_p}{15}Y_p
 +\sum_{\{p,q\}\,\text{adjacent}}
   \left[-\frac{2(s_{pq}+w_e)}9Z_{pq,0}
         +\frac{2(3s_{pq}+7w_e)}{273}Z_{pq,1}\right].}
\tag{4.9}
\] The two constants both equal \(-\sum_pa_p/3\) and cancel. The
edge-disjoint coefficients both equal \(s_{pq}/3\) and cancel. The
repeated coefficient is \(a_p(1/3-1/5)=2a_p/15\). The adjacent
coefficients are \((4s-2w)/9-2s/3=-2(s+w)/9\) and
\((12s+2w)/39-2s/7=2(3s+7w)/273\). These calculations prove every
retained and cancelled channel. In particular the exact physical band is
not obtained by projecting with the free \(P\): its constant,
repeated-face, edge-disjoint, and both adjacent-face dressings in (3.4)
all enter before the subtraction.

Write the following exact nonnegative finite sums: \[
\begin{split}
 \mathsf N_L&=\|\mathbf a\|^2=\sum_pa_p^2,\\
 \mathsf R_{8,L}&=\frac4{225}\mathsf N_L,\\
 \mathsf R_{9/2,L}&=\frac1{81}
      \sum_{\{p,q\}\,\text{adjacent}}(s_{pq}+w_e)^2,\\
 \mathsf R_{13/2,L}&=\frac1{24843}
      \sum_{\{p,q\}\,\text{adjacent}}(3s_{pq}+7w_e)^2,\\
 \mathsf R_L&=\mathsf R_{8,L}+\mathsf R_{9/2,L}+\mathsf R_{13/2,L}.
\end{split}
\tag{4.10}
\] The coefficient \(1/24843\) is \((3/4)\cdot4/273^2\). Orthogonality
and (3.2) prove \(\|\mathbf r\|^2=\mathsf R_L\). In particular
\(\mathsf R_L\ge4\mathsf N_L/225>0\), since (4.2) has positive 12-face
entries. From coupling parity and (4.8), \[
\begin{split}
 \|v_\xi\|^2&=\xi^2\mathsf N_L+O_L(\xi^4),\qquad
 \|P_\xi v_\xi\|^2=\xi^2\mathsf N_L+O_L(\xi^4),\\
 \|\rho_\xi\|^2&=\xi^4\mathsf R_L+O_L(\xi^6),
 \qquad \|\rho_\xi\|=\xi^2\sqrt{\mathsf R_L}+O_L(\xi^4),\\
 \frac{\|\rho_\xi\|^2}{\|v_\xi\|^2}
   &=\xi^2\frac{\mathsf R_L}{\mathsf N_L}+O_L(\xi^4)
       \quad (\xi\downarrow0).
\end{split}
\tag{4.11}
\] Indeed \(\rho_{-\xi}=\mathcal Z\rho_\xi\); its squared norm is even
analytic, so its possible fifth coefficient vanishes. The square root
estimate follows because \(\mathsf R_L>0\). The divisions use
\(\mathsf N_L>0\) at this fixed box.

The mass \(\mathsf N_L\) retains the explicit original polynomial: \[
\boxed{
 \mathsf N_L=\frac{L^2(2L+1)}{60}
       (24L^4+24L^3+8L^2-4L+3).}
\tag{4.12}
\] To prove it directly, for each \(m=n_2\) there are \(2L(2L+1)\)
12-faces and \(4L^2\) 13-faces. Put \(r_m=m^2+m+1/2\), so the two
coefficients are \(r_m/2\) and \(m^2/2\). Hence \[
 \mathsf N_L=\frac14\left[
 2L(2L+1)\sum_{m=-L}^{L-1}r_m^2
 +4L^2\sum_{m=-L}^Lm^4\right].
\tag{4.13}
\] The exact sums are \(\sum_{m=-L}^{L-1}r_m^2=L(4L^4+1)/10\) and
\(\sum_{m=-L}^Lm^4=L(L+1)(2L+1)(3L^2+3L-1)/15\). The first polynomial
has value \(1/2\) at \(L=1\) and increment \(2(L^2+L+1/2)^2\); the
second has value 2 and increment \(2(L+1)^4\). These are exactly the
endpoint increments of the respective sums, proving them by induction.
Substitution proves (4.12). The distinct sum needed in (4.3) is \[
 \sum_pa_p=\frac{L^2(2L+1)(4L^2+2L+1)}3.
\tag{4.14}
\] Its proof uses the same counts and the sums
\(\sum_{m=-L}^{L-1}r_m=L(2L^2+1)/3\) and
\(\sum_{m=-L}^Lm^2=L(L+1)(2L+1)/3\), each verified by its value at 1 and
its two endpoint increments.

The full open-boundary evaluation of the leakage sum (4.10), proved by
the original adjacent-face classification in
\url{native_magnetic_leakage_volume_polynomial.md}, is \[
\boxed{\begin{split}
 \mathsf R_L=\frac1{251535375}\bigl(&153001392L^7+15594888L^6
 +129973560L^5-21634000L^4\\
 &-3201942L^3+4992787L^2-1917900L\bigr).
\end{split}}
\tag{4.15}
\] That proof evaluates each of the two distinct shared-edge sums
separately, retaining their factors \(1/81\) and \(1/24843\). With
\(\mathsf N_L=\|\mathbf a\|^2\) exactly as in (4.12), it gives \[
 \frac{\mathsf R_L}{\mathsf N_L}
 =\frac{12750116}{16769025}-\frac{54848}{51597L}
       +\frac{3233332}{2012283L^2}+O(L^{-3}).
\tag{4.16}
\] This is a limit of the explicitly computed coefficient ratio in
(4.11). Its fixed-box quantum remainder is still \(O_L(\xi^4)\) there;
(4.16) does not assert a uniform estimate for that remainder.

\subsection{5. Finite angles retain further nonlinear
channels}\label{magnetic-part-9-finite-angles-retain-further-nonlinear-channels}

The exact map (4.7) extends without a small-angle approximation: \[
 z_{\xi,\theta}=V_\xi^*\chi_{\xi,\theta}
                =V_\xi^*\mathcal C_\theta\psi_\xi,
 \qquad
 \chi_{\xi,\theta}=V_\xi z_{\xi,\theta}
                     +(I-P_\xi)\chi_{\xi,\theta}.
\tag{5.1}
\] Both terms are orthogonal to the actual vacuum. These exact
expressions, with (1.7)--(1.8), retain all higher electric
representations of that vacuum. In particular no replacement of
\(\psi_\xi\) by its free constant term is made in the map.

We can also exhibit every channel of its second coupling coefficient at
an arbitrary fixed angle. Set \[
 \alpha_p=\prod_{e\in\partial p,\ e=(n,1)}\cos(\theta n_2),
 \qquad \beta_p=\prod_{e\in\partial p,\ e=(n,1)}c_1(\theta n_2),
 \qquad u_p(\theta)=\frac{\alpha_p-1}{3}.
\tag{5.2}
\] Thus \(\alpha_{12}=\cos(n_2\theta)\cos((n_2+1)\theta)\),
\(\alpha_{13}=\cos^2(n_2\theta)\), and \(\alpha_{23}=1\). For adjacent
\(p,q\) sharing \(e\), let \[
 k_{pq,j}(\theta)=
 \prod_{f\in\partial p\mathbin\triangle\partial q,\ f=(n,1)}
                     \cos(\theta n_2)
 \times
 \begin{cases}c_j(\theta n_2(e)),&e\text{ is direction 1},\\1,&\text{otherwise},\end{cases}
 \qquad j=0,1.
\tag{5.3}
\] These are the exact multipliers of \(Z_{pq,j}\); (1.8) and their
joint spins prove this assertion. The multiplier of an edge-disjoint
product is \(\alpha_p\alpha_q\), and that of \(Y_p\) is \(\beta_p\).
Every product includes its original source coordinate.

The vacuum overlap has coefficient
\(c_{\xi,\theta}=1+\xi^2\sum_p(\alpha_p-1)/9+O_L(\xi^4)\), as follows by
inserting (3.3), using the mean of \(B^{\rm vac}\), and
\(\langle A,\mathcal C_\theta A\rangle=\sum_p\alpha_p/9\). Therefore \[
 z_{\xi,\theta}=\xi\sum_pu_p(\theta)W_p+O_{L,\theta}(\xi^3),
 \qquad
 (I-P_\xi)\chi_{\xi,\theta}
       =\xi^2\mathbf r_\theta+O_{H^k,L,\theta}(\xi^3),
\tag{5.4}
\] where the complete finite sum is \[
\begin{split}
 \mathbf r_\theta={}&
 \sum_p\left[\frac{\beta_p-1}{24}-\frac{u_p}{5}\right]Y_p\\
 &+\sum_{\{p,q\}\,\text{edge-disjoint}}
       \frac{(\alpha_p-1)(\alpha_q-1)}9 W_pW_q\\
 &+\sum_{\{p,q\}\,\text{adjacent}}
   \left[\frac4{27}(k_{pq,0}-1)-\frac23(u_p+u_q)\right]Z_{pq,0}\\
 &+\sum_{\{p,q\}\,\text{adjacent}}
   \left[\frac4{39}(k_{pq,1}-1)-\frac27(u_p+u_q)\right]Z_{pq,1}.
\end{split}
\tag{5.5}
\] To prove it, the second centered coefficient is
\((\mathcal C_\theta-I)B^{\rm vac}
 -[\sum_p(\alpha_p-1)/9]1\). Subtract the entire dressing
\(\mathscr D\sum_pu_pW_p\) from (3.4). Its constant equals the same
\(-\sum_p(\alpha_p-1)/9\) and cancels. The other coefficients are
exactly those displayed. In particular the edge-disjoint coefficient is
\((\alpha_p\alpha_q-1)/9-(\alpha_p+\alpha_q-2)/9\), proving its
displayed product, which need not vanish at a finite angle.

Finally (1.10) applied to each of these finitely many smooth channels,
or the explicit character Taylor expansion, proves \[
 \sum_pu_p(\theta)W_p=-\frac23\theta^2\mathbf a+O_{H^k,L}(\theta^4),
 \qquad \mathbf r_\theta=-\frac23\theta^2\mathbf r+O_{H^k,L}(\theta^4).
\tag{5.6}
\] The finite-angle edge-disjoint terms are of order \(\theta^4\) here;
their absence in (4.9) follows from their proved angular order. They are
present in the exact state and in (5.5).

\subsection{6. Full leading first-band spectral
measure}\label{magnetic-part-9-full-leading-first-band-spectral-measure}

For real \(0<\xi\le r_L/2\), define the exact first-band coefficient
matrix and the physical energy map \[
 C_\xi=\frac{B_\xi-e(\xi)I-3I}{\xi^2},\qquad
 t_\xi(\lambda)=\frac{\lambda/\kappa-3}{\xi^2},\qquad
 s_\xi(c)=\kappa(3+\xi^2c).
\tag{6.1}
\] These are mutually inverse affine maps between physical excitation
energy and its second coupling coefficient. Their \(\kappa\) is the same
coefficient in (1.2); no rescaled physical Hamiltonian is substituted.
Equation (3.9) gives the explicit norm estimate \[
 \|C_\xi-F_L\|\le\frac73r_L^{-4}\xi^2.
\tag{6.2}
\]

Let \(E^{Q_L}_q\) be the orthogonal projection onto the eigenspace of
the finite real symmetric \(Q_L\) with eigenvalue \(q\), and define the
exact finite measures \[
 \mu_L^a=\sum_{q\in\operatorname{spec}Q_L}
          \|E^{Q_L}_q(a_p)_p\|_{\ell^2}^2\,\delta_q,
 \qquad
 \eta_L^a=\left(q\mapsto\frac7{15}-\frac q{21}\right)_*\mu_L^a.
\tag{6.3}
\] Repeated eigenvalues are handled by the entire eigenspace projection,
so these definitions do not depend on choices of eigenvectors. Their
total mass is \(\mathsf N_L\), not 1, by completeness. Formula (3.7) and
the unitary \(\mathcal J\) prove the exact coefficient-observable
identity \[
 \int f(c)\,d\eta_L^a(c)
 =\langle\mathbf a,f(F_L)\mathbf a\rangle
 =\sum_q\|E^{Q_L}_q(a_p)_p\|^2f\!\left(\frac7{15}-\frac q{21}\right).
\tag{6.4}
\]

For the actual state define its physical first-band measure by \[
 \mu^{\rm band}_{\xi,\theta}(J)
 =\langle P_\xi\chi_{\xi,\theta},
             1_J(\mathscr A_\xi)P_\xi\chi_{\xi,\theta}\rangle
 \quad (J\text{ Borel in }\mathbb R).
\tag{6.5}
\] Equations (2.7) and (5.1) give its exact pushforward, at every angle,
\[
 (t_\xi)_*\mu^{\rm band}_{\xi,\theta}(J)
     =\langle z_{\xi,\theta},1_J(C_\xi)z_{\xi,\theta}\rangle,
 \qquad
 \mu^{\rm band}_{\xi,\theta}
     =(s_\xi)_*\mu_{C_\xi,z_{\xi,\theta}}.
\tag{6.6}
\] Here \(\mu_{C,z}(J)=\langle z,1_J(C)z\rangle\). This is the exact map
of both the actual vector and all its first-band observables.

For every bounded continuous function \(f:\mathbb R\to\mathbb C\), the
complete leading statement in the requested order of limits is \[
\boxed{
 \lim_{\xi\downarrow0}\frac1{\xi^2}
 \left[\lim_{\theta\to0}\frac9{4\theta^4}
  \int f\bigl(t_\xi(\lambda)\bigr)
             \,d\mu^{\rm band}_{\xi,\theta}(\lambda)\right]
 =\int f(c)\,d\eta_L^a(c).}
\tag{6.7}
\] Indeed (1.10) and boundedness of the relevant finite-dimensional
functional calculus give the inner limit as
\(\langle y_\xi,f(C_\xi)y_\xi\rangle\). Dividing by \(\xi^2\), the
vector tends to \(\mathbf a\) by (4.8), and \(f(C_\xi)\to f(F_L)\) in
matrix norm by (6.2). To justify the last assertion without a
differentiability assumption on \(f\), all spectra are eventually in one
compact interval. Approximate the restriction of the real and imaginary
parts of \(f\) uniformly by polynomials. Polynomial functional calculus
is continuous in the matrix entries; spectral calculus bounds each
uniform approximation error by its scalar supremum. First pass
\(\xi\to0\), then let that error tend to zero. This proves (6.7) for
precisely the stated bounded-continuous class. No estimate of order
\(\xi^2\) for an arbitrary continuous \(f\) is asserted.

Equivalently, keeping the actual norm as the denominator gives \[
 \lim_{\xi\downarrow0}\left[\lim_{\theta\to0}
 \frac{\displaystyle\int f(t_\xi(\lambda))
                     \,d\mu^{\rm band}_{\xi,\theta}(\lambda)}
      {\|\chi_{\xi,\theta}\|^2}\right]
       =\frac1{\mathsf N_L}\int f(c)\,d\eta_L^a(c).
\tag{6.8}
\] This is an explicitly defined probability comparison derived by
division of (6.7) and the independently retained mass formula; it does
not alter the original state or discard its amplitude.

The physical coefficient locations, with every original scale
substituted, are \[
 \kappa\left[3+\xi^2\left(\frac7{15}-\frac q{21}\right)\right]
 =\frac{6g_{\rm YM}^2}{a}
  +\frac1{8g_{\rm YM}^6a}\left(\frac7{15}-\frac q{21}\right).
\tag{6.9}
\] Equation (6.6) is the exact physical pushforward; (6.7) resolves its
leading coefficient measure at every \(q\), including all its weights.
Individual eigenvalues of the actual band differ from the ordered
coefficient predictions by at most \(\kappa(7/3)r_L^{-4}\xi^4\), by
(3.9) and min--max. Within a degenerate coefficient eigenspace higher
orders may split eigenvalues; the spectral measure formulation retains
their aggregate limiting weight without making an unsupported choice of
split eigenvectors.

For fixed \(\kappa>0\), without coefficient resolution one also gets \[
 \lim_{\xi\downarrow0}\frac1{\xi^2}
 \left[\lim_{\theta\to0}\frac9{4\theta^4}
            \int f(\lambda)\,d\mu^{\rm band}_{\xi,\theta}(\lambda)\right]
 =\mathsf N_L f(3\kappa).
\tag{6.10}
\] If \(\kappa\) varies with the original \(g_{\rm YM},a\) while
\(\xi\to0\), (6.6)--(6.9) remain the applicable exact pushforward and
coefficient statements; (6.10) specifically requires fixed \(\kappa\).

\subsection{7. Out-of-band order and its leading physical spectral
measure}\label{magnetic-part-9-out-of-band-order-and-its-leading-physical-spectral-measure}

Set \(\chi^{\rm out}_{\xi,\theta}=(I-P_\xi)\chi_{\xi,\theta}\). It is
smooth, orthogonal to the actual vacuum and to the entire actual first
band. At each fixed positive \(\xi\) in the band interval, (1.10) proves
\[
 \lim_{\theta\to0}\frac9{4\theta^4}
          \|\chi^{\rm out}_{\xi,\theta}\|^2=\|\rho_\xi\|^2,
 \qquad
 \lim_{\theta\to0}\frac9{4\theta^4}
          \|\chi_{\xi,\theta}\|^2=\|v_\xi\|^2.
\tag{7.1}
\] Therefore the precise nonzero leading out-of-band norm is \[
\boxed{
 \lim_{\xi\downarrow0}\frac1{\xi^4}
 \left[\lim_{\theta\to0}\frac9{4\theta^4}
               \|\chi^{\rm out}_{\xi,\theta}\|^2\right]=\mathsf R_L>0,}
\tag{7.2}
\] and its retained fraction is \[
 \lim_{\theta\to0}
   \frac{\|\chi^{\rm out}_{\xi,\theta}\|^2}{\|\chi_{\xi,\theta}\|^2}
  =\xi^2\frac{\mathsf R_L}{\mathsf N_L}+O_L(\xi^4).
\tag{7.3}
\] In unsquared norms the ordered leading amplitude is
\((2/3)\theta^2\xi^2\sqrt{\mathsf R_L}\); for the whole state it is
\((2/3)\theta^2\xi\sqrt{\mathsf N_L}\), with \(\xi>0\). More precisely,
first \(\lim_{\theta\to0}3\|\chi^{\rm out}_{\xi,\theta}\|/(2\theta^2)
=\|\rho_\xi\|\), and then use (4.11). These are ordered limits, so no
unstated uniformity in \(L,\theta,\xi\) is used.

Define \(\mu^{\rm out}_{\xi,\theta}\) by (6.5) with \(P_\xi\chi\)
replaced by \(\chi^{\rm out}\). For fixed \(\kappa\) and every bounded
continuous \(f\), its full leading measure is \[
\boxed{
 \lim_{\xi\downarrow0}\frac1{\xi^4}
 \left[\lim_{\theta\to0}\frac9{4\theta^4}
     \int f(\lambda)\,d\mu^{\rm out}_{\xi,\theta}(\lambda)\right]
 =\mathsf R_{8,L}f(8\kappa)
  +\mathsf R_{9/2,L}f(9\kappa/2)
  +\mathsf R_{13/2,L}f(13\kappa/2).}
\tag{7.4}
\] If physical \(\kappa\) varies, replace the test in the integrand by
\(f(\lambda/\kappa)\); the right side then has arguments \(8,9/2,13/2\).
This is its exact dimensionless-to-physical pushforward.

Here is a proof of the functional-calculus limit, including its
high-energy tail. Put \(u_\xi=\rho_\xi/\xi^2\). By (4.8),
\(u_\xi\to\mathbf r\) in \(H^2\), and
\((K_\xi-e(\xi))u_\xi\to H_0\mathbf r\) in Hilbert norm. Thus these
vectors have uniformly bounded norm and uniformly bounded second
spectral moment. For a real spectral cutoff \(R>0\), the spectral mass
at \(|\lambda|>R\) of either vector is bounded by its second moment
divided by \(R^2\). Also \(K_\xi-e(\xi)\) converges in norm resolvent to
\(H_0\): their difference \(-\xi S-e(\xi)I\) is bounded with norm
tending to zero, and the resolvent identity proves convergence at a
nonreal spectral parameter. The algebra generated by resolvents and
their conjugates is uniformly dense in \(C_0(\mathbb R)\): it separates
real points, is closed under conjugation, and is nonzero at each real
point. Uniform approximation therefore proves norm convergence of
\(f_0(K_\xi-e(\xi))\) for every \(f_0\in C_0(\mathbb R)\). For a general
bounded continuous \(f\), multiply by a continuous compactly supported
cutoff equal to 1 on \([-R,R]\). The discarded expectation is at most
\(\|f\|_\infty\) times the just-proved tail mass. First use
\(u_\xi\to\mathbf r\) and the \(C_0\) convergence, and then let
\(R\to\infty\). This proves convergence for the claimed class. Finally
(4.9), orthogonality, and (3.2) give exactly the three masses and
arguments in (7.4). The angular inner limit follows independently from
(1.10) and boundedness of \(f(\mathscr A_\xi)\).

Since the exact band projection commutes with the full Hamiltonian, the
full state measure is the sum
\(\mu^{\rm band}_{\xi,\theta}+\mu^{\rm out}_{\xi,\theta}\), with no
cross term for any Borel set. Equations (6.7) and (7.2) show that its
first-band coefficient limit also holds when the full state measure is
used on the left of (6.7): the absolute contribution of the omitted
component, after the displayed angular and \(\xi^2\) divisions, is
bounded by \(\|f\|_\infty\|\rho_\xi\|^2/\xi^2=O_L(\xi^2)\). This remains
true although \(t_\xi\) sends its higher energies to large coefficient
values; boundedness of the test function is the precise hypothesis
controlling that contribution.

There is also a direct finite-box two-parameter norm estimate, without
taking a volume limit. On a sufficiently small rectangle of real
\((\xi,\theta)\) around zero, \[
\begin{split}
 z_{\xi,\theta}
   &=-\frac23\xi\theta^2\mathbf a
      +O_L(|\xi|^3\theta^2+|\xi|\theta^4),\\
 \chi^{\rm out}_{\xi,\theta}
   &=-\frac23\xi^2\theta^2\mathbf r
      +O_{H^k,L}(|\xi|^3\theta^2+\xi^2\theta^4).
\end{split}
\tag{7.5}
\] In particular its out-of-band norm is \(O_L(\xi^2\theta^2)\). Here is
the joint remainder justification. By (2.8), the vacuum and its
\(\xi\)-derivatives are bounded in each needed higher Sobolev space on a
smaller closed coupling interval. Averaged translation has bounded
derivatives through order four from \(H^{k+4}\) to \(H^k\), uniformly on
a closed angle interval. The expressions in (5.1) therefore have the
corresponding bounded mixed derivatives. They are even in \(\theta\) and
identically zero at \(\theta=0\). The out-of-band expression and its
first \(\xi\)-derivative vanish at \(\xi=0\) for every angle, by (5.4).
Taylor's integral formula in \(\xi\) factors out \(\xi^2\); Taylor's
formula in \(\theta\) then factors out \(\theta^2\), with a fourth-order
angular remainder. The remaining coefficient at \(\xi=0\) is (5.5),
whose second angular coefficient is (5.6). This proves the second line,
including its mixed remainder factors. The transported coordinate is odd
in \(\xi\), by (3.6), so the same argument factors its coupling
remainder as \(\xi^3\) and proves the first line. All constants here
still depend on the fixed original box; (7.5) does not assert a uniform
volume estimate.

The physical amplitude factors in (6.7) and (7.2) are respectively \[
 \frac49\theta^4\xi^2=\frac{\theta^4}{36g_{\rm YM}^8},
 \qquad
 \frac49\theta^4\xi^4=\frac{\theta^4}{576g_{\rm YM}^{16}},
 \qquad \theta^4=\frac{16\pi^4a^8}{D^4}.
\tag{7.6}
\] These multiply \(\mathsf N_L\) and \(\mathsf R_L\) in the respective
ordered leading squared norms. None of those factors is absorbed into
the graph vector or the physical energy scale.

\subsection{8. Scope of the coefficient
map}\label{magnetic-part-9-scope-of-the-coefficient-map}

Equations (1.7), (2.5), (4.7), (5.1), and (6.6) give exact maps from the
specified continuous magnetic background to the actual physical state,
its first interacting spectral subspace, and all its spectral
observables. The full leading coefficient measure is precisely the
pushforward (6.3), with the original polynomial vector (4.2). The
out-of-band contribution has the explicitly positive coefficient (4.10)
and the three leading physical energies (7.4); all finite-angle
nonlinear channels are retained in (5.5) and in the exact class operator
(1.7).

Every asymptotic assertion above fixes \(L\) first. The sufficient
spectral interval is \(0<\xi\le3/(128M)\), and the explicit matrix
remainder constant grows as \(M^4\). No joint continuum convergence,
fixed-positive-coupling infinite-volume convergence, or exchange of
these limits is used. A separate limit of the completely specified
finite graph measure \(\mu_L^a\) can be pushed through the affine map in
(6.3); the exact quantum statement established here is the displayed
angular-first, coupling-second limit at every fixed original box.

The companion exact checker
\url{check_native_magnetic_physical_band_transport.py} constructs the
original link and face incidence sets, checks the plane weights and
polynomial mass, independently divides the vacuum channel coefficients
by their electric energies, subtracts every spectral dressing
denominator, and verifies both shared-edge coefficients and their
angular factors using rational arithmetic. It also evaluates all
finite-angle second-order channels at \(\theta=\pi/3\) for \(L=2\),
including the nonzero edge-disjoint terms. The analytic and spectral
arguments are the proofs above; these arithmetic checks supplement them.

\clearpage
% Complete next live proof; its original equation tags are retained.
\section{The first separated-link connected electric
covariance}\label{magnetic-part-10-the-first-separated-link-connected-electric-covariance}

8 September 2026. Two links which share no elementary face nevertheless
develop a nonzero connected electric covariance through two adjacent
faces. This note calculates that coefficient in the actual full-box
vacuum. The calculation keeps both intermediate spin channels and the
disconnected expectation subtraction. It does not assume a
volume-uniform remainder.

\subsection{1. Original Hilbert space and the coefficient to be
calculated}\label{magnetic-part-10-original-hilbert-space-and-the-coefficient-to-be-calculated}

Fix \(L\geq2\), vertices \(\{-L,\ldots,L\}^3\), all positive contained
links e, and all contained elementary faces p.~Write their complete
counts as \(N=3(2L)(2L+1)^2\) and \(M=3(2L)^2(2L+1)\). Link variables
are \(U_e\in SU(2)\), with the inverse on reverse traversal. The face
trace is \[
 W_{(n;i,j)}=\operatorname{tr}\{U_i(n)U_j(n+\mathbf e_i)
           U_i(n+\mathbf e_j)^{-1}U_j(n)^{-1}\},\qquad i<j.
\] On product Haar probability space use \(T_a=-i\sigma_a/2\), left
derivatives \(X_{e,a}\), \(E_e=-\sum_aX_{e,a}^2\), and \[
 H=\kappa\sum_eE_e+b\sum_p(2-W_p),\quad
 \kappa=2g^2/a,\quad b=1/(2g^2a),\quad\xi=b/\kappa=1/(4g^4).
\tag{1}
\] The fundamental Casimir is \(3/4\). In the original metric
\(-\operatorname{tr}(AB)/2\), the Laplacian is four times the
sum-of-squares Laplacian in these generators, so the electric
coefficient in that metric is \(\kappa/4\). Neither this factor nor the
Wilson scalar \(2bM\) is changed.

The compact connected configuration manifold and bounded smooth
potential give a self-adjoint compact-resolvent operator of domain
\(H^2\), form domain \(H^1\), and unique smooth positive unit ground
vector \(\psi_\xi\). For precision, the modulus inequality supplies a
nonnegative form minimizer, elliptic regularity and the maximum
principle make it positive, and the identity
\(q_{H-\mathcal E}(\psi f)=\kappa\sum_{e,a}\int\psi^2|X_{e,a}f|^2\)
implies uniqueness by division by \(\psi\). Gauge transformations
preserve this vector and every link Casimir. All following vectors are
thus physical.

Set \(r_e=\#\{p:e\in\partial p\}\),
\(\gamma_e=\langle\psi,E_e\psi\rangle\), \(v_e=(E_e-\gamma_e)\psi\), and
\[
 C_{ef}=\langle v_e,v_f\rangle
   =\langle\psi,E_eE_f\psi\rangle-\gamma_e\gamma_f.
\tag{2}
\] The operators commute for different links. Smoothness makes (2)
legitimate as an operator identity on the vacuum, not only a formal
product of forms.

For distinct e,f that share no face, define the integer \[
 a_{ef}=\#\{(p,q): e\in\partial p,\ f\in\partial q,
                     |\partial p\cap\partial q|=1\}.
\tag{3}
\] The pair in (3) is ordered by which face contains e and which
contains f; there is no double counting because no face contains both
links. All faces are those of the original open box. The result is \[
 \boxed{C_{ef}=\frac{127}{219024}\,a_{ef}\xi^4+O_L(\xi^6)
        \quad(e,f\text{ share no face}).}
\tag{4}
\] In particular the coefficient is positive whenever a two-face channel
exists, even though the energy matrix between these links is exactly
zero. The latter assertion is proved in Section 4.

\subsection{2. The full second vacuum coefficient, including both
shared-link
channels}\label{magnetic-part-10-the-full-second-vacuum-coefficient-including-both-shared-link-channels}

Put \(H_0=\sum E_e\) and \(S=\sum W_p\), so \(H=\kappa(H_0-\xi S)+2bM\).
On \(|z|=3/8\) in the full Hilbert space the resolvent of H\_0 has norm
at most \(8/3\), since its first nonzero eigenvalue is \(3/4\). Also
\(\|S\|=2M\). Therefore the bounded perturbation Neumann series
converges for \(|\xi|<3/(16M)\), gives the analytic rank-one bottom
projection, and gives an analytic positive unit eigenvector for real
\(\xi\) near zero. Elliptic regularity applied repeatedly to the
eigen-equation upgrades this to analyticity in each fixed Sobolev space.
Its radius and estimates may depend on L. This suffices for the finite
derivative calculations below.

Product Haar integration gives \(\int W_p=0\),
\(\langle W_p,W_q\rangle=\delta_{pq}\), and \(H_0W_p=3W_p\). Indeed a
link in only one of two distinct face boundaries has odd central parity;
for the same face integrate its Haar holonomy and use the fundamental
character squared norm one. Comparing coefficients in the actual
equation \((H_0-\xi S)\psi=e(\xi)\psi\), including its unit-vector
condition, gives \[
 \psi=1+\xi A_1+\xi^2A_2+O_{H^k,L}(\xi^3),\quad A_1=S/3,
 \quad e(\xi)=-M\xi^2/3+O_L(\xi^3),
\] \[
 H_0A_2=(S^2-M)/3,\qquad\int A_2=-M/18.
\tag{5}
\] The inverse in (5) can be resolved entirely. For adjacent distinct
faces p,q sharing g put \[
 Z_{pq,0}=\int W_pW_q\,dU_g,\qquad
 Z_{pq,1}=W_pW_q-Z_{pq,0}.
\] After reversing one real trace if needed, their words are
\(\operatorname{tr}(U_gB)\), \(\operatorname{tr}(U_g^{-1}D)\), with six
distinct outer links. The fundamental Haar identity gives
\(Z_{pq,0}=\operatorname{tr}(BD)/2\). The shared-edge spin tensor
\(1/2\otimes1/2\) consists of spins zero and one, while all outer links
have spin \(1/2\). Thus \[
 H_0Z_{pq,0}=\tfrac92 Z_{pq,0},\quad
 H_0Z_{pq,1}=\tfrac{13}2Z_{pq,1},\quad
 \|Z_{pq,0}\|^2=\tfrac14,\quad\|Z_{pq,1}\|^2=\tfrac34.
\tag{6}
\] The norms use orthogonal Haar projection on g and total product norm
one. A repeated face has \(W_p^2-1=\chi_1(U_p)\) with free energy8 and
norm one. An edge-disjoint pair has energy6 and norm one, even when the
faces meet at a vertex. Consequently the full coefficient is \[
 \boxed{A_2=-\frac M{18}+\sum_p\frac{W_p^2-1}{24}
 +\sum_{\{p,q\}\ {\rm edge\!\!-\!\!disjoint}}\frac{W_pW_q}{9}
 +\sum_{\{p,q\}\ {\rm adjacent}}
       \left(\frac4{27}Z_{pq,0}+\frac4{39}Z_{pq,1}\right).}
\tag{7}
\] All distinct-pair sums in (7) are unordered; the factor two in S² has
already been included. No finite-face subsystem is substituted for H.

Different unordered pairs in (7) have different odd-link parity
supports. If their symmetric boundaries agreed, their symmetric face
difference would be a nonempty closed mod-two collection of at most four
squares. A face in that collection needs another face on each of its
four boundary edges; two distinct elementary squares share at most one
edge, so at least four other squares would be needed, a contradiction.
Haar integration on a link where parity differs makes the cross term
zero. Repeated-face terms are even everywhere, and different
repeated-face terms have different joint link spins. Thus all the
displayed components, with their two spin channels, are mutually
orthogonal. This remains true after application of any link Casimir,
which preserves those parity and spin subspaces.

\subsection{3. Exact connected fourth
coefficient}\label{magnetic-part-10-exact-connected-fourth-coefficient}

Suppose e,f share no face. Then \(E_eE_f1=E_eE_fA_1=0\): each term in
A\_1 is a single face and depends on at most one of these links.
Self-adjointness moves \(E_eE_f\) to either of these vectors, so every
term of degree less than four in \(\langle\psi,E_eE_f\psi\rangle\)
vanishes, and its degree-four term is exactly \[
 \langle A_2,E_eE_fA_2\rangle.
\tag{8}
\] In particular no third or fourth unknown coefficient of the vacuum is
being omitted from (8): its possible pairing with 1 or A\_1 is
identically zero for the stated reason. Similarly \[
 \gamma_e=\xi^2\langle A_1,E_eA_1\rangle+O_L(\xi^3)
          =\frac{r_e}{12}\xi^2+O_L(\xi^3).
\tag{9}
\]

Only a pair p,q with \(e\in p\), \(f\in q\) contributes to (8). Repeated
faces contribute zero. There are exactly \(r_er_f\) such distinct
unordered pairs when they are labeled as in (3). For an edge-disjoint
pair, both e and f carry fundamental spin, so its contribution is
\((3/4)^2/9^2=1/144\). For an adjacent pair, neither e nor f can be its
shared edge: otherwise the face containing the other link would contain
both e and f.~They are therefore both outer fundamental edges in each of
the two channels (6), and the contribution is \[
 \left(\frac34\right)^2
 \left[\left(\frac4{27}\right)^2\frac14
       +\left(\frac4{39}\right)^2\frac34\right]
       =\frac1{324}+\frac3{676}.
\tag{10}
\] Equation (9) subtracts \(r_er_f/144\) at degree four in (2). Hence
every edge-disjoint pair cancels its disconnected contribution exactly;
every adjacent pair leaves \[
 \frac1{324}+\frac3{676}-\frac1{144}
          =\frac{127}{219024}>0.
\tag{11}
\] This proves the coefficient in (4).

For its stated even error, the center transformation
\(U_{(n,i)}\mapsto(-1)^{\sum_{j<i}n_j}U_{(n,i)}\) changes every face
trace sign, commutes with every E\_e and preserves Haar measure. It
intertwines H\_0−ξS with H\_0+ξS and hence the positive unit vacua at ξ
and−ξ. Both terms in (2) are consequently even analytic functions of ξ
near zero. Their next possible order after four is six. This proves the
error in (4); it does not supply a volume-uniform analytic radius.

\subsection{4. Exact support map and relation to the energy
matrix}\label{magnetic-part-10-exact-support-map-and-relation-to-the-energy-matrix}

Define the edge adjacency matrix \(B\) by \(B_{eg}=1\) when e,g are
distinct physical links belonging to a common elementary face, and zero
otherwise. Two distinct links belong to at most one elementary square:
perpendicular incident links determine its two plane directions and
base, while parallel links in a square must be opposite sides and also
determine that square. Thus, for the e,f considered above, \[
 \boxed{a_{ef}=(B^2)_{ef}.}
\tag{12}
\] Indeed a pair p,q counted in (3) has a unique shared edge g, which is
a common neighbor of e,f.~Conversely a common neighbor g determines the
unique p containing e,g and q containing g,f; these are distinct because
e,f share no face, and their unique shared edge is g. These two maps are
inverse. The covariance coefficient thus has support at graph distance
two, including boundary incidences, and no support beyond distance two
at this order. It makes no all-orders finite-range claim about C.

There are actual such channels in every L≥2. Take \(e=((0,0,0),1)\) and
\(f=((0,2,0),1)\). They share no face, but the 12 faces based at
\((0,0,0)\) and \((0,1,0)\) share the edge \(((0,1,0),1)\), giving
\(a_{ef}\geq1\). All of their vertices are in the original box,
including at L=2.

Put \(\mathcal N_{ef}=\langle v_e,(H-\mathcal E)v_f\rangle\). Commuting
the link operators on the smooth real vacuum gives \[
 \mathcal N_{ef}=\tfrac12
 \langle\psi,[E_e,[H,E_f]]\psi\rangle.
\tag{13}
\] To check (13), expand the double commutator, use
\(H\psi=\mathcal E\psi\), commute E\_e,E\_f, and use the real symmetry
of their two matrix elements. Only potential faces containing f occur in
\([H,E_f]\). If e shares no face with f, E\_e commutes with every one of
their multiplication and derivative terms. Thus \[
 \boxed{\mathcal N_{ef}=0\text{ exactly, at every positive coupling,
 when e,f share no face}.}
\tag{14}
\] There is no contradiction between (4) and (14): a matrix-valued
spectral measure can have nonzero zeroth moment and zero first moment
off diagonal. Explicitly, for the spectral resolution of A=H−E on the
physical space, \[
 d\Sigma_{ef}(\omega)=\langle v_e,\mathbf1_{d\omega}(A)v_f\rangle,
 \quad C_{ef}=\int d\Sigma_{ef},\quad
 \mathcal N_{ef}=\int\omega\,d\Sigma_{ef}.
\tag{15}
\] The measure matrix is positive semidefinite on every Borel set
because its quadratic form is the squared norm of a projected linear
combination of the v\_e. An individual off-diagonal entry, however, need
not be a nonnegative measure; (15) retains its signs. Smoothness gives
the finite moments used here. These identities are the exact relation
between the nonlocal covariance contribution and the local energy
matrix, not an assumption that one determines the other's support.

Finally \(\xi^4=1/(256g^{16})\). The physical coefficient in (4) is
therefore \(127a_{ef}/(56070144g^{16})\); C has no extra energy factor,
whereas \(\mathcal N\) has one. The scalar energy shift remains in H and
cancels only in H−E. The fixed-box coefficient is a concrete next input
to the extensive denominator calculation; control of all higher orders
and summability at fixed coupling is not inferred from its finite range
at order four.

\subsection{5. The complete fourth-order spectral functional and time
correlation}\label{magnetic-part-10-the-complete-fourth-order-spectral-functional-and-time-correlation}

The same calculation can resolve the energy dependence, rather than only
its zeroth moment. Put \(A_\xi=(H-\mathcal E)/\kappa\), and write \[
 F_{ef,\xi}(h)=\langle v_e,h(A_\xi)v_f\rangle.
\] For each fixed box, for distinct links sharing no face, and every
bounded continuous real or complex function h which is differentiable at
3, \[
\boxed{\lim_{\xi\to0}\xi^{-4}F_{ef,\xi}(h)
 =a_{ef}\left[-\frac{59}{275184}h(3)-\frac1{336}h'(3)
             +\frac1{1296}h(9/2)+\frac3{132496}h(13/2)\right].}
\tag{16}
\] Here the limit is a functional on the specified test functions. In
particular the derivative term must not be omitted or represented as an
ordinary positive spectral atom. Each original matrix-valued spectral
measure is still the exact positive-semidefinite measure (15).

We prove (16), including the relation to the actual first interacting
spectral subspace. Let P project onto the free physical energy-three
space \(\mathcal P=\operatorname{span}\{W_p\}\), and let \[
 R_3=(I-P)((H_0-3)|_{\mathcal P^\perp})^{-1}(I-P).
\] The inverse is on the physical space: its constant channel is -1/3.
To verify isolation, a nonzero physical spin network has no degree-one
support vertex. Its support contains a cycle of length at least four in
the bipartite cubic graph, costing at least 3. With fewer than six
nontrivial edges it must be one elementary square: a fifth edge would
have a new degree-one endpoint, or be a forbidden same-bipartition chord
of the four-cycle. The degree-two invariant contractions require equal
spins on that square. Only spin 1/2 has energy below 9/2. A six-edge
rectangular loop realizes energy 9/2. Thus the first physical free band
is exactly P, with the next energy 9/2.

For sufficiently small real \(\xi\), the contour \(|z-3|=3/4\) defines
the actual spectral projection \(P_\xi\), and the exact isometry \[
 V_\xi=P_\xi P(PP_\xi P|_{\mathcal P})^{-1/2}
       :\mathcal P\longrightarrow\operatorname{ran}P_\xi
\tag{17}
\] uses its full Gram factor. Indeed the free contour resolvent has norm
at most 4/3. For \(|\xi|\le3/(64M)\), its Neumann factor is at most 1/8,
and integrating the resolvent difference gives \(\|P_\xi-P\|\le1/7\).
The displayed inverse square root therefore exists by its binomial
series; idempotence and orthogonality of the real projection show
\(V_\xi^*V_\xi=I\). Its range is smooth, and the finite-rank resolvent
formula and elliptic regularity justify all operator products and fixed
Sobolev derivatives below. It is orthogonal to the actual vacuum for a
smaller interval, for example \(|\xi|\le3/(128M)\), by the free gaps and
the bounded perturbation estimate \(2M|\xi|\).

Let Q be the unchanged open-face signless adjacency matrix, with
diagonal the actual face degrees and off-diagonal 1 for shared-edge
faces. Direct virtual-channel calculation gives \[
 V_\xi^*A_\xi V_\xi=3I+\xi^2F+O_L(\xi^4),\qquad
 F=\frac7{15}I-\frac1{21}Q.
\tag{18}
\] Here is the coefficient calculation rather than an assumption of
(18). Differentiating (17) gives \(V'_0=R_3SP\). The second matrix
coefficient before vacuum subtraction is \(-PSR_3SP\). The vacuum
channel gives +1/3 to every entry; a repeated-face channel gives -1/5 to
its diagonal; an edge-disjoint pair gives -1/3 to its two diagonals and
mutual entries; an adjacent pair gives
\(- (1/4)/(9/2-3)-(3/4)/(13/2-3)=-8/21\) to that same four-entry
pattern. Thus the diagonal is \(7/15-M/3-d_p/21\), and the off-diagonal
is -1/21 exactly for adjacent faces. The actual vacuum coefficient -M/3
from (5) gives (18). The central cochain used in Section 3 acts as -I on
P; it makes the matrix in (18) even and accounts for the fourth-order
remainder.

The electric vector can be treated with arbitrary fixed real link
weights w, retaining its linear dependence on all of them. Define
\(\Gamma_w=\sum w_eE_e\), \(a_p(w)=\sum_{e\in p}w_e/4\), and
\(v(w)=(\Gamma_w-\langle\Gamma_w\rangle)\psi\). Put \[
 y_\xi(w)=V_\xi^*v(w),\qquad
 \rho_\xi(w)=(I-P_\xi)v(w).
\] These are exact orthogonal components, both orthogonal to the actual
vacuum. Coefficient comparison using (7) and \(V'_0=R_3SP\) gives \[
 y_\xi(w)=\xi x(w)+\xi^3z(w)+O_L(\xi^5),\qquad
 x(w)=\sum_pa_p(w)W_p,
\] \[
 \rho_\xi(w)=\xi^2r(w)+O_{H^k,L}(\xi^3),
\tag{19}
\] where, for \(s=a_p(w)+a_q(w)\) and g the shared edge, \[
 r(w)=\sum_p\frac{2a_p(w)}{15}(W_p^2-1)
       +\sum_{\{p,q\}\ {\rm adjacent}}
       \left[-\frac{2(s+w_g)}9Z_{pq,0}
             +\frac{2(3s+7w_g)}{273}Z_{pq,1}\right].
\tag{20}
\] For completeness, \(\Gamma_w\) has eigenvalues \(8a_p\) on the
repeated face, \(3s\) on an edge-disjoint pair, and
\(3s-3w_g/2,3s+w_g/2\) on the adjacent zero and one channels. The shared
edge is counted twice as fundamental in 3s; replacing that by its actual
spin accounts for precisely the last two corrections. In
\(\Gamma_w A_2\), the respective coefficients are
\(a_p/3,s/3,(4s-2w_g)/9,(12s+2w_g)/39\). Subtract the full dressing
\(R_3Sx(w)\), whose coefficients are \(a_p/5,s/3,2s/3,2s/7\). The two
constant coefficients are both \(-\sum_pa_p/3\) and also cancel. This
proves (20) with no omitted channel. The centered projection identity
\(V_\xi^*\psi_\xi=0\) is exact. The same central cochain makes
\(y_\xi(w)\) odd, proving the first expansion (19).

For a single link e use the original weight \(w^{(e)}_g=\delta_{eg}\),
and denote the resulting x,z,r by \(x_e,z_e,r_e^{\rm vec}\). This
superscript distinguishes the vector from the incidence count r\_e. Then
\(x_e=\sum_{p\ni e}W_p/4\). If e,f share no face,
\(\langle x_e,x_f\rangle=0\), and (18) gives exactly \[
 \langle x_e,Fx_f\rangle=-a_{ef}/336.
\tag{21}
\] Each ordered adjacent pair contributes \(-1/(4\cdot21\cdot4)\). In
the cross spectral measure of \(r_e^{\rm vec},r_f^{\rm vec}\), only the
pairs counted by (3) occur. Repeated-face cross terms vanish. For every
contributing pair, e and f are outer links of different faces, \(s=1/4\)
for each weight, and \(w_g=0\). Thus their zero-channel coefficients in
(20) are both -1/18, and their one-channel coefficients are both 1/182.
Orthogonality and (6) give the complete free cross measure \[
 \langle r_e^{\rm vec},h(H_0)r_f^{\rm vec}\rangle
 =a_{ef}\left[\frac1{1296}h(9/2)+\frac3{132496}h(13/2)\right].
\tag{22}
\] No edge-disjoint channel survives in (20).

As \(A_\xi-H_0=-\xi S-e(\xi)I\) tends to zero in operator norm, its
resolvents converge in norm. For bounded continuous h, this implies
convergence of the cross expectations on \(\rho_\xi/\xi^2\) to (22):
first use continuous functions vanishing at infinity and norm resolvent
convergence, then approximate h on a compact interval and control the
discarded tail by the uniformly bounded second spectral moments. Those
moments follow from the \(H^2\) convergence in (19). Polarization gives
the cross statement from the four corresponding diagonal statements.
Hence (22) is the entire fourth-order out-of-band contribution.

The fourth coefficient of the band mass is obtained without knowing the
third vacuum coefficient: subtract (22) at h=1 from (4). It is \[
 \langle x_e,z_f\rangle+\langle z_e,x_f\rangle
 =a_{ef}\left(\frac{127}{219024}-\frac1{1296}-\frac3{132496}\right)
 =-\frac{59}{275184}a_{ef}.
\tag{23}
\] Differentiability of h at 3 and spectral calculus of the finite
matrix (18) give \(h(V_\xi^*A_\xi V_\xi)=h(3)I+\xi^2h'(3)F+o_L(\xi^2)\).
Combine this with (19), (21), (23), and the exact orthogonal spectral
decomposition into the band and its complement. This proves (16).

In particular the complete dimensionful ground-state time correlation,
with physical \(\kappa>0\) and physical time \(t\ge0\) fixed, is \[
\boxed{\begin{split}
 \langle v_e,e^{-t(H-\mathcal E)}v_f\rangle
 =a_{ef}\xi^4\biggl[&-\frac{59}{275184}e^{-3\kappa t}
       +\frac{\kappa t}{336}e^{-3\kappa t}\\
       &+\frac1{1296}e^{-9\kappa t/2}
       +\frac3{132496}e^{-13\kappa t/2}\biggr]+O_{L,t,\kappa}(\xi^6).
\end{split}}
\tag{24}
\] Here \(h(q)=e^{-\kappa tq}\) is used on the nonnegative excitation
spectrum, with any bounded continuous extension below zero, so it
belongs to the test class of (16). The bounded potential perturbation
gives \(e^{-\kappa tK_\xi}\) its norm-convergent Duhamel series: each
order is bounded by \((2M|\xi|\kappa t)^n/n!\). The excitation semigroup
retains the additional analytic scalar \(e^{\kappa t e(\xi)}\). The
vacuum and vectors are analytic in fixed Sobolev norms, and the central
cochain makes the whole expression even. These facts justify the even
remainder in (24), not only a pointwise little-o limit. For varying
physical \(\kappa\), the same proof uses the explicit parameter
\(\tau=\kappa t\) and keeps its dependence; no estimate uniform in
\(\tau\) or L is implied.

At t=0 the bracket is \(127/219024\). Its derivative at t=0 is exactly
zero, because \[
3\left(-\frac{59}{275184}\right)-\frac1{336}
 +\frac{9}{2\cdot1296}+\frac{39}{2\cdot132496}=0.
\tag{25}
\] This agrees with the exact identity (14). It exhibits which
separated-link spectral contributions cancel in the first moment: the
actual band-mass correction, the actual band-energy derivative, and both
higher spin channels. The term proportional to t is essential to that
cancellation. Equations (16)--(25) describe the fixed-box fourth
coefficient of the original interacting correlations; they do not
exchange a volume limit with the coupling expansion or assert a
continuum spectrum.

\subsection{6. Primary context}\label{magnetic-part-10-primary-context}

The original Hamiltonian framework is J. Kogut and L. Susskind,
\emph{Hamiltonian formulation of Wilson's lattice gauge theories},
Physical Review D \textbf{11} (1975), 395--408,
\url{https://doi.org/10.1103/PhysRevD.11.395}. Full conventions appear
in (1). For historical linked-cluster calculations in the same spatial
dimension, see A. C. Irving, T. E. Preece and C. J. Hamer, \emph{Cluster
expansion approach to non-abelian lattice gauge theory in (3+1)D (I).
SU(2)}, Nuclear Physics B \textbf{270} (1986), 536--552,
\url{https://doi.org/10.1016/0550-3213(86)90567-5}. This note derives
(4) directly, with no historical novelty assertion or imported
coefficient.

\clearpage
% Complete next live proof; its original equation tags are retained.
\section{Exact open-boundary polynomial for native magnetic spectral
leakage}\label{magnetic-part-11-exact-open-boundary-polynomial-for-native-magnetic-spectral-leakage}

8 September 2026. The coefficient calculated here is the first nonzero
out-of-band spectral mass of the original native magnetic state. Every
positive physical link, every elementary face, all coplanar and hinged
shared-edge pairs, and both shared-edge spin channels are retained. The
polynomial is derived from their original coordinates and boundary
counts. The physical state is the one built from the actual interacting
vacuum.

\subsection{1. Exact object whose coefficient is
counted}\label{magnetic-part-11-exact-object-whose-coefficient-is-counted}

Fix \(L\ge2\), with vertices \(\{-L,\ldots,L\}^3\). Include every
positive nearest-neighbor link \(e=(n,i)\) and elementary face
\(p=(n;i,j)\), \(i<j\), contained in that box. Denote the sets by
\(\mathsf E_L,\mathsf P_L\), with cardinalities \[
 N_L^{\rm edges}=3(2L)(2L+1)^2,\qquad M_L=3(2L)^2(2L+1).
\tag{1.1}
\] Each face has its original word \[
 W_{n;i,j}(U)=\operatorname{tr}\bigl[
 U_i(n)U_j(n+\mathbf e_i)U_i(n+\mathbf e_j)^{-1}U_j(n)^{-1}\bigr].
\tag{1.2}
\] There is one \(SU(2)\) variable per physical link; inverse traversal
uses the inverse of that same variable. Use product Haar probability
measure, the gauge-invariant closed subspace \(\mathcal H_{\rm phys}\),
\(T_a=-i\sigma_a/2\), the left derivatives \(X_{e,a}\) generated by
\(T_a\), and the exact link Casimirs \(E_e=-\sum_aX_{e,a}^2\). Retain \[
\begin{gathered}
 H=\kappa K_\xi+2bM_LI,
 \quad K_\xi=H_0-\xi S,
 \quad H_0=\sum_eE_e,
 \quad S=\sum_pW_p,\\
 \kappa=\frac{2g_{\rm YM}^2}{a},\qquad
 b=\frac1{2g_{\rm YM}^2a},\qquad
 \xi=\frac b\kappa=\frac1{4g_{\rm YM}^4},\qquad a>0.
\end{gathered}
\tag{1.3}
\] The original metric \(c(A,B)=-\operatorname{tr}(AB)/2\) has
\(-\Delta_e^c=4E_e\), included in \(\kappa\). The operator and form
domains of \(H\) are \(H^2\cap\mathcal H_{\rm phys}\) and
\(H^1\cap\mathcal H_{\rm phys}\), respectively: the product Laplacian is
elliptic on the compact configuration manifold and \(S\) is bounded
smooth multiplication. Its positive unit vacuum \(\psi_\xi\) is smooth,
with \(K_\xi\psi_\xi=e(\xi)\psi_\xi\). Positivity and uniqueness follow
from the modulus inequality, elliptic maximum principle, and the
identity \[
 \langle\psi_\xi F,(H-\mathcal E_\xi)\psi_\xi F\rangle
 =\kappa\sum_{e,a}\int\psi_\xi^2|X_{e,a}F|^2dU,
 \qquad \mathcal E_\xi=\kappa e(\xi)+2bM_L.
\tag{1.4}
\] A second ground vector divided by the strictly positive \(\psi_\xi\)
would have every derivative zero, hence be constant. Gauge
transformations therefore fix this same actual vacuum.

The native original weights, face averages, and angle are \[
\begin{gathered}
 w_{(n,1)}=n_2^2,\qquad w_{(n,2)}=w_{(n,3)}=0,
 \qquad a_p=\frac14\sum_{e\in\partial p}w_e,\\
 a_{n;1,2}=\frac{n_2^2+(n_2+1)^2}{4},\qquad
 a_{n;1,3}=\frac{n_2^2}{2},\qquad a_{n;2,3}=0,\\
 \Gamma=\sum_ew_eE_e,\qquad
 h_1(n)=\exp(-\theta n_2H_c),\quad h_2=h_3=I,
 \qquad H_c=i\sigma_3=-2T_3,\qquad \theta=\frac{2\pi a^2}{D},\quad D>0.
\end{gathered}
\tag{1.5}
\] Here \(a\) without a face subscript is the original spatial spacing.
All integer origins and angles in (1.5) are unchanged. The operator
\(\Gamma\) is diagonal on the joint link-spin decomposition, with
eigenvalue \(\sum_ew_ej_e(j_e+1)\); its self-adjoint domain is the set
of coefficient families square summable after multiplication by that
eigenvalue. The inequality \(0\le\Gamma\le L^2H_0\) on these coordinates
proves \(\Gamma:H^{k+2}\to H^k\). Thus all expressions below applied to
the smooth vacuum lie in their asserted domains.

Let \(\Pi\) be the orthogonal vertex gauge projection and let
\((\mathcal T_\theta f)(U)=f((h_e^{-1}U_e)_e)\). The actual centered
physical state and its angular coefficient are \[
\begin{split}
 \chi_{\xi,\theta}
  &=\Pi\mathcal T_\theta\psi_\xi
    -\langle\psi_\xi,\mathcal T_\theta\psi_\xi\rangle\psi_\xi,\\
 v_\xi&=(\Gamma-\gamma_\xi)\psi_\xi,
 \qquad\gamma_\xi=\langle\psi_\xi,\Gamma\psi_\xi\rangle,\\
 \chi_{\xi,\theta}&=-\frac23\theta^2v_\xi+O_{H^k,L,\xi}(\theta^4).
\end{split}
\tag{1.6}
\] The factor \(-2/3\) follows from the exact full-representation gauge
average. Each translated direction-1 link has a different canonical
source. Independent source twirls therefore act on its spin-\(j\) sector
by \(c_j(\theta n_2)=(2j+1)^{-1}\sum_{m=-j}^je^{2im\theta n_2}\). Its
second angular term is \(-2\theta^2n_2^2j(j+1)/3\), and its odd terms
vanish. The product over links proves (1.6). Taylor's integral formula
for the averaged translation flow bounds four derivatives
\(H^{k+4}\to H^k\), so this is a vector-domain expansion of the actual
vacuum. All spin sectors of that vacuum remain in the exact average.

Write \(P\) for the free physical energy-three face projection. Its
basis \(W_p\) is orthonormal: a link of unmatched central parity kills
an unequal-face Haar inner product, and an equal-face holonomy is Haar
with fundamental character squared integral one. Thus
\(\mathcal J_L(a_p)_p=\mathbf a_L=\sum_pa_pW_p\) is the exact isometry
from face coefficients into that subspace. A nonconstant physical spin
support has minimum vertex degree two, hence a cycle of at least four
links. Below energy \(9/2\) it can only be one elementary square with
four equal spin-\(1/2\) links. This proves that the free physical bottom
is simple, its next eigenvalue is 3 on this face space, and the next
possible energy is \(9/2\).

For real \(0<\xi\le3/(128M_L)\), let \[
\begin{gathered}
 P_\xi=\frac1{2\pi i}\int_{|z-3|=3/4}(z-K_\xi)^{-1}dz,
 \qquad T_\xi=PP_\xi P|_{\operatorname{ran}P},\\
 V_\xi=P_\xi P T_\xi^{-1/2}:
 \operatorname{ran}P\longrightarrow\operatorname{ran}P_\xi.
\end{gathered}
\tag{1.7}
\] The Neumann resolvent estimate using \(\|S\|=2M_L\) gives
\(\|P_\xi-P\|\le1/7\) on the larger complex disk \(|\xi|\le3/(64M_L)\).
Thus the inverse square root is its convergent binomial series. At real
coupling \(P_\xi\) is orthogonal and \(V_\xi^*V_\xi=I\); its range has
the same finite dimension as that of \(P\). The min--max bound
\(2M_L|\xi|\le3/64\) and the free isolation prove that \(P_\xi\) is the
actual entire first excitation band and that \(P_\xi\psi_\xi=0\)
exactly. In particular \[
 v_\xi=V_\xi(V_\xi^*v_\xi)+\rho_\xi,
 \quad \rho_\xi=(I-P_\xi)v_\xi,
 \quad \rho_\xi\perp\psi_\xi,
 \quad \rho_\xi\perp\operatorname{ran}P_\xi.
\tag{1.8}
\] This is the exact map to the interacting band; the same vacuum is
used on both sides. The finite-dimensional range is smooth by elliptic
regularity, so the map and the products in (1.8) respect the vector
domains used here.

For adjacent faces \(p,q\) sharing their unique physical edge \(e\),
define \[
 Y_p=W_p^2-1,\quad
 Z_{pq,0}=\int W_pW_q\,dU_e,\quad
 Z_{pq,1}=W_pW_q-Z_{pq,0},\quad s_{pq}=a_p+a_q.
\tag{1.9}
\] Haar integration gives \(Z_{pq,0}\) as one half of the fundamental
trace on their six-edge boundary. Hence the free electric energies and
squared norms of \(Y_p,Z_{pq,0},Z_{pq,1}\) are respectively
\((8,1),(9/2,1/4),(13/2,3/4)\). The last norm is the total product norm
one minus the orthogonal Haar projection norm \(1/4\). Different
unordered pairs have different outer-boundary central parity. Otherwise
their symmetric difference would be a closed mod-two collection of at
most four elementary squares; one square would require four distinct
other squares to cover its four edges, which is impossible. This proves
the pairwise orthogonalities; repeated faces have distinct joint spins
and even parity, giving the remaining orthogonalities.

The exact transport calculation yields \[
\begin{split}
 V_\xi^*v_\xi&=\xi\mathbf a_L+O_L(\xi^3),\\
 \rho_\xi&=\xi^2\mathbf r_L+O_{H^k,L}(\xi^3),\\
 \mathbf r_L&=\sum_p\frac{2a_p}{15}Y_p
 +\sum_{\{p,q\}\text{ adjacent}}
 \left[-\frac{2(s_{pq}+w_e)}9Z_{pq,0}
       +\frac{2(3s_{pq}+7w_e)}{273}Z_{pq,1}\right].
\end{split}
\tag{1.10}
\] Here is a direct coefficient verification that also specifies what
was subtracted. The actual vacuum has \[
\begin{split}
 \psi_\xi&=1+\frac\xi3S+\xi^2B^{\rm vac}+O_{H^k,L}(\xi^3),\\
 B^{\rm vac}&=-\frac{M_L}{18}+\sum_p\frac{Y_p}{24}
 +\sum_{\{p,q\}\text{ edge-disjoint}}\frac{W_pW_q}{9}
 +\sum_{\{p,q\}\text{ adjacent}}
     \left(\frac4{27}Z_{pq,0}+\frac4{39}Z_{pq,1}\right).
\end{split}
\tag{1.11}
\] Indeed coefficient comparison gives \(H_0B^{\rm vac}=(S^2-M_L)/3\)
and the unit norm gives its Haar mean \(-M_L/18\). Resolve the repeated
and unordered pair terms in \(S^2\) and divide by their energies
\(8,6,9/2,13/2\) to get (1.11). Analyticity in each fixed Sobolev space
follows from the isolated vacuum contour and repeated elliptic
improvement in its eigen-equation. Its scalar \(\gamma_\xi\) has second
coefficient \(\sum_pa_p/3\). The weighted electric eigenvalues of those
same four channels are \(8a_p,3s_{pq},3s_{pq}-3w_e/2,3s_{pq}+w_e/2\).

Let \(R_3=(I-P)(H_0-3)^{-1}(I-P)\), with inverse on
\(\operatorname{ran}(I-P)\). Differentiating (1.7) gives \(V'_0=R_3SP\).
Its value on \(\mathbf a_L\) has constant \(-\sum_pa_p/3\),
repeated-face coefficients \(a_p/5\), edge-disjoint coefficients
\(s_{pq}/3\), and adjacent coefficients \(2s_{pq}/3,2s_{pq}/7\). The
second coefficient of \(v_\xi\) is \(\Gamma B^{\rm vac}-(\sum_pa_p)/3\).
Subtracting the displayed spectral dressing cancels the identical
constant and edge-disjoint terms. The repeated difference is
\(a_p/3-a_p/5=2a_p/15\). The two adjacent differences are
\((4s-2w)/9-2s/3=-2(s+w)/9\) and \((12s+2w)/39-2s/7=2(3s+7w)/273\). This
proves (1.10), including the actual vacuum cancellation and both
nonlinear spin channels. Finite-angle edge-disjoint terms need not
vanish: their second coupling coefficient is
\((\alpha_p-1)(\alpha_q-1)W_pW_q/9\), where
\(\alpha_p=\prod_{f\in\partial p,\ f=(n,1)}\cos(\theta n_2)\). They have
angular order \(\theta^4\), explaining their absence in (1.10) while
retaining them in the exact state (1.6).

Define the unchanged mass and leakage coefficient by \[
\begin{split}
 A_L&=\|\mathbf a_L\|^2=\sum_pa_p^2,\\
 \mathsf R_L&=\|\mathbf r_L\|^2
 =\frac4{225}A_L
  +\frac1{81}\sum_{\{p,q\}\text{ adjacent}}(s_{pq}+w_e)^2
  +\frac1{24843}\sum_{\{p,q\}\text{ adjacent}}(3s_{pq}+7w_e)^2.
\end{split}
\tag{1.12}
\] The last denominator includes the squared norm \(3/4\):
\((3/4)\cdot4/273^2=1/24843\). Thus (1.12) is exactly the unnormalized
Hilbert norm, not a changed representation amplitude. The symbol
\(\mathsf R_L\) here denotes leakage; it is distinct from the
graph-profile residual named \(R_L\) in the companion graph-measure
calculation. The angular coefficient named \(\mathcal A_L\) in
MAGNETIC\_TRANSLATION\_PLAQUETTE.md is \(\mathcal A_L=4A_L\).

\subsection{2. Exhaustive adjacent-face classification in the original
box}\label{magnetic-part-11-exhaustive-adjacent-face-classification-in-the-original-box}

For an original transverse coordinate \(m\in\{-L,\ldots,L\}\), write \[
 \Sigma_L(m)=\{+1:m<L\}\cup\{-1:m>-L\},
 \qquad k_L(m)=|\Sigma_L(m)|
       =\begin{cases}1,&|m|=L,\\2,&|m|<L.\end{cases}
\tag{2.1}
\] The notation lists available physical sides; it creates no wrapping
side. Its exact sum is \(\sum_{m=-L}^Lk_L(m)=4L\). For arbitrary real
scalar parameters \(c,d\), introduce the polynomial sum \[
 H_L(c,d)=\sum_{\{p,q\}\text{ adjacent}}
                 [c(a_p+a_q)+d w_{e(p,q)}]^2.
\tag{2.2}
\] We will evaluate it at \((c,d)=(1,1)\) and \((3,7)\), precisely the
two original sums in (1.12). The parameters do not change the link
weights, face vectors, or Hamiltonian. Put \(t=c+d\).

Each adjacent pair has a unique shared edge, since two distinct
elementary squares share at most one edge. For an edge in direction
\(i\), its incident faces have one of the two transverse directions,
with the available signs in (2.1). An unordered pair is either the two
opposite faces of one plane, or one face of each transverse plane. These
are exactly the two coplanar types and the hinged type listed below.

For a direction-1 edge with original source \((x,m,z)\), \(-L\le x<L\),
\(-L\le m,z\le L\), its incident 12-face on side
\(\sigma\in\Sigma_L(m)\) has weight \[
 t_\sigma(m)=\frac{m^2+(m+\sigma)^2}{4}
       =\frac{m^2}{2}+\frac{\sigma m}{2}+\frac14;
\tag{2.3}
\] its incident 13-faces each have weight \(m^2/2\). Here \(w_e=m^2\).
The three types therefore contribute \[
\begin{split}
 H_{1,12}(c,d)&=2L(2L+1)
       \sum_{m=-L+1}^{L-1}(t m^2+c/2)^2,\\
 H_{1,13}(c,d)&=2L(2L-1)t^2\sum_{m=-L}^{L}m^4,\\
 H_{1,\rm hinged}(c,d)&=8L^2\sum_{m=-L}^L
       \sum_{\sigma\in\Sigma_L(m)}
               (t m^2+c\sigma m/2+c/4)^2.
\end{split}
\tag{2.4}
\] For the first line both 12-sides exist exactly when \(|m|<L\); their
weights add to \(m^2+1/2\), with \(2L\) values of \(x\) and \(2L+1\) of
\(z\). For the second line both 13-sides exist exactly when \(|z|<L\);
their weights add to \(m^2\), with \(2L(2L-1)\) choices of \((x,z)\). In
the hinged line their sum is \(m^2+\sigma m/2+1/4\). The number of
available 13-sides summed over \(z\) is \(4L\); multiplying by the
\(2L\) values of \(x\) gives \(8L^2\). Thus every boundary side has its
actual multiplicity.

For a direction-2 edge with source \((x,m,z)\), \(-L\le m<L\),
\(-L\le x,z\le L\), one has \(w_e=0\). Both possible 12-faces have
weight \(r_m/2\), where \(r_m=m^2+m+1/2\), and each possible 23-face has
weight zero. The coplanar 12-pair requires \(|x|<L\) and has sum
\(r_m\). A hinged 12/23-pair has sum \(r_m/2\), with
\((\sum_xk_L(x))(\sum_zk_L(z))=16L^2\) possible sides at each \(m\). The
coplanar 23-pair contributes zero. Consequently \[
 H_2(c,d)=c^2(2L-1)(2L+1)\sum_{m=-L}^{L-1}r_m^2
        +4L^2c^2\sum_{m=-L}^{L-1}r_m^2
        =c^2(8L^2-1)\sum_{m=-L}^{L-1}r_m^2.
\tag{2.5}
\] The dependence on \(d\) vanishes because the original shared-edge
weight is zero, not because the faces or their interactions were
removed.

For a direction-3 edge with source \((x,m,z)\), \(-L\le z<L\),
\(-L\le x,m\le L\), again \(w_e=0\). Its possible 13-faces each have
weight \(m^2/2\), and its possible 23-faces each have weight zero. A
coplanar 13-pair contributes \(c^2m^4\) with \(2L(2L-1)\) choices of
\((z,x)\). A hinged 13/23-pair contributes \(c^2m^4/4\), and its
\((z,x,\text{sides})\) count is \((2L)(4L)k_L(m)=8L^2k_L(m)\). The
coplanar 23-pair contributes zero. Thus \[
 H_3(c,d)=2L(2L-1)c^2\sum_{m=-L}^Lm^4
       +2L^2c^2\sum_{m=-L}^Lk_L(m)m^4.
\tag{2.6}
\]

As a count independent of all weights, each edge direction has
\(2\cdot(2L)(2L-1)(2L+1)\) coplanar pairs and \((2L)(4L)^2\) hinged
pairs. Their sum is \(48L^3-4L\) per direction. Summing the three
directions gives exactly \(144L^3-12L\) unordered adjacent pairs. This
verifies that (2.4)--(2.6) include every pair, including the zero-weight
types.

\subsection{3. Exact finite sums and complete polynomial
evaluation}\label{magnetic-part-11-exact-finite-sums-and-complete-polynomial-evaluation}

Define the original one-sided and symmetric sums \[
 S_j(L)=\sum_{m=-L}^{L-1}m^j,
 \qquad F_j(L)=\sum_{m=-L}^{L}m^j,
 \qquad I_j(L)=\sum_{m=-L+1}^{L-1}m^j.
\tag{3.1}
\] The needed exact values are \[
\begin{array}{c|ccccc}
 j&0&1&2&3&4\\\hline
 S_j&2L&-L&(2L^3+L)/3&-L^3&(6L^5+10L^3-L)/15.
\end{array}
\tag{3.2}
\] Also \[
\begin{split}
 F_4&=(6L^5+15L^4+10L^3-L)/15,\\
 I_4&=(6L^5-15L^4+10L^3-L)/15,
 \qquad I_2=(2L^3-3L^2+L)/3,\qquad I_0=2L-1,\\
 \sum_{m=-L}^{L-1}r_m^2&=\frac{L(4L^4+1)}{10},\qquad
 \sum_{m=-L}^{L}k_L(m)m^4=2F_4-2L^4.
\end{split}
\tag{3.3}
\] For a direct proof of (3.2), all five formulas have the correct value
at \(L=1\). The increment of the defining sum from \(L\) to \(L+1\) adds
the two original endpoint values \((-L-1)^j+L^j\). Expanding each
displayed polynomial's increment gives those same two values, proving
the formulas by induction. Adding the endpoint \(L\) and removing the
two endpoints \(\pm L\) proves the formulas for \(F_4\) and \(I_j\).
Finally \(r_m^2=m^4+2m^3+2m^2+m+1/4\) gives its sum from (3.2); the last
formula uses precisely the two endpoint deficiencies of \(k_L\). These
evaluations retain the negative odd sums \(-L,-L^3\).

In the last sum in (2.4), the \(\sigma=+1\) range is \(-L\le m<L\). Its
\(\sigma=-1\) range is \(-L<m\le L\). The change of summation index
\(m\mapsto-m\) in the latter shows that its value equals that of the
former, including their boundary terms. This is an exact evaluation of
both original ranges. Therefore \[
\begin{split}
 H_{1,\rm hinged}(c,d)
 =16L^2\bigl[t^2S_4+tcS_3
       +(tc/2+c^2/4)S_2+(c^2/4)S_1+(c^2/16)S_0\bigr],\\
 H_{1,12}(c,d)=2L(2L+1)
         \bigl[t^2I_4+tcI_2+(c^2/4)I_0\bigr],\qquad
 H_{1,13}(c,d)=2L(2L-1)t^2F_4.
\end{split}
\tag{3.4}
\] Substitution of (3.2)--(3.3) yields the three complete direction
polynomials, with \(H_1=H_{1,12}+H_{1,13}+H_{1,\rm hinged}\): \[
\begin{split}
 H_1(c,d)={}&\frac{48}{5}(c+d)^2L^7
 +(20c^2/3+16cd+12d^2)L^5
 -\frac83(c^2+cd)L^4\\
 &+(26c^2/15-6cd/5-8d^2/5)L^3
 +\frac23(c^2+cd)L^2-\frac12c^2L,\\
 H_2(c,d)={}&c^2(16L^7/5-2L^5/5+4L^3/5-L/10),\\
 H_3(c,d)={}&c^2(16L^7/5+16L^6/5+10L^5/3
                       -4L^4/3-8L^3/15+2L^2/15).
\end{split}
\tag{3.5}
\] For example, (2.6) becomes \(c^2[(8L^2-2L)F_4-4L^6]\) before
expansion, proving its direction-3 boundary term. Summing (3.5) gives
the full pair polynomial \[
\boxed{\begin{split}
 H_L(c,d)={}&(16c^2+96cd/5+48d^2/5)L^7+(16c^2/5)L^6\\
 &+(48c^2/5+16cd+12d^2)L^5-(4c^2+8cd/3)L^4\\
 &+(2c^2-6cd/5-8d^2/5)L^3
 +(4c^2/5+2cd/3)L^2-(3c^2/5)L.
\end{split}}
\tag{3.6}
\] Every coefficient in (3.6) is thus obtained from the explicit edge
classification and exact finite sums, independently of any
interpolation. The two required evaluations are \[
\boxed{\begin{split}
 H_L(1,1)&=\frac{224}{5}L^7+\frac{16}{5}L^6
 +\frac{188}{5}L^5-\frac{20}{3}L^4-\frac45L^3
 +\frac{22}{15}L^2-\frac35L,\\
 H_L(3,7)&=\frac{5088}{5}L^7+\frac{144}{5}L^6
 +\frac{5052}{5}L^5-92L^4-\frac{428}{5}L^3
 +\frac{106}{5}L^2-\frac{27}{5}L.
\end{split}}
\tag{3.7}
\]

For completeness, the unchanged face-vector mass follows from the same
coordinate counts. At each \(-L\le m<L\) there are \(2L(2L+1)\) 12-faces
of weight \(r_m/2\). At each \(-L\le m\le L\) there are \(4L^2\)
13-faces of weight \(m^2/2\). Hence \[
\begin{split}
 A_L&=\frac14\left[
 2L(2L+1)\frac{L(4L^4+1)}{10}+4L^2F_4\right]\\
 &=\frac{L^2(2L+1)}{60}(24L^4+24L^3+8L^2-4L+3)\\
 &=\frac45L^7+\frac65L^6+\frac23L^5+\frac1{30}L^3+\frac1{20}L^2.
\end{split}
\tag{3.8}
\] In particular \(A_L\ge4L^7/5>0\).

\subsection{4. The complete leakage polynomial and sharp
ratio}\label{magnetic-part-11-the-complete-leakage-polynomial-and-sharp-ratio}

The three physically distinct channel masses are exactly \[
 \mathsf R_{8,L}=\frac4{225}A_L,\qquad
 \mathsf R_{9/2,L}=\frac{H_L(1,1)}{81},\qquad
 \mathsf R_{13/2,L}=\frac{H_L(3,7)}{24843}.
\tag{4.1}
\] Equations (3.7)--(3.8), with these distinct factors retained, give \[
\boxed{\begin{split}
 \mathsf R_L
 =\frac1{251535375}\bigl(&153001392L^7+15594888L^6
 +129973560L^5-21634000L^4\\
 &-3201942L^3+4992787L^2-1917900L\bigr).
\end{split}}
\tag{4.2}
\] Thus its leading coefficient is \[
 \lim_{L\to\infty}\frac{\mathsf R_L}{L^7}
       =\frac{51000464}{83845125}.
\tag{4.3}
\] The negative coefficients at degrees 4, 3, and 1 are retained
open-boundary terms. Positivity of the full polynomial for every allowed
\(L\) follows from the original orthogonal sum (1.12), with the
quantitative bound \(\mathsf R_L\ge4A_L/225>0\).

Set \[
 C_{\rm leak}=\frac{12750116}{16769025}
       =\frac4{225}+\frac{56}{81}+\frac{424}{8281},
 \qquad b_{\rm edge}=\frac{54848}{51597}.
\tag{4.4}
\] The exact ratio and its sharp first boundary correction are \[
\boxed{
 \frac{\mathsf R_L}{A_L}
 =C_{\rm leak}-\frac{b_{\rm edge}}L
      +\frac{3233332}{2012283L^2}+O(L^{-3}).}
\tag{4.5}
\] In particular the limiting constant is approximately
\(0.760337348176176\). There is no normalization of the original state
in (4.5): it is the ratio of its two already specified
Taylor-coefficient masses.

To verify both the limit and the correction exactly, define the two
explicit polynomials \[
\begin{split}
 P_5(L)&=2139072L^5-24724L^4+216340L^3
                          +95770L^2+45698L+19179,\\
 P_4(L)&=3233332L^4+1566220L^3-95770L^2
                                      +43430L+114513.
\end{split}
\tag{4.6}
\] Multiplication of (3.8) and (4.2) gives the exact identities \[
 C_{\rm leak}A_L-\mathsf R_L
       =\frac{4L}{10061415}P_5(L),
 \qquad
 \mathsf R_L-C_{\rm leak}A_L+\frac{b_{\rm edge}}L A_L
       =\frac{4L}{10061415}P_4(L).
\tag{4.7}
\] Their coefficients provide an independent arithmetic check on every
term in the ratio. For \(L\ge2\), the first polynomial is positive
because its first two terms are \(L^4(2139072L-24724)>0\) and the others
are positive. The second is positive because
\(L^2(3233332L^2+1566220L-95770)>0\) and its remaining terms are
positive. Consequently \[
 C_{\rm leak}-\frac{b_{\rm edge}}L
   <\frac{\mathsf R_L}{A_L}<C_{\rm leak}
      \quad(L\ge2),
\tag{4.8}
\] and, using \(A_L\ge4L^7/5\) and omitting just the negative term when
bounding \(P_4\), \[
 0<\frac{\mathsf R_L}{A_L}-C_{\rm leak}+\frac{b_{\rm edge}}L
 \le\frac{24787475}{10061415L^2}.
\tag{4.9}
\] Indeed \(P_4(L)\le(3233332+1566220+43430+114513)L^4
=4957495L^4\) for \(L\ge1\). Finally the last quotient in (4.7) has
leading term \((4\cdot3233332/10061415)/(4/5)\,L^{-2}
=3233332/(2012283L^2)\); division of its explicit polynomials has next
remainder \(O(L^{-3})\), since the denominator divided by \(L^7\) tends
to \(4/5\) and is bounded away from zero. This proves (4.5), rather than
assuming the rate from numerical values.

If instead the denominator is the original angular coefficient
\(\mathcal A_L=4A_L\), the same exact map gives \[
 \frac{\mathsf R_L}{\mathcal A_L}
  =\frac{3187529}{16769025}-\frac{13712}{51597L}
        +\frac{808333}{2012283L^2}+O(L^{-3}).
\tag{4.10}
\] This separate formula records the factor of four rather than silently
using the two coefficient names interchangeably.

\subsection{5. Original amplitudes, exact spectral measures, and order
of
limits}\label{magnetic-part-11-original-amplitudes-exact-spectral-measures-and-order-of-limits}

Let \(\mathscr A_\xi=H-\mathcal E_\xi I
=\kappa(K_\xi-e(\xi)I)\). Define the actual out-of-band physical measure
of the native state by \[
 \chi^{\rm out}_{\xi,\theta}=(I-P_\xi)\chi_{\xi,\theta},
 \qquad
 \mu^{\rm out}_{L,\xi,\theta}(J)
 =\langle\chi^{\rm out}_{\xi,\theta},
                 1_J(\mathscr A_\xi)\chi^{\rm out}_{\xi,\theta}\rangle.
\tag{5.1}
\] The state in (5.1) is exactly orthogonal to the actual vacuum and
band, by (1.7)--(1.8). Define the fully evaluated dimensionless
coefficient measure \[
 \Lambda_L=\frac4{225}A_L\,\delta_8
       +\frac{H_L(1,1)}{81}\,\delta_{9/2}
       +\frac{H_L(3,7)}{24843}\,\delta_{13/2}.
\tag{5.2}
\] Its total mass is precisely \(\mathsf R_L\). The exact map from
dimensionless excitation energy to physical excitation energy is
\(u_\kappa(x)=\kappa x\), since (1.3)--(1.4) give
\(\mathscr A_\xi=\kappa(K_\xi-e(\xi)I)\). On the coefficient measure it
gives the physical atoms \[
 8\kappa=\frac{16g_{\rm YM}^2}{a},\qquad
 \frac92\kappa=\frac{9g_{\rm YM}^2}{a},\qquad
 \frac{13}2\kappa=\frac{13g_{\rm YM}^2}{a}.
\tag{5.3}
\] Functional calculus proves this pushforward for every Borel set, not
only for the three moments: \(1_J(\kappa A)=1_{u_\kappa^{-1}(J)}(A)\)
for a self-adjoint \(A\), by applying the spectral theorem to the
composed indicator.

For every bounded continuous \(f:\mathbb R\to\mathbb C\), the precise
ordered coefficient statement is \[
 \lim_{\xi\downarrow0}\frac1{\xi^4}
 \left[\lim_{\theta\to0}\frac9{4\theta^4}
       \int f(\lambda/\kappa)\,d\mu^{\rm out}_{L,\xi,\theta}(\lambda)
 \right]=\int f\,d\Lambda_L.
\tag{5.4}
\] To prove it, (1.6) supplies the inner limit as the spectral
expectation of \(\rho_\xi\), and (1.10) gives
\(\rho_\xi/\xi^2\to\mathbf r_L\) in \(H^2\). The bounded difference
\(K_\xi-e(\xi)-H_0\) tends to zero, giving norm-resolvent convergence by
the resolvent identity. Uniform approximation by the algebra generated
by nonreal resolvents proves convergence of \(C_0(\mathbb R)\)
functional calculus. Moreover the \(H^2\) convergence gives a bounded
second spectral moment of \(\rho_\xi/\xi^2\), so the mass outside
\([-R,R]\) is bounded by a constant times \(R^{-2}\). Approximate a
bounded continuous \(f\) by its continuous compactly supported cutoffs
in the expectations; this tail bound gives the same conclusion for that
full class. The energies and orthogonal norms in (1.9)--(1.12) yield
exactly the three terms of (5.2). This proves (5.4).

The parity cochain \(z_i(n)=(-1)^{\sum_{j<i}n_j}I\) sends every \(W_p\)
to \(-W_p\), fixes all link Casimirs, and commutes with the native
translation. Its unitary \(\mathcal Z\) satisfies
\(\psi_{-\xi}=\mathcal Z\psi_\xi\),
\(P_{-\xi}=\mathcal ZP_\xi\mathcal Z\), and
\(\rho_{-\xi}=\mathcal Z\rho_\xi\). Thus the squared norm is even
analytic. In particular the fixed-box scalar remainder and the fraction
of the original state's mass are \[
\begin{split}
 \lim_{\theta\to0}\frac9{4\theta^4}
               \|\chi^{\rm out}_{\xi,\theta}\|^2
       &=\xi^4\mathsf R_L+O_L(\xi^6),\\
 \lim_{\theta\to0}\frac9{4\theta^4}
               \|\chi_{\xi,\theta}\|^2
       &=\xi^2 A_L+O_L(\xi^4),\\
 \lim_{\theta\to0}
  \frac{\|\chi^{\rm out}_{\xi,\theta}\|^2}{\|\chi_{\xi,\theta}\|^2}
       &=\xi^2\frac{\mathsf R_L}{A_L}+O_L(\xi^4).
\end{split}
\tag{5.5}
\] The denominators are nonzero on a sufficiently small positive
coupling interval because \(A_L>0\). All vacuum centering and dressing
operations precede this division. The fixed-box mixed Sobolev Taylor
formula further gives \[
 \chi^{\rm out}_{\xi,\theta}
 =-\frac23\theta^2\xi^2\mathbf r_L
       +O_{H^k,L}(\theta^2|\xi|^3+\theta^4\xi^2).
\tag{5.6}
\] For its justification the smooth vacuum is analytic in \(\xi\) in all
needed higher Sobolev spaces, and averaged translations have four
bounded angular derivatives between those spaces. The out-of-band state
and its first coupling derivative vanish at \(\xi=0\) for every angle;
it is even and vanishes at \(\theta=0\). Taylor integral remainders
therefore factor \(\xi^2\theta^2\), with the indicated next coupling and
angular terms. The leading coefficient is (1.10).

The original squared-norm prefactors are exactly \[
\begin{split}
 \frac49\theta^4\xi^4\mathsf R_L
   &=\frac{\theta^4}{576g_{\rm YM}^{16}}\mathsf R_L
     =\frac{\pi^4a^8}{36D^4g_{\rm YM}^{16}}\mathsf R_L,\\
 \frac49\theta^4\xi^2 A_L
   &=\frac{\theta^4}{36g_{\rm YM}^{8}}A_L
     =\frac{4\pi^4a^8}{9D^4g_{\rm YM}^{8}}A_L.
\end{split}
\tag{5.7}
\] The unsquared leading out-of-band amplitude is
\((2/3)\theta^2\xi^2\sqrt{\mathsf R_L}
=\theta^2\sqrt{\mathsf R_L}/(24g_{\rm YM}^8)
=\pi^2a^4\sqrt{\mathsf R_L}/(6D^2g_{\rm YM}^8)\). Its physical
excitation-form leading coefficient, when needed, is \[
 \frac49\theta^4\xi^4\kappa
       \left[8\mathsf R_{8,L}
             +\frac92\mathsf R_{9/2,L}
             +\frac{13}2\mathsf R_{13/2,L}\right].
\tag{5.8}
\] This follows from (1.10) in \(H^1\) and the channel electric
energies; it retains their different energies and does not infer an
energy merely from a small unnormalized norm.

Since (5.2) has fixed three-point support, the coefficient-measure
limits are exact weak limits and also limits in total variation: \[
\boxed{\begin{split}
 \frac{\Lambda_L}{L^7}&\longrightarrow
  \frac{16}{1125}\delta_8+\frac{224}{405}\delta_{9/2}
                         +\frac{1696}{41405}\delta_{13/2},\\
 \frac{\Lambda_L}{A_L}&\longrightarrow
  \frac4{225}\delta_8+\frac{56}{81}\delta_{9/2}
                         +\frac{424}{8281}\delta_{13/2}.
\end{split}}
\tag{5.9}
\] Indeed every atom's coefficient converges by (3.7)--(3.8), and the
total variation of the difference is the sum of the three absolute
coefficient differences. The second limiting measure has mass
\(C_{\rm leak}\), rather than 1, because its denominator is the
unchanged \emph{first-band} face-vector mass \(A_L\). Composing (5.4)
and (5.9) is precisely the order \(\theta\to0\), then \(\xi\to0\), then
\(L\to\infty\). For example (5.5) proves \[
 \lim_{L\to\infty}\left[
 \lim_{\xi\downarrow0}\frac1{\xi^2}
  \left(\lim_{\theta\to0}
     \frac{\|\chi^{\rm out}_{\xi,\theta}\|^2}{\|\chi_{\xi,\theta}\|^2}
  \right)\right]=C_{\rm leak}.
\tag{5.10}
\]

All \(O_L\) terms in (5.5)--(5.6) remain fixed-box estimates. The proved
band interval itself requires \(\xi\le3/(128M_L)\), and the higher
Sobolev and coupling coefficients have no volume-uniform bound
established here. Neither (4.5) nor (5.9)--(5.10) asserts a
fixed-positive-coupling infinite-volume leakage fraction or a joint
continuum limit. What is completely evaluated is the specified first
native leakage coefficient, its full boundary polynomial, all its
spectral channels, and the original amplitude and physical-energy maps.

The companion checker
\url{check_native_magnetic_leakage_volume_polynomial.py} verifies the
polynomial algebra coefficient by coefficient with exact rational
polynomial operations and compares each edge-direction class with direct
original-box incidence enumeration. No interpolation is used to produce
or validate a purported proof of the polynomial.

\end{document}
